%%
%% Copyright 2024 OXFORD UNIVERSITY PRESS
%%
%% This file is part of the 'oup-authoring-template Bundle' for the Royal Statistical Society.
%% ---------------------------------------------
%%
%% It may be distributed under the conditions of the LaTeX Project Public
%% License, either version 1.2 of this license or (at your option) any
%% later version.  The latest version of this license is in
%%    http://www.latex-project.org/lppl.txt
%% and version 1.2 or later is part of all distributions of LaTeX
%% version 1999/12/01 or later.
%%
%% The list of all files belonging to the 'oup-authoring-template Bundle' is
%% given in the file `manifest.txt'.
%%
%% Template article for OXFORD UNIVERSITY PRESS's document class `oup-authoring-template'
%% with bibliographic references
%%

\documentclass[numsec,webpdf,modern,large,namedate]{oup-authoring-template}

\onecolumn

\graphicspath{{Fig/}}

\usepackage[doublespacing]{setspace}
\usepackage[fontsize=12pt]{fontsize}

\RequirePackage{amsthm,amsmath,amsfonts,amssymb}
\RequirePackage[authoryear]{natbib}%
\RequirePackage{graphicx}%
\usepackage{subfigure}
\usepackage{float}
\usepackage{multirow}
\usepackage{diagbox}
\usepackage[figuresright]{rotating}
\usepackage{pdflscape}
\usepackage{bm}
\usepackage{algorithm}
\usepackage{enumerate}
\usepackage{booktabs}
\usepackage{mathabx}

\theoremstyle{plain}
\newtheorem{axiom}{Axiom}
\newtheorem{claim}[axiom]{Claim}
\newtheorem{stheorem}{Theorem}%
\newtheorem{slemma}{Lemma}%
\newtheorem{sremark}{Remark}%
\newtheorem{scorollary}{Corollary}%
\renewcommand{\theslemma}{S.\arabic{slemma}}
\renewcommand{\thesremark}{S.\arabic{sremark}}
\renewcommand{\thestheorem}{S.\arabic{stheorem}}
\renewcommand{\thescorollary}{S.\arabic{scorollary}}
\renewcommand{\thealgorithm}{S.\arabic{algorithm}}

\renewcommand{\thesection}{S.\arabic{section}}
\theoremstyle{remark}
\newtheorem{definition}{Definition}[section]
\newtheorem*{example}{Example}
\newtheorem*{fact}{Fact}
\newtheorem*{pft}{Proof of Theorem}
\newtheorem{sassumption}{Assumption}
\newtheorem{corollary}{Corollary}[section]
\renewcommand{\thesassumption}{S.\arabic{sassumption}}
\newtheorem{assertion}{Assertion}

\newcounter{cnstcnt}
\newcommand{\xc}{%
   \refstepcounter{cnstcnt}%
   \ensuremath{c_{\thecnstcnt}}
}
\newcommand{\jc}[1]{\ensuremath{c_{\ref{#1}}}}

\newcounter{bnstcnt}
\newcommand{\xb}{%
   \refstepcounter{bnstcnt}%
   \ensuremath{b_{\thebnstcnt}}
}
\newcommand{\jb}[1]{\ensuremath{b_{\ref{#1}}}}

\newcounter{dbnstcnt}
\newcommand{\xdb}{%
   \refstepcounter{dbnstcnt}%
   \ensuremath{B_{\thedbnstcnt}}
}
\newcommand{\jdb}[1]{\ensuremath{B_{\ref{#1}}}}

\makeatletter
\newcommand*{\centerfloat}{%
  \parindent \z@
  \leftskip \z@ \@plus 1fil \@minus \marginparwidth
  \rightskip \leftskip
  \parfillskip \z@skip}
\makeatother

\newtheorem{manualtheoreminner}{Assumption}
\newenvironment{manualtheorem}[1]{%
  \renewcommand\themanualtheoreminner{#1}%
  \manualtheoreminner
}{\endmanualtheoreminner}


\newcommand{\M}{\bm{M}}
\newcommand{\bH}{\bm{H}}
\newcommand{\Y}{\bm{Y}}
\newcommand{\X}{\bm{X}}
\newcommand{\bOmega}{\bm{\Omega}}
\newcommand{\U}{\bm{U}}
\newcommand{\V}{\bm{V}}
\newcommand{\bu}{\bm{u}}
\newcommand{\bv}{\bm{v}}
\newcommand{\x}{\bm{x}}
\newcommand{\y}{\bm{y}}
\newcommand{\N}{\bm{N}}
\newcommand{\bL}{\bm{L}}
\newcommand{\R}{\bm{R}}

\providecommand{\op}{\mathrm{op}}
\providecommand{\F}{\mathrm{F}}
\providecommand{\PP}{P_{\bOmega}}
\providecommand{\TT}{T_\Omega}
\providecommand{\norm}[1]{\left\|#1\right\|}
\providecommand{\inner}[2]{\left\langle #1,#2\right\rangle}
\DeclareMathOperator{\SVT}{SVT}
\DeclareMathOperator{\Clip}{OpClip}
\DeclareMathOperator{\rank}{rank}
\DeclareMathOperator{\diag}{diag}
\usepackage{algpseudocode}
\usepackage{xr-hyper}
\externaldocument[main-]{Residual_spectral_matching}


\AtBeginDocument{\setlength{\headheight}{24pt}\addtolength{\topmargin}{-8pt}}


\usepackage{xcolor}

% BEGIN MAJOR-REVISION MARKUP DEFINITIONS
% Set \markchangesfalse to suppress the selective revision colouring.
\newif\ifmarkchanges
\markchangesfalse
\DeclareRobustCommand{\rev}[1]{{\ifmarkchanges\color{red}\fi#1}}
\newenvironment{revision}{\begingroup\ifmarkchanges\color{red}\fi}{\endgroup}
\DeclareRobustCommand{\revisionfloatcolor}{\ifmarkchanges\color{red}\fi}
% END MAJOR-REVISION MARKUP DEFINITIONS

\begin{document}
\raggedbottom
\vspace*{8pt}

\bigskip
\bigskip
\bigskip
\begin{center}
{\LARGE\bfseries Supplement to ``Residual spectral matching: a new perspective for matrix completion''}
\end{center}
\begin{center}
Ziyuan Chen and Fang Yao
\end{center}

% BEGIN MAJOR REVISION S01
\begin{revision}
\section{Notation and organization}
In this supplementary, the notation and assumptions are the same as those of the main paper. Section~\ref{sec:supp_sparse_rmt} establishes spectral and scale calibration. Sections~\ref{sec:proof_factor} and \ref{sec:nuclear_estimation} prove the estimation results. Sections~\ref{sec:initialization_bounds} and \ref{sec:estimator_algorithms} provide the algorithmic analysis. The heterogeneous noise extension is proved in Section~\ref{sec:heterogeneous_noise}. Section~\ref{sec:weighted_certification} establishes recovery with unknown profile distributions.

For the homogeneous model, $Y=P_{\bOmega}(M_0+H)$ with a fixed rank-$r$ matrix $M_0\in\mathbb R^{m\times n}$, $m\ge n$, and $0<\rho_-\le n/m\le1$.  The entries of $H$ are independent, centered, have common variance $\sigma^2>0$ and satisfy $\mathbb P(|H_{ij}|>t\sigma)\le2e^{-t/K}$ uniformly. The entries of $\bOmega$ are independent Bernoulli$(p)$ and independent of $H$; the Gaussian reference draws are independent of the observations. The singular spaces of $M_0$ are $\mu_0$-incoherent. Write $s_\star=s_r(M_0)>0$, $s_1=s_1(M_0)$, $\kappa=s_1/s_\star$, and $h=\sigma\sqrt{m/p}$.
The homogeneous spectral results assume $mp\ge\log^{\alpha_0}m$ for fixed $\alpha_0>15$ and sufficiently large dimensions. Each statistical or algorithmic result states its additional structural sampling condition. We assume $12r\le n$ throughout the rank constrained results.

The notation $\mathbb M_{r,\mu_1}$ denotes the class of matrices of rank at most $r$ represented as $M=LR'$ with $r$ columns in each factor and projector leverage bounds $\mu_1r/m,\mu_1r/n$ on its nonzero singular spaces. The zero matrix has zero leverage. For any matrix $A$, $s_i(A)=\psi_i(A)$ denotes its $i$th singular value, in decreasing order. Write $I=\{\lceil n/3\rceil,\ldots,\lfloor2n/3\rfloor\}$ and
$\delta_m=(\log m)^{-1}$. Weights are nonnegative, sum to one, and are fixed during optimization. The reference middle sum is positive on the stated events. Statements about global minimizers assume their existence. Candidates with zero bulk scale are assigned infinite loss. The heterogeneous model and its distinct fine/coarse assumptions are specified separately in Section~\ref{sec:heterogeneous_noise}.
\end{revision}
% END MAJOR REVISION S01


% BEGIN MAJOR REVISION S02
\begin{revision}
\section{Sparse spectral calibration at the true matrix}
\label{sec:supp_sparse_rmt}
We establish the spectral approximation used throughout the proofs. The observation probability in this paper may be as small as a polylogarithmic multiple of $m^{-1}$. Consequently, the entries of the variance-normalized matrix $p^{-1/2}P_{\bOmega}(\bH)$ have diverging higher moments as $p\to0$ and dense sample covariance results cannot be applied directly to $P_{\bOmega}(\bH)$.  We instead use the universality theory of \cite{brailovskaya2024universality}, which applies to sparse sums of
independent random matrices. Dense Gaussian/Wishart theory enters only for the covariance matched Gaussian comparator supplied by that universality principle.

Throughout this block we consider the aspect-ratio regime
\begin{align}
\label{eq:supp_aspect_ratio}
    0<\rho_-\le \rho=\frac{n}{m}\le1
\end{align}
for a fixed constant $\rho_->0$.  This is the usual regime underlying the rooted Marchenko--Pastur approximation used in the paper. We write $s_1(\bm A)\ge\cdots\ge s_n(\bm A)$ for the ordered singular values of an $m\times n$ matrix $\bm A$.  Let
\begin{align}
\label{eq:supp_rooted_mp_density}
    \nu_{\rho}(x)
    = \frac{1}{\pi}\frac{
    \sqrt{(\lambda_+^2-x^2)(x^2-\lambda_-^2)}}{\rho x}
    \mathbf 1_{\{\lambda_-\le x\le\lambda_+\}},
    \qquad
    \lambda_{\pm}=1\pm\sqrt{\rho},
\end{align}
be the rooted Marchenko--Pastur density, and let $\gamma_1\ge\cdots\ge\gamma_n$ be the deterministic quantiles defined by
\begin{align}
\label{eq:supp_mp_quantiles}
    \int_{\gamma_i}^{\lambda_+}\nu_{\rho}(x)\,dx
    =\frac{i-1/2}{n},\qquad 1\le i\le n.
\end{align}

\begin{slemma}[Coarse rigidity for sparse rectangular matrices]
\label{lem:sparse_coarse_rigidity}
Let $\{\xi_{ij}:1\le i\le m,1\le j\le n\}$ be independent random variables satisfying uniformly in $i,j$,
\begin{align}
\label{eq:supp_sparse_entry_assumption}
    \mathbb E\xi_{ij}=0,\qquad
    \mathbb E\xi_{ij}^2=1,\qquad
    \mathbb P(|\xi_{ij}|>t)\le 2\exp(-t/K),\quad t>0.
\end{align}
Let $\Omega_{ij}\stackrel{\mathrm{i.i.d.}}{\sim}\operatorname{Bernoulli}(p)$ be independent of the $\xi_{ij}$, and define
\begin{align*}
    \bm A=\frac{1}{\sqrt{mp}}P_{\bOmega}(\bm\Xi),
    \qquad \bm\Xi=(\xi_{ij})_{m\times n}.
\end{align*}
For every fixed $\alpha_0>15$, if $mp\ge\log^{\alpha_0}m$, then for some constant $D>0$,
\begin{align}
\label{eq:supp_sparse_coarse_rigidity}
    \mathbb P\left
    \{\max_{1\le i\le n}|s_i(\bm A)-\gamma_i|
    >C\log^{-1}m\right\}
    \le c n^{-D}.
\end{align}
In particular, with the same probability,
\begin{align}
\label{eq:supp_sparse_edge_bound}
    s_1(\bm A)\le1+\sqrt\rho+C\log^{-1}m.
\end{align}
\end{slemma}
\begin{proof}[Proof.]
Set $\delta_m=\log^{-1}m$. We use the following sparse universality comparison.
Let $\bm X=\sum_{a=1}^N\bm Z_a$ be a centered self-adjoint random matrix with
independent summands, and let $\bm G$ be the centered Gaussian matrix with the
same covariance as $\bm X$.  If $\operatorname{tr}=d^{-1}\operatorname{Tr}$
denotes normalized trace, then Remark 6.13 of
\cite{brailovskaya2024universality} implies, whenever the third moments on the
right-hand side are finite,
\begin{align}
\label{eq:supp_bvh_trace_resolvent}
    \left|
    \mathbb E\operatorname{tr}(z\bm I-\bm X)^{-1}
    -\mathbb E\operatorname{tr}(z\bm I-\bm G)^{-1}
    \right|
    \lesssim
    \frac{1}{(\operatorname{Im}z)^4}
    \sum_{a=1}^N\mathbb E\operatorname{tr}|\bm Z_a|^3,
    \qquad \operatorname{Im}z>0.
\end{align}
For spectral edges we use the bounded-summand edge universality result,
Corollary 2.7 of \cite{brailovskaya2024universality}, after logarithmic
truncation.  For the lower singular-value edge we use Theorem 3.16 of the
same paper.  These results are applied only after verifying their sparse
matrix parameters below.

We divide the argument into four steps.

\textit{Step 1: Hermitian dilation and sparse universality parameters.}
Consider
\begin{align*}
    \bm X
    =\begin{pmatrix}
      \bm0&\bm A\\
      \bm A'&\bm0
    \end{pmatrix}
    =\sum_{i=1}^m\sum_{j=1}^n\bm Z_{ij},
\end{align*}
where
\begin{align*}
    \bm Z_{ij}
    =\frac{\Omega_{ij}\xi_{ij}}{\sqrt{mp}}
      (\bm e_i\bm e_{m+j}'+\bm e_{m+j}\bm e_i').
\end{align*}
The positive eigenvalues of $\bm X$ are exactly the singular values of $\bm A$.  Since
$\operatorname{Var}(\Omega_{ij}\xi_{ij}/\sqrt{mp})=m^{-1}$, following the notation in \cite{brailovskaya2024universality}, a direct calculation gives
\begin{align*}
    \sigma(\bm X)=1,\qquad
    \sigma_*(\bm X)\lesssim m^{-1/2},\qquad
    v(\bm X)\lesssim m^{-1/2}.
\end{align*}
Moreover, using normalized trace on $\mathbb R^{m+n}$ and the uniformly bounded third moments implied by \eqref{eq:supp_sparse_entry_assumption},
\begin{align}
\label{eq:supp_sparse_third_trace}
    \sum_{i=1}^m\sum_{j=1}^n
    \mathbb E\operatorname{tr}|\bm Z_{ij}|^3
    \lesssim (mp)^{-1/2}.
\end{align}
Summing over $mn$ locations gives exactly the order $(mp)^{-1/2}$ under \eqref{eq:supp_aspect_ratio}. The centered Gaussian matrix with the same covariance as $\bm X$ is the Hermitian dilation of
\begin{align*}
    \bm G_0=m^{-1/2}(g_{ij})_{m\times n},
    \qquad g_{ij}\stackrel{\mathrm{i.i.d.}}{\sim}N(0,1).
\end{align*}
Thus the Gaussian comparator is dense, but it is used only after the sparse
universality reduction.

\textit{Step 2: mesoscopic spectral counting.}
Let $F_{\bm A}$ be the empirical distribution function of the $n$ singular
values of $\bm A$, and let $F_\rho$ be the distribution function associated
with \eqref{eq:supp_rooted_mp_density}.  From
\eqref{eq:supp_bvh_trace_resolvent} and
\eqref{eq:supp_sparse_third_trace}, for $z=E+\mathrm{i}\eta$,
\begin{align}
\label{eq:supp_sparse_stieltjes_compare}
    \left|
    \mathbb E m_{\bm X}(z)-\mathbb E m_{\bm G}(z)
    \right|
    \lesssim\frac{1}{\sqrt{mp}\,\eta^4},
\end{align}
where $m_{\bm X}$ and $m_{\bm G}$ denote the normalized Stieltjes transforms
of the two Hermitian dilations.  The ordinary dense random matrix input is used only for the comparator
$\bm G_0$.  The classical Marchenko--Pastur convergence theory for dense
sample covariance matrices gives that in the Gaussian case, an essentially
$m^{-1}$ Kolmogorov rate up to logarithmic factors \citep{gotze2011rate}.  Applying that high-probability bound to the dense
Gaussian comparator and using $d_{\rm K}\le1$ on the exceptional event gives
\begin{align}
\label{eq:supp_dense_gaussian_mp}
    \sup_x\left|\mathbb E F_{\bm G_0}(x)-F_\rho(x)\right|
    \le \mathbb E d_{\rm K}(F_{\bm G_0},F_\rho)
    \le C m^{-1}\log^{C_0}m.
\end{align}
The same rate transfers from the sample-covariance eigenvalue distribution to
the rooted singular-value distribution by the monotone transformation
$x\mapsto\sqrt{x}$.  The smoothing argument is as follows.

The Hermitian dilation has $m-n$ deterministic zero eigenvalues.  To avoid
introducing this common atom into the smoothing inequality, define instead
the symmetrized singular value measures
\begin{align*}
    \bar\mu_{\bm A}
    &=\frac1{2n}\sum_{i=1}^n
      \{\delta_{s_i(\bm A)}+\delta_{-s_i(\bm A)}\},\\
    \bar\mu_{\bm G_0}
    &=\frac1{2n}\sum_{i=1}^n
      \{\delta_{s_i(\bm G_0)}+\delta_{-s_i(\bm G_0)}\},
\end{align*}
and let $\bar\mu_\rho$ be the deterministic symmetric probability measure
whose positive half has density $\nu_\rho$.  Since the zero atoms of the two
Hermitian dilations are identical,
\begin{align}
\label{eq:supp_sym_stieltjes_conversion}
    \left|
    \mathbb E m_{\bar\mu_{\bm A}}(z)
    -\mathbb E m_{\bar\mu_{\bm G_0}}(z)
    \right|
    \le C_{\rho_-}
    \left|
    \mathbb E m_{\bm X}(z)-\mathbb E m_{\bm G}(z)
    \right|.
\end{align}
Moreover, \eqref{eq:supp_dense_gaussian_mp} and integration by parts imply
\begin{align}
\label{eq:supp_dense_gaussian_stieltjes_mp}
    \left|
    \mathbb E m_{\bar\mu_{\bm G_0}}(E+\mathrm i\eta)
    -m_{\bar\mu_\rho}(E+\mathrm i\eta)
    \right|
    \lesssim \frac{m^{-1}\log^{C_0}m}{\eta}.
\end{align}
% Indeed, $\int_{\mathbb R}|x-(E+\mathrm i\eta)|^{-2}\,dx=\pi/\eta$.

A finite-interval version of Bai's smoothing inequality
\citep[Chapter 8]{bai2010spectral}, applied directly against the
deterministic measure $\bar\mu_\rho$, yields for $B=4$ and $0<\eta<1$,
\begin{align}
\label{eq:supp_bai_smoothing}
    d_{\rm K}(\mathbb E\bar\mu_{\bm A},\bar\mu_\rho)
    \lesssim&
    \int_{-B}^{B}
       \left|
       \mathbb E m_{\bar\mu_{\bm A}}(u+\mathrm i\eta)
       -m_{\bar\mu_\rho}(u+\mathrm i\eta)
       \right|\,du
       +\eta \notag\\
    &+\eta^{-1}\int_{|x|>B}
       |F_{\mathbb E\bar\mu_{\bm A}}(x)
       -F_{\bar\mu_\rho}(x)|\,dx .
\end{align}
The $\eta$ term is uniform because the symmetrized rooted
Marchenko--Pastur density is uniformly bounded under
\eqref{eq:supp_aspect_ratio}.  The tail integral is negligible.  Indeed, the
sparse edge bound in Step 3 is proved independently of the present counting
argument and gives $\mathbb P\{s_1(\bm A)>3\}\le cn^{-D_0}$ for
arbitrarily large $D_0$, while
$\mathbb E s_1(\bm A)^2\le\mathbb E\|\bm A\|_F^2=n$.
Cauchy--Schwarz therefore makes the tail contribution smaller than any
prescribed polynomial order.

Combining \eqref{eq:supp_sparse_stieltjes_compare},
\eqref{eq:supp_sym_stieltjes_conversion},
\eqref{eq:supp_dense_gaussian_stieltjes_mp}, and
\eqref{eq:supp_bai_smoothing} gives
\begin{align}
\label{eq:supp_sparse_expected_cdf}
    d_{\rm K}(\mathbb E\bar\mu_{\bm A},\bar\mu_\rho)
    \lesssim
    \eta+\frac{1}{\sqrt{mp}\,\eta^4}
    +\frac{m^{-1}\log^{C_0}m}{\eta}+n^{-D}.
\end{align}
Finally, for every $x\ge0$,
\begin{align*}
    F_{\bar\mu_{\bm A}}(x)
    =\frac12+\frac12F_{\bm A}(x),
    \qquad
    F_{\bar\mu_\rho}(x)
    =\frac12+\frac12F_\rho(x).
\end{align*}
Thus \eqref{eq:supp_sparse_expected_cdf} yields the same bound for
$\sup_{x\ge0}|\mathbb E F_{\bm A}(x)-F_\rho(x)|$.

Choosing $\eta=(mp)^{-1/10}$, the first two terms in \eqref{eq:supp_sparse_expected_cdf} are both of
order $(mp)^{-1/10}$.  The third term is polynomially smaller than the first two, since
$m^{-1}\log^{C_0}m/\eta\le m^{-9/10}\log^{C_0}m$.  Hence, if $mp\ge\log^{\alpha_0}m$ with any
fixed $\alpha_0>15$,
\begin{align}
\label{eq:supp_sparse_cdf_small}
    \sup_{x\ge0}
    |\mathbb E F_{\bm A}(x)-F_\rho(x)|
    =o(\delta_m^{3/2}).
\end{align}

We next pass from expectation to high probability without using any dense
random-matrix result.  Replacing one row of $\bm A$ changes
$\bm A'\bm A$ by a matrix of rank at most two.  Therefore the rank
inequality for empirical distribution functions shows that, for each fixed
$x$, $F_{\bm A}(x)$ changes by at most $2/n$ under one-row replacement.
Since the rows are independent, the bounded-difference inequality gives
\begin{align}
\label{eq:supp_sparse_cdf_concentration}
    \mathbb P\left(
    |F_{\bm A}(x)-\mathbb E F_{\bm A}(x)|>u\right)
    \le2\exp(-cnu^2).
\end{align}
Fix a sufficiently small constant $c_0>0$ and apply
\eqref{eq:supp_sparse_cdf_concentration} at the $O(n)$ deterministic points
$\gamma_i\pm C\delta_m$, with
$u=c_0\delta_m^{3/2}$.  A union bound together with
\eqref{eq:supp_sparse_cdf_small} yields, with probability at least
$1-cn^{-D}$,
\begin{align}
\label{eq:supp_sparse_cdf_grid}
    \max_{1\le i\le n}
    \max_{\varsigma\in\{-1,1\}}
    \left|
    F_{\bm A}(\gamma_i+\varsigma C\delta_m)
    -F_\rho(\gamma_i+\varsigma C\delta_m)
    \right|
    \le2c_0\delta_m^{3/2}.
\end{align}
Here $c_0$ is fixed first; the displacement constant $C$ will be chosen
large enough in Step 4.

\textit{Step 3: exclusion of sparse spectral outliers.}
The counting estimate \eqref{eq:supp_sparse_cdf_grid} alone cannot exclude a
small number of outliers, so we control the two edges separately using
sparse-efficient matrix concentration/universality.

Let $L=C_D\log(mn)$ and define
\begin{align*}
    \mu_{ij}
    &=\mathbb E\{\xi_{ij}\mathbf1(|\xi_{ij}|\le L)\},\\
    v_{ij}
    &=\operatorname{Var}\{\xi_{ij}\mathbf1(|\xi_{ij}|\le L)\},\\
    \widetilde\xi_{ij}
    &=\frac{\xi_{ij}\mathbf1(|\xi_{ij}|\le L)-\mu_{ij}}
      {\sqrt{v_{ij}}}.
\end{align*}
The sub-exponential tail implies
\begin{align}
\label{eq:supp_trunc_moments}
    \max_{i,j}\{|\mu_{ij}|+|v_{ij}-1|\}
    \le C L^2e^{-cL},
\end{align}
and
\begin{align}
\label{eq:supp_sparse_trunc_prob}
    \mathbb P\left(\max_{i,j}|\xi_{ij}|>L\right)
    \le2mn e^{-L/K}\le n^{-D-10}
\end{align}
when $C_D$ is large enough.  On the complement of the bad event in
\eqref{eq:supp_sparse_trunc_prob}, the matrices constructed from
$\xi_{ij}$ and $\widetilde\xi_{ij}$ differ in operator norm by
$o(n^{-D-5})$ after the normalization $(mp)^{-1/2}$; this follows from
\eqref{eq:supp_trunc_moments} and the Frobenius-norm bound.  Thus it suffices
for the edge analysis to consider
\begin{align*}
    \widetilde{\bm A}
    =\frac{1}{\sqrt{mp}}P_{\bOmega}(\widetilde{\bm\Xi}),
\end{align*}
whose entries are centered, variance one before masking, and bounded by
$CL$.

For the Hermitian dilation of $\widetilde{\bm A}$ we have
\begin{align*}
    \sigma\le1+o(1),\qquad
    \sigma_*\lesssim m^{-1/2},\qquad
    R\lesssim\frac{L}{\sqrt{mp}}.
\end{align*}
Corollary 2.7 of \cite{brailovskaya2024universality}, with
$t=C_D\log n$, therefore gives
\begin{align*}
    s_1(\widetilde{\bm A})
    \le \mathbb E s_1(\bm G_0)
    +C\left\{
      \sqrt{\frac{\log n}{m}}
      +\left(\frac{L}{\sqrt{mp}}\right)^{1/3}(\log n)^{2/3}
      +\frac{L\log n}{\sqrt{mp}}
    \right\}
\end{align*}
with probability at least $1-cn^{-D}$.  For the dense Gaussian comparator,
the standard Gaussian comparison inequality gives
$\mathbb E s_1(\bm G_0)\le1+\sqrt\rho$.  Hence, if
$mp\ge\log^{\alpha_0}m$ with $\alpha_0>15$, the bracketed term is
$o(\delta_m)$ and
\begin{align}
\label{eq:supp_sparse_top_edge}
    s_1(\bm A)\le1+\sqrt\rho+C_{\rm e}\delta_m
\end{align}
with probability at least $1-cn^{-D}$.

When $\rho<1$, apply Theorem 3.16 of
\cite{brailovskaya2024universality} directly to the rectangular matrix
$\widetilde{\bm A}$.  In the present normalization,
\begin{align}
\label{eq:supp_rectangular_parameters}
    \sigma(\widetilde{\bm A})\le1+o(1),\qquad
    \sigma_*(\widetilde{\bm A})\lesssim m^{-1/2},\qquad
    v(\widetilde{\bm A})\lesssim m^{-1/2},\qquad
    R(\widetilde{\bm A})\lesssim\frac{L}{\sqrt{mp}}.
\end{align}
Moreover,
\begin{align*}
    \mathbb E\widetilde{\bm A}'\widetilde{\bm A}=\bm I_n,
    \qquad
    \mathbb E\widetilde{\bm A}\widetilde{\bm A}'=\rho\bm I_m.
\end{align*}
Hence the user-friendly free lower bound accompanying Theorem 3.16 gives
$1-\sqrt\rho$ as the leading lower-edge term.  With $t=C_D\log n$, the
remainder in that theorem is bounded by
\begin{align*}
    C\left\{
    m^{-1/4}(\log n)^{3/4}
    +\frac{\log n}{\sqrt m}
    +\left(\frac{L}{\sqrt{mp}}\right)^{1/3}(\log n)^{2/3}
    +\frac{L\log n}{\sqrt{mp}}
    \right\}
    =o(\delta_m).
\end{align*}
Consequently,
\begin{align}
\label{eq:supp_sparse_bottom_edge}
    s_n(\bm A)\ge1-\sqrt\rho-C\delta_m
\end{align}
with probability at least $1-cn^{-D}$.  When $\rho=1$, the lower edge of the
rooted Marchenko--Pastur law is zero, and the outward lower-edge bound follows
simply from $s_n(\bm A)\ge0$.

\textit{Step 4: deterministic inversion of the spectral counts.}
The rooted Marchenko--Pastur law satisfies the one-sided mass lower bounds
needed for quantile inversion.  Under \eqref{eq:supp_aspect_ratio}, there are
constants $c_\rho,t_\rho>0$ such that, uniformly in
$\rho\in[\rho_-,1]$, for $0<t\le t_\rho$,
\begin{align}
\label{eq:supp_mp_mass_lower}
    \int_x^{x+t}\nu_\rho(u)\,du
    &\ge c_\rho t^{3/2},
    &&\lambda_-\le x<x+t\le\lambda_+,\notag\\
    \int_{x-t}^{x}\nu_\rho(u)\,du
    &\ge c_\rho t^{3/2},
    &&\lambda_-\le x-t<x\le\lambda_+.
\end{align}
Indeed, this follows directly from the explicit density
\eqref{eq:supp_rooted_mp_density}: at either regular edge the one-sided mass
is of order $t^{3/2}$, while in the bulk it is of order $t$.  When
$\rho=1$, the rooted density is bounded away from zero near its lower edge,
so the same lower bound remains valid uniformly.

We next invert the spectral counts.  By
\eqref{eq:supp_mp_quantiles},
\begin{align*}
    F_\rho(\gamma_i)=1-\frac{i-1/2}{n}.
\end{align*}
Suppose first that
$s_i(\bm A)>\gamma_i+C\delta_m$ and that
$\gamma_i+C\delta_m\le\lambda_+$.  At
$x_i^+=\gamma_i+C\delta_m$, at least $i$ singular values are larger than
$x_i^+$, so
\begin{align*}
    F_{\bm A}(x_i^+)\le1-\frac{i}{n}.
\end{align*}
On the other hand, \eqref{eq:supp_mp_mass_lower} gives
\begin{align*}
    F_\rho(x_i^+)
    &\ge1-\frac{i-1/2}{n}+c_\rho(C\delta_m)^{3/2}\\
    &=1-\frac{i}{n}+\frac{1}{2n}
      +c_\rho C^{3/2}\delta_m^{3/2}.
\end{align*}
Thus
\begin{align*}
    |F_{\bm A}(x_i^+)-F_\rho(x_i^+)|
    \ge c_\rho C^{3/2}\delta_m^{3/2},
\end{align*}
which contradicts \eqref{eq:supp_sparse_cdf_grid} once $C$ is chosen large
enough relative to the fixed $c_0$.  If
$\gamma_i+C\delta_m>\lambda_+$, then
$\gamma_i>\lambda_+-C\delta_m$, while
\eqref{eq:supp_sparse_top_edge} gives
$s_i(\bm A)\le\lambda_++C_{\rm e}\delta_m$.  Therefore
$s_i(\bm A)-\gamma_i\le(C+C_{\rm e})\delta_m$.  Thus the same final
$O(\delta_m)$ conclusion holds after enlarging the universal constant.

Similarly, suppose that
$s_i(\bm A)<\gamma_i-C\delta_m$ and
$\gamma_i-C\delta_m\ge\lambda_-$.  At
$x_i^-=\gamma_i-C\delta_m$ we have
\begin{align*}
    F_{\bm A}(x_i^-)\ge1-\frac{i-1}{n},
\end{align*}
whereas
\begin{align*}
    F_\rho(x_i^-)
    \le1-\frac{i-1/2}{n}-c_\rho C^{3/2}\delta_m^{3/2}.
\end{align*}
This again contradicts \eqref{eq:supp_sparse_cdf_grid}.  If
$x_i^-<\lambda_-$ and $\rho<1$, then
$\gamma_i<\lambda_-+C\delta_m$, while
\eqref{eq:supp_sparse_bottom_edge} gives
$s_i(\bm A)\ge\lambda_--C_{\rm e}\delta_m$ after enlarging
$C_{\rm e}$ if necessary.  Consequently
$\gamma_i-s_i(\bm A)\le(C+C_{\rm e})\delta_m$.  When $\rho=1$, the
outward lower edge is zero and the same conclusion follows from
$s_i(\bm A)\ge0$.  We conclude that
\begin{align*}
    \max_{1\le i\le n}|s_i(\bm A)-\gamma_i|
    \le C\delta_m,
\end{align*}
which proves \eqref{eq:supp_sparse_coarse_rigidity} and
\eqref{eq:supp_sparse_edge_bound}.
\end{proof}

Let $\bm G^{(t)}=(g^{(t)}_{ij})_{m\times n}$, $1\le t\le T$, be independent
standard Gaussian matrices, independent of $(\bOmega,\bH)$, and define the
ideal same-mask reference
\begin{align}
\label{eq:supp_ideal_reference}
    \lambda_i^{\circ}(\bOmega)
    =\mathbb E_{\bm G}\left[
      s_i\left(\frac{P_{\bOmega}(\bm G)}{\sqrt m}\right)
      \biggm|\bOmega\right],
    \qquad1\le i\le n,
\end{align}
and its Monte Carlo approximation
\begin{align}
\label{eq:supp_mc_reference}
    \widehat\lambda_i
    =\frac1T\sum_{t=1}^T
      s_i\left(\frac{P_{\bOmega}(\bm G^{(t)})}{\sqrt m}\right).
\end{align}

\begin{slemma}[Same-mask spectral and scale calibration]
\label{lem:loss_for_true}
Under the homogeneous-noise assumptions and the sampling condition of
Lemma~\ref{lem:sparse_coarse_rigidity}, the following bounds hold with
probability at least $1-cn^{-3}$:
\begin{align}
 \max_i\left|\frac{s_i(P_{\bOmega}\bH)}{\sigma\sqrt m}
                  -\widehat\lambda_i\right|
 &\le C\sqrt p\log^{-1}m,\label{eq:supp_noise_reference_uniform}\\
 \frac{\|P_{\bOmega}\bH\|}{\sigma\sqrt m}
       \vee\widehat\lambda_1
 &\le\sqrt p(1+\sqrt\rho+C\log^{-1}m),
       \label{eq:control_theta_1}\\
 \sum_{i\in I}\widehat\lambda_i&\ge cn\sqrt p,
       \label{eq:bound_for_lambda_sum}\\
 \left|\frac{\widehat\sigma(\M_0)}{\sigma}-1\right|
 &\le C\log^{-1}m,\label{eq:supp_sigma_consistency}\\
 \mathcal L(\M_0;\Y,\bOmega)&\le C_0\log^{-2}m.
       \label{eq:calibration_true_loss}
\end{align}
On the same event, simultaneously for all $\M\in\mathbb R^{m\times n}$,
\begin{align}
 \left|\frac{\widehat\sigma(\M)}{\sigma}-1\right|
 \le C\log^{-1}m+
 \frac{C\|P_{\bOmega}(\M-\M_0)\|_F}{\sigma\sqrt{mnp}}.
 \label{eq:working_scale_bound}
\end{align}
The conditional reference and Monte Carlo bounds are
\eqref{eq:supp_conditional_reference_rigidity} and
\eqref{eq:supp_mc_reference_error}, respectively.
\end{slemma}
\begin{proof}[Proof.]
\textit{Same-mask reference.} For every fixed $D>0$, if $mp\ge\log^{\alpha_0}m$ with $\alpha_0$
sufficiently large, then with probability at least $1-cn^{-D}$ over the
mask,
\begin{align}
\label{eq:supp_conditional_reference_rigidity}
    \max_{1\le i\le n}
    |\lambda_i^{\circ}(\bOmega)-\sqrt p\,\gamma_i|
    \le C\sqrt p\,\log^{-1}m.
\end{align}
Conditionally on such a mask, with probability at least $1-cn^{-D}$ over the
Monte Carlo Gaussian matrices,
\begin{align}
\label{eq:supp_mc_reference_error}
    \max_{1\le i\le n}
    |\widehat\lambda_i-\lambda_i^{\circ}(\bOmega)|
    \le C\sqrt{\frac{\log n}{mT}}.
\end{align}

Apply Lemma \ref{lem:sparse_coarse_rigidity} to the normalized sparse
Gaussian matrix
\begin{align*}
    \bm B_{\bOmega}(\bm G)
    =\frac{P_{\bOmega}(\bm G)}{\sqrt{mp}}.
\end{align*}
For an arbitrarily large $D_1$,
\begin{align}
\label{eq:supp_joint_gaussian_bad}
    \mathbb P_{\bOmega,\bm G}\left(
    \max_i|s_i(\bm B_{\bOmega}(\bm G))-\gamma_i|
    >C\delta_m\right)
    \le cn^{-D_1}.
\end{align}
Let $q(\bOmega)$ be the conditional probability of the bad event in
\eqref{eq:supp_joint_gaussian_bad} given $\bOmega$.  Fubini's theorem and
Markov's inequality give
\begin{align}
\label{eq:supp_q_mask}
    \mathbb P_{\bOmega}\{q(\bOmega)>n^{-D_1/2}\}
    \le cn^{-D_1/2}.
\end{align}
A Chernoff bound also gives
\begin{align}
\label{eq:supp_mask_size_event}
    \mathbb P_{\bOmega}\{|\bOmega|>2mnp\}
    \le \exp(-cmnp).
\end{align}
Hence, on $|\bOmega|\le2mnp$, for a fixed mask,
\begin{align*}
    \mathbb E_{\bm G}\|\bm B_{\bOmega}(\bm G)\|_F^2
    =\frac{|\bOmega|}{mp}\le2n.
\end{align*}
Therefore, on the intersection of $|\bOmega|\le2mnp$ with
$q(\bOmega)\le n^{-D_1/2}$, Cauchy--Schwarz gives
for every $i$,
\begin{align*}
    &\left|
    \mathbb E_{\bm G}
      s_i(\bm B_{\bOmega}(\bm G))-\gamma_i
    \right|\\
    &\qquad\le C\delta_m
      +\left\{
      \mathbb E_{\bm G}
      (\|\bm B_{\bOmega}(\bm G)\|_F+\gamma_1)^2
      \right\}^{1/2}q(\bOmega)^{1/2}\\
    &\qquad\le C\delta_m
      +C\sqrt n\,q(\bOmega)^{1/2}.
\end{align*}
After multiplying by $\sqrt p$, the second term is
$C\sqrt{np}\,q(\bOmega)^{1/2}$.  Relative to the desired scale
$\sqrt p\,\delta_m$, its ratio is
$C\sqrt n\,q(\bOmega)^{1/2}/\delta_m$.  Since $D_1$ can be chosen
arbitrarily large, \eqref{eq:supp_q_mask} and
\eqref{eq:supp_mask_size_event} imply that this ratio is $o(1)$ on an event
of probability at least $1-cn^{-D}$.  This proves
\eqref{eq:supp_conditional_reference_rigidity}.  The conditional exceptional set is controlled by its Frobenius second moment.

For the Monte Carlo error, condition on $\bOmega$.  The singular-value
Lipschitz inequality gives, for every fixed $i$,
\begin{align*}
    |s_i(P_{\bOmega}(\bm G_1)/\sqrt m)
      -s_i(P_{\bOmega}(\bm G_2)/\sqrt m)|
    \le m^{-1/2}\|\bm G_1-\bm G_2\|_F.
\end{align*}
Hence the average of $T$ independent copies is
$(mT)^{-1/2}$-Lipschitz in the joint Gaussian vector.  Gaussian
concentration followed by a union bound over $i=1,\ldots,n$ gives
\eqref{eq:supp_mc_reference_error}.  In particular,
\begin{align}
\label{eq:supp_mc_absorb}
    \sqrt{\frac{\log n}{mT}}
    \le C\sqrt p\,\log^{-1}m
\end{align}
whenever $mpT\gtrsim\log^3m$, which is automatic under the above
polylogarithmic sparsity assumption for every $T\ge1$.

\textit{Noise spectrum and bulk scale.} Put
$\theta_i=s_i(P_{\bOmega}\bH)/(\sigma\sqrt m)$ and $\xi_{ij}=H_{ij}/\sigma$.  By the standing noise assumption, the $\xi_{ij}$ satisfy
\eqref{eq:supp_sparse_entry_assumption}.  Lemma
\ref{lem:sparse_coarse_rigidity} therefore yields
\begin{align}
\label{eq:supp_noise_to_mp}
    \max_i\left|
    \frac{\theta_i}{\sqrt p}-\gamma_i
    \right|\le C\log^{-1}m
\end{align}
with probability at least $1-cn^{-4}$.  Equations~\eqref{eq:supp_conditional_reference_rigidity},
\eqref{eq:supp_mc_reference_error}, and \eqref{eq:supp_mc_absorb} give, on another event of probability at least
$1-cn^{-4}$,
\begin{align}
\label{eq:supp_reference_to_mp}
    \max_i\left|
    \frac{\widehat\lambda_i}{\sqrt p}-\gamma_i
    \right|\le C\log^{-1}m.
\end{align}
The triangle inequality proves
\eqref{eq:supp_noise_reference_uniform}.  Since
$\gamma_1\le1+\sqrt\rho$, the same two displays imply
\eqref{eq:control_theta_1}.

By definition,
\begin{align}
\label{eq:supp_sigma_M0_ratio}
    \frac{\widehat\sigma(\M_0)}{\sigma}-1
    =\frac{\sum_{i=n/3}^{2n/3}
      (\theta_i-\widehat\lambda_i)}
    {\sum_{i=n/3}^{2n/3}\widehat\lambda_i}.
\end{align}
Equation~\eqref{eq:supp_noise_reference_uniform} gives
\begin{align}
\label{eq:supp_sigma_M0_numerator}
    \left|\sum_{i=n/3}^{2n/3}
      (\theta_i-\widehat\lambda_i)\right|
    \lesssim n\sqrt p\,\log^{-1}m.
\end{align}
Under \eqref{eq:supp_aspect_ratio}, the middle third of the rooted
Marchenko--Pastur quantiles is uniformly bounded away from zero.  Hence
equation~\eqref{eq:supp_reference_to_mp} implies
\begin{align*}
    \sum_{i=n/3}^{2n/3}\widehat\lambda_i
    \ge\sqrt p\sum_{i=n/3}^{2n/3}
      (\gamma_i-C\log^{-1}m)
    \gtrsim n\sqrt p
\end{align*}
with probability at least $1-cn^{-3}$.  Combining
\eqref{eq:supp_sigma_M0_ratio}, \eqref{eq:supp_sigma_M0_numerator}, and \eqref{eq:bound_for_lambda_sum} proves
\eqref{eq:supp_sigma_consistency}.

\textit{Candidate-dependent scale.} Let
\begin{align*}
    \Delta=\M-\M_0,\qquad
    \bm E(\M)=\Y-P_{\bOmega}(\M)
    =P_{\bOmega}(\bH-\Delta).
\end{align*}
By Mirsky's inequality,
\begin{align*}
    \sum_{i=1}^n
    \left\{s_i(\bm E(\M))-s_i(P_{\bOmega}(\bH))\right\}^2
    \le \|P_{\bOmega}(\Delta)\|_F^2.
\end{align*}
Consequently,
\begin{align}
\label{eq:mirsky_middle_sum}
    \left|
    \sum_{i\in I}s_i(\bm E(\M))
    -\sum_{i\in I}s_i(P_{\bOmega}(\bH))
    \right|
    \le C\sqrt n\,\|P_{\bOmega}(\Delta)\|_F.
\end{align}
The comparison already established in \eqref{eq:supp_noise_reference_uniform} gives
\begin{align*}
    \max_i\left|
    \frac{s_i(P_{\bOmega}(\bH))}{\sigma\sqrt m}
    -\widehat\lambda_i
    \right|
    \le C\sqrt p\log^{-1}m,
\end{align*}
while \eqref{eq:bound_for_lambda_sum} gives
$\sum_{i\in I}\widehat\lambda_i\ge cn\sqrt p$.
Using the definition
\begin{align*}
    \widehat\sigma(\M)
    =\frac{\sum_{i\in I}s_i(\bm E(\M))}
    {\sqrt m\sum_{i\in I}\widehat\lambda_i},
\end{align*}
we obtain from \eqref{eq:mirsky_middle_sum}
\begin{align*}
    |\widehat\sigma(\M)-\sigma|
    &\le C\sigma\log^{-1}m
      +C(mnp)^{-1/2}\|P_{\bOmega}(\Delta)\|_F,
\end{align*}
which proves \eqref{eq:working_scale_bound}.

\textit{Loss at the truth.} At $\M=\M_0$,
\begin{align*}
    \mathcal L(\M_0;\Y,\bOmega)
    &=p^{-1}\sum_{i=1}^n\omega_i
      \left(
      \frac{\sigma}{\widehat\sigma(\M_0)}\theta_i
      -\widehat\lambda_i
      \right)^2\\
    &\le2p^{-1}\sum_{i=1}^n\omega_i
      (\theta_i-\widehat\lambda_i)^2
      +2p^{-1}\left(
      \frac{\sigma}{\widehat\sigma(\M_0)}-1
      \right)^2
      \sum_{i=1}^n\omega_i\theta_i^2\\
    &\le2p^{-1}\max_i
      |\theta_i-\widehat\lambda_i|^2
      +2p^{-1}\theta_1^2
      \left(
      \frac{\sigma}{\widehat\sigma(\M_0)}-1
      \right)^2.
\end{align*}
By \eqref{eq:supp_noise_reference_uniform}, \eqref{eq:supp_sigma_consistency},
and \eqref{eq:control_theta_1}, both terms on the right-hand side are
$O(\log^{-2}m)$.  Therefore, with probability at least $1-cn^{-3}$,
\begin{align*}
    \mathcal L(\M_0;\Y,\bOmega)
    \le C_0\log^{-2}m.
\end{align*}
\end{proof}
\end{revision}
% END MAJOR REVISION S02


% BEGIN MAJOR REVISION S03
\begin{revision}
\section{Proofs for rank-constrained estimation}
\label{sec:proof_factor}

Throughout this section, write
\begin{align*}
    \Delta=\M-\M_0,\qquad
    d=\|\Delta\|_F,\qquad
    Q_{\Omega}(\Delta)=\|P_{\bOmega}(\Delta)\|_F^2,
\end{align*}
and let
\begin{align*}
    \bar\mu=\max\{\mu_0,\mu_1\}.
\end{align*}
We work on the high-probability event constructed in Section~\ref{sec:supp_sparse_rmt}
for the sparse spectral comparison at the true matrix; in particular,
\begin{align}
\label{eq:factor_noise_edge_event}
    \|P_{\bOmega}(\bH)\|
    &\le C\sigma\sqrt{mp},\\
\label{eq:factor_reference_edge_event}
    \widehat\lambda_1&\le C\sqrt p,\\
\label{eq:factor_reference_bulk_event}
    \sum_{i=n/3}^{2n/3}\widehat\lambda_i&\ge c n\sqrt p,
\end{align}
with probability at least $1-cn^{-3}$.  Constants may depend on the fixed lower bound for $n/m$.

We use three uniform sampling bounds before proving
Theorem~\ref{main-thm:upper_bound_for_rankr} and its corollaries.

\subsection{Uniform sampling bounds}
\label{sec:uniform_sampling}
\begin{slemma}[Uniform Bernoulli sampling bounds]
\label{lem:mask_geometry}
Let $J$ denote the $m\times n$ all-ones matrix, and put
$W_\Omega=p^{-1}\bOmega-J$ and $T_\Omega=p^{-1}P_{\bOmega}-\operatorname{Id}$.
Under the standing sparse sampling condition and $np\ge C\mu_0r\log m$,
with probability at least $1-cn^{-4}$,
\begin{align}
 \|W_\Omega\|&\le\Lambda:=C\sqrt{m/p},\label{eq:sampling_mask_norm}\\
 |\langle T_\Omega A,B\rangle|
 &\le\theta\|A\|_F\|B\|_F,\qquad A,B\in T_\star,
 \quad\theta=C\sqrt{\frac{\mu_0r\log m}{np}},
 \label{eq:recovery_tangent-event}
\end{align}
where $T_\star$ is the fixed tangent space at $\M_0$.
On this event the following hold for all conformable factors, including
data-dependent factors:
\begin{align}
 |\langle T_\Omega(AB'),CD'\rangle|
 &\le\Lambda b(A,C)b(B,D),
 \quad b(A,C)=\min\{\|A\|_F\|C\|_{2,\infty},
                            \|A\|_{2,\infty}\|C\|_F\},
 \label{eq:recovery_bilinear-event}\\
 \|T_\Omega(LR')\|
 &\le\Lambda\|L\|_{2,\infty}\|R\|_{2,\infty}.
 \label{eq:sampling_schur_bound}
\end{align}
The two pairs of factors in the bilinear inequality may have different
numbers of columns.
\end{slemma}
\begin{proof}[Proof.]
For the centered Bernoulli matrix divided by $\sqrt{mp}$, the Hermitian
summands have norm at most $(mp)^{-1/2}$, variance parameter at most one,
and weak variance parameter at most $Cm^{-1/2}$.  Its covariance-matched
Gaussian matrix has independent entries of variance $(1-p)/m$ and expected
operator norm at most two.  The bounded-summand edge comparison used in
Section~\ref{sec:supp_sparse_rmt}, with $t=C\log n$, bounds its norm by
\[
 2+C\left\{\sqrt{\log n/m}+(mp)^{-1/6}(\log n)^{2/3}
                  +(mp)^{-1/2}\log n\right\}\le C.
\]
This proves \eqref{eq:sampling_mask_norm}; for $p=1$ the centered matrix is zero.
For \eqref{eq:recovery_tangent-event}, let $z_{ij}=P_{T_\star}(e_ie_j')$.
Then $\|z_{ij}\|_F^2\le\mu_0r/m+\mu_0r/n$ and
$\sum_{ij}z_{ij}\otimes z_{ij}=\mathcal I_{T_\star}$.
Apply matrix Bernstein \citep{tropp2012userfriendly} to
$\sum_{ij}(\Omega_{ij}/p-1)z_{ij}\otimes z_{ij}$.
Its summand and variance bounds are $C\mu_0r/(np)$, and its dimension
is at most $r(m+n)$.  The resulting bound is the displayed $\theta$
plus $C\mu_0r\log m/(np)$, which is absorbed by $\theta$.

To prove the bilinear bound, form matrices whose $i$th and $j$th rows are
$A_{i,:}\otimes C_{i,:}$ and $B_{j,:}\otimes D_{j,:}$.
Their Frobenius norms are at most $b(A,C)$ and
$b(B,D)$, respectively.  Pairing them through $W_\Omega$
and using its operator norm gives \eqref{eq:recovery_bilinear-event}.
For \eqref{eq:sampling_schur_bound}, define the block-diagonal maps
$D_Lx=(x_iL_{i,:})_i$ and $D_Ry=(y_jR_{j,:})_j$.
The Schur product $W_\Omega\odot(LR')$ is
$D_L'(W_\Omega\otimes I)D_R$.
Since $\|D_L\|=\max_i\|L_{i,:}\|_2$ and similarly for $D_R$, the claim
follows.  The bound holds without summing the individual singular directions.
\end{proof}

\begin{slemma}[Relative sampling for incoherent candidates]
\label{lem:global_relative_sampling}
Let $C_M\ge1$, possibly depending on the dimensions.  On the event of
Lemma~\ref{lem:mask_geometry}, simultaneously for every
$\M\in\mathbb M_{r,\mu_1}$ satisfying $\|\M\|\le C_Ms_1$, with
$\Delta=\M-\M_0$ and $f=\|\Delta\|_F$,
\begin{align}
 \left|p^{-1}\|P_{\bOmega}\Delta\|_F^2-f^2\right|
 &\le a_g f^2,\label{eq:sampling_global_relative_sampling}\\
 a_g&=C\kappa\left\{\theta+(1+C_M)
                   \Lambda\frac{\bar\mu r}{\sqrt{mn}}\right\}
 \le C_{C_M}\kappa\left\{
      \sqrt{\frac{\bar\mu r\log m}{np}}+
      \frac{\bar\mu r}{\sqrt{np}}\right\}.
 \label{eq:sampling_global_sampling_constant}
\end{align}
In particular, for fixed $C_M$, if
\begin{align}
 np\ge C_{C_M}\kappa^2\{\bar\mu r\log m+\bar\mu^2r^2\},
 \label{eq:relative_sampling_condition}
\end{align}
then, on the same event,
\begin{align}
 \frac p2\|\M-\M_0\|_F^2
 \le\|P_{\bOmega}(\M-\M_0)\|_F^2
 \le\frac{3p}2\|\M-\M_0\|_F^2.
 \label{eq:factor_local_rsc}
\end{align}
\end{slemma}
\begin{proof}[Proof.]
Take balanced factors $\M=XY'$ and $\M_0=X_0Y_0'$, padding the first
pair with zero columns when its rank is smaller than $r$.  A common
orthogonal rotation of $X,Y$ can be chosen so that, with
$E_X=X-X_0$ and $E_Y=Y-Y_0$,
\begin{align}
 \|E_X\|_F^2+\|E_Y\|_F^2
 \le\frac{f^2}{(\sqrt2-1)s_\star}.
 \label{eq:sampling_global_factor_distance}
\end{align}
To prove this inequality, for $\operatorname{Dil}(M)=\left(\begin{smallmatrix}0&M\\M'&0\end{smallmatrix}\right)$
and $F=(X',Y')'$, balancedness gives
$FF'=|\operatorname{Dil}(M)|+\operatorname{Dil}(M)$.
For any symmetric $A,B$, expansion in their eigenbases shows
\[
 \||A|-|B|\|_F^2
 =\sum_{ij}(|a_i|-|b_j|)^2|\langle u_i,v_j\rangle|^2
 \le\sum_{ij}(a_i-b_j)^2|\langle u_i,v_j\rangle|^2
 =\|A-B\|_F^2.
\]
The diagonal and off-diagonal blocks of the balanced lift are orthogonal,
so $\|FF'-F_0F_0'\|_F^2\le4f^2$.
Apply the global Gram-factor inequality
\citep[Lemma~5.4]{tu2016procrustes}
\[
 \min_{O'O=I}\|FO-F_0\|_F^2
 \le\frac{\|FF'-F_0F_0'\|_F^2}
 {2(\sqrt2-1)\sigma_r(F_0)^2},
 \qquad \sigma_r(F_0)^2=2s_\star,
\]
which proves \eqref{eq:sampling_global_factor_distance}.

Let $b_L=\sqrt{\bar\mu r/m}$ and $b_R=\sqrt{\bar\mu r/n}$.
Balancedness, endpoint incoherence, and the operator norm bound imply
\begin{align*}
 \|X_0\|_{2,\infty}+\|E_X\|_{2,\infty}&\le C_{C_M}\sqrt{s_1}\,b_L,\\
 \|Y_0\|_{2,\infty}+\|E_Y\|_{2,\infty}&\le C_{C_M}\sqrt{s_1}\,b_R.
\end{align*}
Write
$\Delta_T=X_0E_Y'+E_XY_0'\in T_\star$ and
$\Delta_Q=E_XE_Y'$.  By \eqref{eq:sampling_global_factor_distance},
$\|\Delta_T\|_F+\|\Delta_Q\|_F\le C_{C_M}\sqrt\kappa f$;
for the quadratic term use
$\|E_X\|\le(1+\sqrt{C_M})\sqrt{s_1}$.
More specifically, $\|\Delta_T\|_F\le C\sqrt\kappa f$ independently of $C_M$.
The tangent--tangent form is therefore at most $C\kappa\theta f^2$.
For the two tangent--quadratic forms use, respectively,
\begin{align*}
 &\Lambda\|X_0\|_{2,\infty}\|E_X\|_F
                  \|E_Y\|_F\|E_Y\|_{2,\infty},\\
 &\Lambda\|E_X\|_F\|E_X\|_{2,\infty}
                  \|Y_0\|_{2,\infty}\|E_Y\|_F.
\end{align*}
Each is at most $C_{C_M}\Lambda\kappa b_Lb_R f^2$.
The quadratic--quadratic form is bounded by
\[
 \Lambda\|E_X\|_F\|E_X\|_{2,\infty}
 \|E_Y\|_F\|E_Y\|_{2,\infty},
\]
with the same upper bound.
Each row-norm product introduces at most $(1+\sqrt{C_M})^2\le2(1+C_M)$.
The tangent--tangent bound is independent of $C_M$.
Adding the four forms proves \eqref{eq:sampling_global_relative_sampling} with the
explicit dependence stated in \eqref{eq:sampling_global_sampling_constant}.

Choosing the constant in \eqref{eq:relative_sampling_condition} so that $a_g\le1/2$ gives \eqref{eq:factor_local_rsc}.
\end{proof}

\begin{slemma}[Sampling with a spikiness remainder]
\label{lem:coarse_sampling_rsc}\label{lem:nuclear_effective_rsc}
Assume $m\ge n$, $n/m\ge\rho_->0$, and $mp\ge C\log n$.  With probability
at least $1-cn^{-4}$,
\begin{align}
\label{eq:nuclear_effective_rsc}
 \left|
 \|P_{\bOmega}(\bm A)\|_F^2
 -
 p\|\bm A\|_F^2
 \right|
 \le&
 \frac p8\|\bm A\|_F^2
 +
 C\sqrt{mp}\,
 \|\bm A\|_\infty\|\bm A\|_*
 \nonumber\\
 &+
 C\log n\,\|\bm A\|_\infty^2
\end{align}
simultaneously for all
$\bm A\in\mathbb R^{m\times n}$.
Consequently, for every matrix with rank at most $k\le n$,
\begin{align}
 \|P_{\bOmega}(\bm A)\|_F^2
 \ge\frac p2\|\bm A\|_F^2-Ckm\|\bm A\|_\infty^2.
 \label{eq:coarse_rsc}
\end{align}
\end{slemma}
\begin{proof}[Proof.]
By homogeneity it is enough to consider $\|\bm A\|_F=1$.
For
\[
a\in[(mn)^{-1/2},1],
\qquad
R\in[1,\sqrt n],
\]
define
\[
 \mathcal A(a,R)
 =
 \{
 \bm A:
 \|\bm A\|_F=1,\
 \|\bm A\|_\infty\le a,\
 \|\bm A\|_*\le R
 \}
\]
and
\[
 Z(a,R)
 =
 \sup_{\bm A\in\mathcal A(a,R)}
 \left|
 \sum_{i,j}
 (\Omega_{ij}-p)A_{ij}^2
 \right|.
\]
Let $\varepsilon_{ij}$ be independent Rademacher variables.  Symmetrization
and contraction for $x\mapsto x^2$ on $[-a,a]$ give
\begin{align*}
 \mathbb EZ(a,R)
 &\le
 Ca\,
 \mathbb E
 \sup_{\bm A\in\mathcal A(a,R)}
 \left|
 \sum_{i,j}
 \varepsilon_{ij}\Omega_{ij}A_{ij}
 \right|\\
 &\le
 CaR\,
 \mathbb E\|\bm E\odot\bOmega\|.
\end{align*}
Seginer's inequality~\citep{seginer2000expected} and the maximal
binomial-degree bound under $mp\ge C\log n$ yield
\[
 \mathbb E\|\bm E\odot\bOmega\|
 \le
 C\sqrt{mp}.
\]
Hence
\begin{align}
\label{eq:nuclear_effective_rsc_expectation}
 \mathbb EZ(a,R)
 \le
 CaR\sqrt{mp}.
\end{align}

For every matrix in $\mathcal A(a,R)$, the centered summands are bounded by
$a^2$ and have total variance at most $pa^2$.  Bousquet's inequality
therefore gives, with probability at least $1-e^{-t}$,
\begin{align*}
 Z(a,R)
 \le
 \mathbb EZ(a,R)
 +
 C\sqrt{
 \{pa^2+a^2\mathbb EZ(a,R)\}t
 }
 +
 Ca^2t.
\end{align*}
Using
$2\sqrt{xy}\le x+y$
twice, and taking $t=C\log n$, gives
\begin{align}
\label{eq:nuclear_effective_rsc_shell}
 Z(a,R)
 \le
 \frac p{16}
 +
 CaR\sqrt{mp}
 +
 Ca^2\log n
\end{align}
with probability at least $1-n^{-7}$.

Take dyadic grids in both $a$ and $R$.  Their joint cardinality is
$O(\log^2(mn))$, so a union bound preserves probability at least
$1-cn^{-5}$.  For an arbitrary normalized matrix choose the smallest dyadic
$a$ and $R$ above its actual max norm and nuclear norm.  Each selected value
is at most twice the corresponding actual quantity.  Weakening numerical
constants gives, for all normalized matrices,
\[
 \left|
 \|P_{\bOmega}(\bm A)\|_F^2-p
 \right|
 \le
 \frac p8
 +
 C\sqrt{mp}\,
 \|\bm A\|_\infty\|\bm A\|_*
 +
 C\log n\,\|\bm A\|_\infty^2.
\]
Restore homogeneity to prove
\eqref{eq:nuclear_effective_rsc}.
For the rank-$k$ consequence, use $\|\bm A\|_*\le\sqrt k\|\bm A\|_F$
and Young's inequality to bound the middle term in
\eqref{eq:nuclear_effective_rsc} by
$p\|\bm A\|_F^2/8+Ckm\|\bm A\|_\infty^2$.
Since $\log n\le Ckm$, the last term is absorbed into the same remainder.
Weakening the leading numerical constant gives \eqref{eq:coarse_rsc}.
\end{proof}

\subsection{Proofs of the estimation results}
\begin{proof}[Proof of Theorem~\ref{main-thm:upper_bound_for_rankr}.]
Work on the intersection of the calibration and sampling events above.
Fix $C_L>0$ and $\omega_1>0$, and consider a finite-loss
candidate $\M\in\mathbb M_{r,\mu_1}$ satisfying
\begin{align}
 \mathcal L(\M;\Y,\bOmega)\le C_L\omega_1.
 \label{eq:factor_low_loss_condition}
\end{align}
All subsequent bounds hold on this common event, uniformly in the candidate.

\textit{Observed energy.} We first prove
\begin{align}
 Q_\Omega(\Delta)
 \le C_{C_L}\sigma\sqrt{rmp}\,d
       +C_{C_L}\frac rn d^2.
 \label{eq:factor_energy_upper}
\end{align}
The first weighted singular-value term in the loss and
\eqref{eq:factor_low_loss_condition} imply
\begin{align*}
    \left|
    \frac{\|\bm E(\M)\|}{\sqrt m\widehat\sigma(\M)}
    -\widehat\lambda_1
    \right|
    \le \sqrt{C_Lp}.
\end{align*}
Together with \eqref{eq:factor_reference_edge_event},
\begin{align}
\label{eq:residual_operator_low_loss}
    \|\bm E(\M)\|
    \le C_{C_L}\widehat\sigma(\M)\sqrt{mp}.
\end{align}
Since $\bm E(\M)=P_{\bOmega}(\bH-\Delta)$, equations
\eqref{eq:factor_noise_edge_event} and
\eqref{eq:residual_operator_low_loss} yield
\begin{align}
\label{eq:sampled_delta_operator}
    \|P_{\bOmega}(\Delta)\|
    \le C_{C_L}
    \{\widehat\sigma(\M)+\sigma\}\sqrt{mp}.
\end{align}
Because $\operatorname{rank}(\Delta)\le2r$,
\begin{align*}
    Q_{\Omega}(\Delta)
    &=\langle P_{\bOmega}(\Delta),\Delta\rangle
    \le \|P_{\bOmega}(\Delta)\|\,\|\Delta\|_*
    \le \sqrt{2r}\|P_{\bOmega}(\Delta)\|d.
\end{align*}
Applying \eqref{eq:working_scale_bound} to \eqref{eq:sampled_delta_operator} gives
\begin{align*}
    Q_{\Omega}(\Delta)
    \le C_{C_L}\sigma\sqrt{rmp}\,d
       +C_{C_L}\sqrt{\frac rn}
        Q_{\Omega}(\Delta)^{1/2}d.
\end{align*}
Young's inequality absorbs one half of the last
$Q_{\Omega}(\Delta)$ term and proves
\eqref{eq:factor_energy_upper}.

\textit{Operator localization.} Put $\beta=\bar\mu r/\sqrt{np}$.
We next show that, if $C_{C_L}\beta\le1/2$,
\begin{align}
 \|\M-\M_0\|\le C_{C_L}h+C_{C_L}\beta s_1.
 \label{eq:factor_coarse_localization}
\end{align}
Let $N=p^{-1}P_{\bOmega}\bH$, $\Delta=\M-\M_0$, and
$E_\Omega=T_\Omega\Delta$.  Applying
\eqref{eq:sampling_schur_bound} to balanced factorizations of the two endpoints
separately gives
\begin{align}
 \|E_\Omega\|\le C\beta(s_1+\|\Delta\|).
 \label{eq:sampling_endpoint_mask_bound}
\end{align}
The residual divided by $p$ is $N-\Delta-E_\Omega$.
Since $\rank(\Delta)\le2r$ and $i\in I$ implies $i>2r$, the rank
perturbation singular-value inequality gives
$s_i(N-\Delta)\le s_{i-2r}(N)\le\|N\|\le Ch$.
Weyl's inequality and $\sum_{i\in I}\widehat\lambda_i/\sqrt p\ge cn$
therefore give
\[
 \widehat\sigma(\M)\sqrt{m/p}
 =\frac{\sum_{i\in I}s_i\{N-\Delta-E_\Omega\}}
        {\sum_{i\in I}\widehat\lambda_i/\sqrt p}
 \le C(h+\|E_\Omega\|).
\]
The first weighted term of the low-loss condition implies
$\|\bm E(\M)/p\|\le C_{C_L}
\widehat\sigma(\M)\sqrt{m/p}$.
Thus $\|\Delta\|\le\|N\|+\|E_\Omega\|+
\|\bm E(\M)/p\|\le C_{C_L}(h+\|E_\Omega\|)$.
Use \eqref{eq:sampling_endpoint_mask_bound} and absorb its last term.

\textit{Frobenius error and minimizers.} Write $d=\|\M-\M_0\|_F$ and $\delta=\|\M-\M_0\|$.
The sampling condition makes $C_{C_L}\beta\le1/2$.
If $h\ge s_1$, \eqref{eq:factor_coarse_localization} gives $\delta\le Ch$, and
$d\le\sqrt{2r}\delta\le C\sqrt r h$ already proves the claim.
If $h<s_1$, \eqref{eq:factor_coarse_localization} gives $\|\M\|\le C_{C_L}s_1$.
Use \eqref{eq:factor_local_rsc} with this fixed $C_M$ and combine
it with \eqref{eq:factor_energy_upper}:
\[
 \frac p2 d^2\le C_{C_L}\sigma\sqrt{rmp}\,d
                    +C_{C_L}\frac r n d^2.
\]
Since $np\ge C_{C_L}r$, absorb the last term and divide by
$pd$ when $d>0$.  This gives $d\le C_{C_L}\sigma\sqrt{mr/p}$;
$d=0$ is immediate.  All events are uniform, so the argument applies to
data-dependent candidates.

Choose, for instance, $C_L=1$.  Lemma~\ref{lem:loss_for_true}
shows that
\begin{align*}
    \mathcal L(\M_0;\Y,\bOmega)\le C_0\log^{-2}m.
\end{align*}
Since $\omega_1\gg\log^{-2}m$, for all sufficiently large $m$,
\begin{align*}
    \mathcal L(\widetilde\M;\Y,\bOmega)
    \le \mathcal L(\M_0;\Y,\bOmega)
    \le \omega_1.
\end{align*}
The uniform bound just proved therefore gives
\begin{align*}
    \|\widetilde\M-\M_0\|_F
    \le C\sigma\sqrt{\frac{mr}{p}}
\end{align*}
with probability at least $1-cn^{-3}$.

The same argument also calibrates the working scale at a low-loss candidate.
For the minimizer, define
\begin{align}
 \widehat\sigma_{\rm pil}=\widehat\sigma(\widetilde\M_{\rm fac}).
 \label{eq:nuclear_pilot_definition}
\end{align}
Substitution of the Frobenius bound into \eqref{eq:factor_energy_upper}
gives $\|P_{\bOmega}(\widetilde\M_{\rm fac}-\M_0)\|_F^2\le C\sigma^2mr$.
Hence \eqref{eq:working_scale_bound} yields
\begin{align}
 \left|\widehat\sigma_{\rm pil}/\sigma-1\right|
 \le C\{\log^{-1}m+\sqrt{r/(np)}\}.
 \label{eq:nuclear_pilot_consistency}
\end{align}
Together with \eqref{eq:supp_sigma_consistency}, this establishes the pilot event
\begin{align}
 \mathcal E_{\rm pil}
 =\left\{\frac12\sigma\le\widehat\sigma_{\rm pil}\le\frac32\sigma,
          \quad\frac12\sigma\le\widehat\sigma(\M_0)\le\frac32\sigma\right\}
 \label{eq:nuclear_pilot_event}
\end{align}
for sufficiently large dimensions. The same bounds hold for any candidate
satisfying the low-loss premise, with constants depending on $C_L$.
\end{proof}

\begin{proof}[Proof of Corollary~\ref{main-cor:large_M0}.]
The approximate-minimizer assertion in Corollary~\ref{main-cor:large_M0}
follows immediately from Theorem~\ref{main-thm:upper_bound_for_rankr}:
truth feasibility and $\mathcal L(M_0)\le\omega_1/2$ put every candidate
within $\omega_1/2$ of the infimum below $\omega_1$.
\end{proof}

\label{sec:approximate_rank_bounds}
\begin{proof}[Proof of Corollary~\ref{main-cor:mis_ranks}.]
Let the chosen rank cap be $s\ge r$, put $k=r+s\le n$,
and suppose $M\in\mathbb M_{s,\mu_1}$ has loss at most $C_L\omega_1$.
The argument giving \eqref{eq:factor_energy_upper} uses only
$\rank(M-M_0)\le k$, and hence gives
\begin{align*}
 Q_\Omega(\Delta)\le C\sigma\sqrt{kmp}\,d+C(k/n)d^2,
 \qquad d=\|\Delta\|_F.
\end{align*}
The two endpoint leverage bounds imply
\begin{align*}
 \|\Delta\|_\infty
 \le\frac{(\mu_0r+\mu_1s)\|M_0\|_F+\mu_1s\,d}{\sqrt{mn}}.
\end{align*}
Combine these inequalities with Lemma~\ref{lem:coarse_sampling_rsc}
at rank $k$.  If $np\ge Ck\{1+\mu_1^2s^2\}$, the terms proportional
to $d^2$ are absorbed, leaving
\begin{align*}
 d^2\le C\sigma\sqrt{mk/p}\,d
       +C\frac{k(\mu_0r+\mu_1s)^2}{np}\|M_0\|_F^2.
\end{align*}
Solving this quadratic proves Corollary~\ref{main-cor:mis_ranks}.
For a minimizer, its low-loss premise follows from truth feasibility and
$\omega_1\gg\log^{-2}m$.
\end{proof}
\end{revision}
% END MAJOR REVISION S03


% BEGIN MAJOR REVISION S04
\begin{revision}
\section{Proofs for nuclear-norm estimation}
\label{sec:nuclear_estimation}

We use the clipped loss and objective comparison of Section~\ref{main-sec:4.2}.
On the pilot event
\eqref{eq:nuclear_pilot_event}, fix
\begin{align}
\label{eq:nuclear_clip_definition}
 \widehat\sigma_{\rm cl}(\M)
 :=
 \min\left\{
 \widehat\sigma(\M),\,
 4\widehat\sigma_{\rm pil}
 \right\}.
\end{align}
Then, uniformly over all candidate matrices,
\begin{align}
\label{eq:nuclear_clip_uniform}
 \widehat\sigma_{\rm cl}(\M)
 \le
 6\sigma,
\end{align}
whereas at the truth clipping is inactive:
\begin{align}
\label{eq:nuclear_clip_inactive_truth}
 \widehat\sigma_{\rm cl}(\M_0)
 =
 \widehat\sigma(\M_0).
\end{align}
The numerical constant $4$ is inessential; any fixed constant that makes
the cap strictly larger than the true working scale on
$\mathcal E_{\rm pil}$ can be used.

Define the clipped spectral loss
\begin{align}
\label{eq:nuclear_clipped_loss}
 \mathcal L_{\rm cl}(\M;\Y,\bOmega)
 :=
 \frac1p
 \sum_{i=1}^n
 \omega_i
 \left\{
 \frac{
 s_i\{\Y-P_{\bOmega}(\M)\}
 }{
 \sqrt m\,\widehat\sigma_{\rm cl}(\M)
 }
 -
 \widehat\lambda_i
 \right\}^2.
\end{align}
Because of \eqref{eq:nuclear_clip_inactive_truth},
Lemma~\ref{lem:loss_for_true} gives
\begin{align}
\label{eq:nuclear_clipped_true_loss}
 \mathcal L_{\rm cl}(\M_0;\Y,\bOmega)
 \le
 C_0\log^{-2}m
\end{align}
on the same high-probability event.

Let $\widehat{\M}_\lambda$ satisfy the approximate clipped objective comparison
\begin{align}
\label{eq:nuclear_approx_objective}
 \mathcal L_{\rm cl}(\widehat{\M}_\lambda;\Y,\bOmega)
 +\lambda\|\widehat{\M}_\lambda\|_*
 \le
 \mathcal L_{\rm cl}(\M_0;\Y,\bOmega)
 +\lambda\|\M_0\|_*
 +\varepsilon_{\rm opt},
\end{align}
and write
\begin{align*}
 \Delta=\widehat{\M}_\lambda-\M_0,
 \qquad
 d=\|\Delta\|_F,
 \qquad
 Q_\Omega(\Delta)=\|P_{\bOmega}(\Delta)\|_F^2.
\end{align*}
Let $\M_0=\U\bm\Sigma\V'$ and define
\begin{align*}
 T
 =
 \{\U A'+B\V':
 A\in\mathbb R^{n\times r},\
 B\in\mathbb R^{m\times r}\},
\end{align*}
with orthogonal projections $P_T$ and
$P_{T^\perp}$.  For a matrix $\bm A$, define
\begin{align*}
 \|\bm A\|_{F,k}
 =
 \left(
 \sum_{i=1}^k s_i(\bm A)^2
 \right)^{1/2}.
\end{align*}
Finally put
\begin{align}
\label{eq:nuclear_ellopt}
 \ell_{\rm opt}
 :=
 \mathcal L_{\rm cl}(\M_0;\Y,\bOmega)
 +\varepsilon_{\rm opt},
 \qquad
 \tau_{\lambda,\rm opt}
 :=
 \frac{\ell_{\rm opt}}{\lambda}.
\end{align}

\begin{proof}[Proof of Theorem~\ref{main-thm:upper_bound_for_convexlambda}.]
\textit{Cone inequality.} The objective comparison and the truth-loss bound give
\begin{align}
 \|P_{T^\perp}\Delta\|_*
 &\le
 \|P_T\Delta\|_*
 +\tau_{\lambda,\rm opt},
 \label{eq:nuclear_cone_slack}\\
 \|\Delta\|_*
 &\le
 2\sqrt{2r}\,\|\Delta\|_F
 +\tau_{\lambda,\rm opt},
 \label{eq:nuclear_effective_rank}
\end{align}
and
\begin{align}
 \mathcal L_{\rm cl}(\widehat{\M}_\lambda;\Y,\bOmega)
 \le
 \ell_{\rm opt}
 +\lambda\sqrt{2r}\,\|\Delta\|_F.
 \label{eq:nuclear_loss_upper}
\end{align}
Moreover,
\begin{align}
\label{eq:nuclear_ellopt_bound}
 \ell_{\rm opt}
 \le
 C_0\log^{-2}m+\varepsilon_{\rm opt}.
\end{align}

From \eqref{eq:nuclear_approx_objective} and the decomposability inequality
\begin{align*}
 \|\M_0+\Delta\|_*
 \ge
 \|\M_0\|_*
 +
 \|P_{T^\perp}\Delta\|_*
 -
 \|P_T\Delta\|_*,
\end{align*}
and using nonnegativity of the clipped loss, we obtain
\eqref{eq:nuclear_cone_slack}.  Since
$\operatorname{rank}(P_T\Delta)\le2r$,
\begin{align*}
 \|P_T\Delta\|_*
 \le
 \sqrt{2r}\|P_T\Delta\|_F
 \le
 \sqrt{2r}\|\Delta\|_F,
\end{align*}
which gives \eqref{eq:nuclear_effective_rank}.  Dropping the favorable
$-\lambda\|P_{T^\perp}\Delta\|_*$ term in the same objective
comparison gives \eqref{eq:nuclear_loss_upper}.  Finally,
\eqref{eq:nuclear_ellopt_bound} follows from
\eqref{eq:nuclear_clipped_true_loss}.

\textit{Residual spectrum and sampled energy.} Let
\begin{align}
\label{eq:nuclear_omega_r}
 \bar\omega_r
 =
 \min_{1\le i\le r}\omega_i
 >0.
\end{align}
On the calibration and pilot events, for any matrix $\M$ with residual
$\bm R=\Y-P_{\bOmega}(\M)$,
\begin{align}
 \|\bm R\|_{F,r}
 \le
 C\sigma\sqrt{rmp}
 +
 C\sigma
 \sqrt{
 \frac{pm}{\bar\omega_r}
 \mathcal L_{\rm cl}(\M;\Y,\bOmega)
 }.
 \label{eq:nuclear_residual_topr}
\end{align}
Consequently, for
$\bm A=P_{\bOmega}(\M-\M_0)$,
\begin{align}
 \|\bm A\|_{F,r}
 \le
 C\sigma\sqrt{rmp}
 +
 C\sigma
 \sqrt{
 \frac{pm}{\bar\omega_r}
 \mathcal L_{\rm cl}(\M;\Y,\bOmega)
 }.
 \label{eq:nuclear_sampled_topr}
\end{align}

The clipped loss definition gives
\begin{align*}
 \sum_{i=1}^r
 \left\{
 s_i(\bm R)
 -
 \sqrt m\,\widehat\sigma_{\rm cl}(\M)
 \widehat\lambda_i
 \right\}^2
 \le
 \frac{
 pm\widehat\sigma_{\rm cl}(\M)^2
 }{\bar\omega_r}
 \mathcal L_{\rm cl}(\M;\Y,\bOmega).
\end{align*}
On $\mathcal E_{\rm pil}$,
$\widehat\sigma_{\rm cl}(\M)\le6\sigma$ uniformly, while the sparse spectral
event gives $\max_i\widehat\lambda_i\le C\sqrt p$.  This proves
\eqref{eq:nuclear_residual_topr}.  Write
$\bm E=P_{\bOmega}(\bH)$ and
$\bm A=\bm E-\bm R$.  Since the Ky Fan--Frobenius norm is unitarily
invariant,
\[
 \|\bm A\|_{F,r}
 \le
 \|\bm E\|_{F,r}+\|\bm R\|_{F,r},
\]
and the masked-noise event gives
$\|\bm E\|_{F,r}\le C\sigma\sqrt{rmp}$.

We next establish
\begin{align}
 Q_\Omega(\Delta)
 \le
 C\,
 a_{\lambda,\rm opt}(d)
 \left(
 d+\frac{\tau_{\lambda,\rm opt}}{\sqrt r}
 \right),
 \label{eq:nuclear_Q_upper}
\end{align}
where
\begin{align}
 a_{\lambda,\rm opt}(d)
 :=
 C\sigma\sqrt{rmp}
 +
 C\sigma
 \sqrt{
 \frac{pm}{\bar\omega_r}
 \left(
 \ell_{\rm opt}
 +\lambda\sqrt{2r}\,d
 \right)
 }.
 \label{eq:nuclear_a_lambda}
\end{align}

Let
\[
 \bm A=P_{\bOmega}(\Delta).
\]
Use the orthogonal tangent decomposition
\begin{align}
\label{eq:nuclear_tangent_split}
 P_T\Delta
 =
 \underbrace{
 \U\U'\Delta
 }_{\Delta_{T,1}}
 +
 \underbrace{
 (\operatorname{Id}-\U\U')\Delta\V\V'
 }_{\Delta_{T,2}},
\end{align}
where both matrices have rank at most $r$ and are Frobenius-orthogonal.
Write
\[
 P_{T^\perp}\Delta
 =
 \Delta_1+\Delta_2+\cdots
\]
by grouping its singular values into consecutive blocks of size $r$.
Every $\Delta_j$ has rank at most $r$.

For every rank-$r$ matrix $\bm B$, von Neumann's trace inequality gives
\[
 |\langle\bm A,\bm B\rangle|
 \le
 \|\bm A\|_{F,r}\|\bm B\|_F.
\]
Hence
\begin{align*}
 Q_\Omega(\Delta)
 &=
 \langle\bm A,\Delta\rangle\\
 &\le
 \|\bm A\|_{F,r}
 \left\{
 \|\Delta_{T,1}\|_F
 +
 \|\Delta_{T,2}\|_F
 +
 \|\Delta_1\|_F
 +
 \sum_{j\ge2}\|\Delta_j\|_F
 \right\}.
\end{align*}
The first two terms satisfy
\[
 \|\Delta_{T,1}\|_F+\|\Delta_{T,2}\|_F
 \le
 \sqrt2\,\|P_T\Delta\|_F
 \le
 \sqrt2\,d,
\]
while $\|\Delta_1\|_F\le d$.  The standard singular-value block inequality
gives
\[
 \sum_{j\ge2}\|\Delta_j\|_F
 \le
 \frac{
 \|P_{T^\perp}\Delta\|_*
 }{\sqrt r}.
\]
By \eqref{eq:nuclear_cone_slack},
\[
 \frac{
 \|P_{T^\perp}\Delta\|_*
 }{\sqrt r}
 \le
 \sqrt2\,d+
 \frac{\tau_{\lambda,\rm opt}}{\sqrt r}.
\]
Therefore
\begin{align}
\label{eq:nuclear_energy_topr_intermediate}
 Q_\Omega(\Delta)
 \le
 C\|\bm A\|_{F,r}
 \left(
 d+\frac{\tau_{\lambda,\rm opt}}{\sqrt r}
 \right).
\end{align}
Apply \eqref{eq:nuclear_sampled_topr} and
\eqref{eq:nuclear_loss_upper}.

\textit{Error bound.} Use the assumed realized sampling curvature
\begin{align}
\label{eq:nuclear_realized_curvature}
 Q_\Omega(\Delta)
 \ge
 c_{\rm curv}\,p d^2
\end{align}
for the fixed $c_{\rm curv}>0$. We show that
\begin{align}
\label{eq:nuclear_quantitative_rate}
 d
 \le
 C\Bigg\{&
 \frac{\tau_{\lambda,\rm opt}}{\sqrt r}
 +
 \sigma\sqrt{\frac{mr}{p}}
 +
 \sigma
 \sqrt{
 \frac{m\ell_{\rm opt}}
 {p\bar\omega_r}
 }
 \nonumber\\
 &+
 \frac{
 \sigma^2m\lambda\sqrt r
 }{
 p\bar\omega_r
 }
 \Bigg\}.
\end{align}

If
$d\le2\tau_{\lambda,\rm opt}/\sqrt r$, the first term on the right of
\eqref{eq:nuclear_quantitative_rate} proves the claim.  Otherwise,
\[
 d+\frac{\tau_{\lambda,\rm opt}}{\sqrt r}
 \le
 \frac32d.
\]
Combining \eqref{eq:nuclear_Q_upper} with
\eqref{eq:nuclear_realized_curvature}, dividing by $d>0$, and using
$\sqrt{x+y}\le\sqrt x+\sqrt y$ gives
\begin{align*}
 c_{\rm curv}pd
 \le&
 C\sigma\sqrt{rmp}
 +
 C\sigma
 \sqrt{
 \frac{pm\ell_{\rm opt}}
 {\bar\omega_r}
 }\\
 &+
 C\sigma
 \sqrt{
 \frac{pm\lambda\sqrt r}
 {\bar\omega_r}
 }
 d^{1/2}.
\end{align*}
Young's inequality yields
\[
 C\sigma
 \sqrt{
 \frac{pm\lambda\sqrt r}
 {\bar\omega_r}
 }
 d^{1/2}
 \le
 \frac{c_{\rm curv}p}{2}d
 +
 C
 \frac{
 \sigma^2m\lambda\sqrt r
 }{
 \bar\omega_r
 }.
\]
Move the first term to the left and divide by $p$.
\end{proof}

\begin{proof}[Proof of Corollary~\ref{main-cor:convex_tuning_lambda}.]
Let
\[
 \delta_m=\log^{-1}m.
\]
Assume
\begin{align}
\label{eq:nuclear_optimized_tuning}
 \lambda
 =
 c_\lambda\delta_m^2
 \sigma^{-1}\sqrt{p/m},
 \qquad
 \bar\omega_r
 \ge
 c_\omega\delta_m^2,
 \qquad
 \ell_{\rm opt}
 \le
 C_\ell r\delta_m^2.
\end{align}
On $\mathcal E_{\rm pil}$, if
\eqref{eq:nuclear_realized_curvature} holds, then
\begin{align}
\label{eq:nuclear_conditional_rate}
 \|\widehat{\M}_\lambda-\M_0\|_F
 \le
 C\sigma\sqrt{\frac{mr}{p}}.
\end{align}
In particular, because
$\mathcal L_{\rm cl}(\M_0)\le C_0\delta_m^2$, the optimization tolerance may
be as large as
\begin{align}
\label{eq:nuclear_opt_tolerance}
 \varepsilon_{\rm opt}
 \le
 C r\delta_m^2
\end{align}
without changing the statistical order.

Substitute \eqref{eq:nuclear_optimized_tuning} into
Theorem~\ref{main-thm:upper_bound_for_convexlambda}.  Each of the four terms
in \eqref{eq:nuclear_quantitative_rate} is bounded by
$C\sigma\sqrt{mr/p}$.
\end{proof}

Balancing the two $\lambda$-dependent terms in
\eqref{eq:nuclear_quantitative_rate} gives
\[
 \lambda\asymp\sigma^{-1}
 \sqrt{\frac{\ell_{\rm opt}p\bar\omega_r}{mr}}.
\]
The choice \eqref{eq:nuclear_optimized_tuning} applies uniformly over the
stated optimization tolerances and spectral weights.

\begin{scorollary}[Sampling curvature under a spikiness bound]
\label{cor:nuclear_geometric_certificate}
Assume the hypotheses of Corollary~\ref{main-cor:convex_tuning_lambda},
except for
\eqref{eq:nuclear_realized_curvature}, let $mr\ge C\log n$, and suppose
\begin{align}
\label{eq:nuclear_spikiness_localization}
 \|\Delta\|_\infty
 \le
 c_{\rm sp}
 \sqrt{\frac{p}{mr}}\,
 \|\Delta\|_F
\end{align}
for a sufficiently small numerical $c_{\rm sp}$.  Then, on the intersection
of the pilot, sparse-spectral, and Bernoulli sampling events,
\begin{align}
\label{eq:nuclear_geometric_rate}
 \|\widehat{\M}_\lambda-\M_0\|_F
 \le
 C\sigma\sqrt{\frac{mr}{p}}.
\end{align}
\end{scorollary}
\begin{proof}[Proof.]
If
$d\le2\tau_{\lambda,\rm opt}/\sqrt r$, the desired conclusion follows
directly from \eqref{eq:nuclear_optimized_tuning}.  Otherwise
\eqref{eq:nuclear_effective_rank} gives
$\|\Delta\|_*\le C\sqrt r\,d$.  Applying
Lemma~\ref{lem:nuclear_effective_rsc},
\begin{align*}
 C\sqrt{mp}\,
 \|\Delta\|_\infty\|\Delta\|_*
 &\le
 Cc_{\rm sp}\,pd^2,\\
 C\log n\,\|\Delta\|_\infty^2
 &\le
 Cc_{\rm sp}^2
 \frac{p\log n}{mr}d^2.
\end{align*}
For $mr\ge C\log n$ and sufficiently small $c_{\rm sp}$,
\begin{align}
\label{eq:nuclear_curvature_from_spikiness}
 \frac p2d^2
 \le
 Q_\Omega(\Delta)
 \le
 2pd^2.
\end{align}
The lower inequality supplies
\eqref{eq:nuclear_realized_curvature}.  Apply
Corollary~\ref{main-cor:convex_tuning_lambda}.
\end{proof}
\end{revision}
% END MAJOR REVISION S04


% BEGIN MAJOR REVISION S05
\begin{revision}
\section{Perturbation bounds for initialization}
\label{sec:initialization_bounds}
This section gives the uniform sampling and deterministic perturbation
bounds used in the proof of Theorem~\ref{main-thm:track_interation}.

\subsection{Projection, balancing, and spectral scores}
For a factor pair $(\bL,\R)$ and a numerical radius $B>0$, define the scaled
projection $P_B$ by
\begin{align}
\label{eq:scaled_projection_formula}
 \bL_{i,\cdot}
 &=
 \left(
 1\wedge
 \frac{B}{
 \sqrt m\,
 \|\widetilde\bL_{i,\cdot}\widetilde\R'\|_2
 }
 \right)
 \widetilde\bL_{i,\cdot},
 \qquad 1\le i\le m,\\
 \R_{j,\cdot}
 &=
 \left(
 1\wedge
 \frac{B}{
 \sqrt n\,
 \|\widetilde\R_{j,\cdot}\widetilde\bL'\|_2
 }
 \right)
 \widetilde\R_{j,\cdot},
 \qquad 1\le j\le n.
 \nonumber
\end{align}
This is the closed-form scaled incoherence projection of
\cite{tong2021scaled}.  The projection radius is
\begin{align}
\label{eq:rsm_projection_radius}
 B=C_B\sqrt{\mu_0r}\,\sigma_1(\M_0),
\end{align}
which has the same matrix scale as
$\sqrt m\|\M_0\|_{2,\infty}$ and
$\sqrt n\|\M_0'\|_{2,\infty}$.

For every full-rank factor pair $(\bL,\R)$ with
$\operatorname{rank}(\bL\R')=r$, define the balancing map
\begin{align}
\label{eq:balancing_map}
 \operatorname{Bal}(\bL,\R)
 =
 \left(
 \U_{\M}\bm\Sigma_{\M}^{1/2},
 \V_{\M}\bm\Sigma_{\M}^{1/2}
 \right),
 \qquad
 \bL\R'
 =
 \U_{\M}\bm\Sigma_{\M}\V_{\M}'.
\end{align}
The balancing map changes only the factor gauge and leaves the represented
matrix unchanged.

For any matrix $\bm A$, define
\begin{align}
\label{eq:rsm_opclip_definition}
 \operatorname{OpClip}_{L}(\bm A)
 =
 \begin{cases}
 \bm A,&\|\bm A\|\le L,\\[1mm]
 L\bm A/\|\bm A\|,&\|\bm A\|>L.
 \end{cases}
\end{align}

At an iterate $\M=\bL\R'$, let
\begin{align*}
 \bm E(\M)
 &=
 \Y-P_{\bOmega}(\M)
 =
 \U_E\bm D_E\V_E',\\
 \widehat\sigma(\M)
 &=
 \frac{
 \sum_{i=n/3}^{2n/3}D_{E,ii}
 }{
 \sqrt m\sum_{i=n/3}^{2n/3}\widehat\lambda_i
 },\\
 Q(\M)
 &=
 \sqrt m\,\widehat\sigma(\M)\,
 \U_E
 \operatorname{diag}
 (\widehat\lambda_1,\ldots,\widehat\lambda_n)
 \V_E'.
\end{align*}
Define the masked reference contribution
\begin{align}
\label{eq:rsm_masked_reference}
 \bm\Gamma(\M)=p^{-1}P_{\bOmega}\{Q(\M)\}.
\end{align}
For the initialization recurrence, set
\begin{align}
\label{eq:rsm_factor_reference_cap}
 L_{\rm fac}(\M)
 =C_Q\widehat\sigma(\M)\sqrt{m/p}
\end{align}
and define the recovery score
\begin{align}
\label{eq:rsm_stabilized_score}
 G_{\rm RSM}(\M)
 =
 -p^{-1}\bm E(\M)
 +\operatorname{OpClip}_{L_{\rm fac}(\M)}
 \{\bm\Gamma(\M)\}.
\end{align}
Whenever the operator cap is inactive,
\begin{align}
\label{eq:rsm_literal_when_inactive}
 G_{\rm RSM}(\M)
 =p^{-1}P_{\bOmega}
 \{Q(\M)-\bm E(\M)\}.
\end{align}
The operator cap is part of the recovery iteration.  When it is inactive,
the score equals the outer-masked equal-weight RSM direction.

The reference is generated once, independently of the data conditional
on the mask, and is held fixed in the initialization recurrence.

\subsection{Matrix representation of the updates}
\label{sec:recovery_theory}
We use the notation $s_\star=\sigma_r(M_0)$, $s_1=\sigma_1(M_0)$,
$\kappa=s_1/s_\star$, $h=\sigma\sqrt{m/p}$, and $\mu=\mu_0$.
Write $M_0=U_\star\Sigma V_\star^\top$ and let
$G_{\rm safe}=G_{\rm RSM}$ be the equal-weight recovery score
in \eqref{eq:rsm_stabilized_score}.  The reference is generated once
from the observed mask.

For balanced factors $M=LR^\top=U_M S V_M^\top$, use
\begin{align}
 \widetilde L&=L-\eta G_{\rm safe}(M)R(R^\top R+\gamma I_r)^{-1},\label{eq:recovery_updateL}\\
 \widetilde R&=R-\eta G_{\rm safe}(M)^\top L(L^\top L+\gamma I_r)^{-1}.
 \label{eq:recovery_updateR}
\end{align}
The undamped initialization recurrence takes $\gamma=0$.
The convergence results also allow the stated range of damping parameters
$\gamma\ge0$.

Put $\widetilde M=\widetilde L\widetilde R^\top$ and
\[
 (D_L)_{ii}=1\wedge\frac{B}{\sqrt m\norm{\widetilde M_{i,:}}_2},\qquad
 (D_R)_{jj}=1\wedge\frac{B}{\sqrt n\norm{\widetilde M_{:,j}}_2}.
\]
The projected output is $M^+=D_L\widetilde M D_R$, followed by balanced
refactorization.  The convention for a zero row or column is multiplier one.
The theoretical radius is
\begin{equation}
 B=C_B\sqrt{\mu r}\,s_1,\qquad C_B\ge2.\label{eq:recovery_radius}
\end{equation}
Lemma~\ref{lem:computable_initialization} uses a sufficiently large fixed
$C_B=K_0$.  Section~\ref{sec:initialization} supplies a data-based radius
using the observed signal scale and the supplied incoherence envelope.

\subsection{Uniform matrix-level sampling control}
Let $\Delta=M-M_0$, $\delta=\|\Delta\|_{\rm op}$, and
$f=\|\Delta\|_F$.  Put $T_\Omega=p^{-1}P_\Omega-\operatorname{Id}$.
We work on the common event of Lemma~\ref{lem:mask_geometry}, where
\eqref{eq:recovery_tangent-event}--\eqref{eq:recovery_bilinear-event} hold with
$\theta\le C\sqrt{\mu r\log m/(np)}$ and
$\Lambda\le C\sqrt{m/p}$.  The event is uniform over the factors.

\begin{slemma}[Uniform sampling bounds in matrix projectors]
\label{lem:recovery_sampling}
Suppose $\rank(M)=r$, $\delta\le\varepsilon s_\star$ for a sufficiently
small fixed $\varepsilon$, and
$\sqrt m\|M\|_{2,\infty}\vee\sqrt n\|M^\top\|_{2,\infty}\le B$.
On the common event \eqref{eq:recovery_tangent-event}--\eqref{eq:recovery_bilinear-event}, put
\[
 a=C\left\{\theta+\Lambda\frac{C_B^2\kappa^2\mu r}{\sqrt{mn}}\right\}.
\]
If $P_U,P_V$ are the current matrix-space projectors, then
\begin{align}
 \norm{\TT(\Delta)P_V}_\F+\norm{P_U\TT(\Delta)}_\F&\le a f,
 \label{eq:recovery_projected-sampling}\\
 \left|p^{-1}\norm{\PP\Delta}_\F^2-f^2\right|&\le a f^2.
 \label{eq:recovery_energy}
\end{align}
In particular,
$a\le C\{\sqrt{\mu r\log m/(np)}+C_B^2\kappa^2\mu r/\sqrt{np}\}$
on the stated Bernoulli event.
\end{slemma}
\begin{proof}[Proof.]
\textit{Coordinates and row bounds.} There exist full-column-rank $\bar U,\bar V$ such that
$M=\bar U\Sigma\bar V^\top$.  Put
$A=\bar U-U_\star$, $C=\bar V-V_\star$, and
$Z_U=\bar U(\bar U^\top\bar U)^{-1}$,
$Z_V=\bar V(\bar V^\top\bar V)^{-1}$.
Then
\begin{align}
 \norm{A}_\op+\norm{C}_\op&\le C\delta/s_\star,
 &\norm{A\Sigma}_\F+\norm{C\Sigma}_\F&\le C f,\label{eq:recovery_canonical-main}\\
 \norm{Z_U}_\op+\norm{Z_V}_\op&\le C,
 &\norm{Z_U-U_\star}_\F+\norm{Z_V-V_\star}_\F&\le C f/s_\star.
 \label{eq:recovery_canonical-Z}
\end{align}
Writing $b_L=\sqrt{\mu r/m}$ and $b_R=\sqrt{\mu r/n}$, we also have
\begin{align}
 \norm{A\Sigma}_{2,\infty}&\le C C_B s_1 b_L,
 &\norm{C\Sigma}_{2,\infty}&\le C C_B s_1 b_R,\label{eq:recovery_weightedrows}\\
 \norm{A}_{2,\infty}+\norm{Z_U}_{2,\infty}&\le C C_B\kappa b_L,
 &\norm{C}_{2,\infty}+\norm{Z_V}_{2,\infty}&\le C C_B\kappa b_R.
 \label{eq:recovery_unweightedrows}
\end{align}

Let $P_U$ be the projector onto the column space of $M$, and define
$\bar U=P_U U_\star$.  Since
\[
 (I-P_U)U_\star=-(I-P_U)\Delta V_\star\Sigma^{-1},
\]
its norm is at most $\delta/s_\star$, and $\bar U$ has smallest singular
value bounded below by $\sqrt{1-\varepsilon^2}$.  Define
$\bar V^\top=\Sigma^{-1}\bar U^\dagger M$.  Then
$M=\bar U\Sigma\bar V^\top$ and $\bar U^\dagger U_\star=I_r$.
The useful exact identities are
\[
 A\Sigma=(I-P_U)\Delta V_\star,\qquad
 \Sigma C^\top=\bar U^\dagger\Delta.
\]
They prove \eqref{eq:recovery_canonical-main}; they also show that both $\bar U$
and $\bar V$ have singular values within a fixed small neighborhood of one.
Writing $Z_V=(V_\star+C)\{(V_\star+C)^\top(V_\star+C)\}^{-1}$ gives
\eqref{eq:recovery_canonical-Z} by the inverse identity and the corresponding
Frobenius bound for $C$.  The calculation for $Z_U$ is identical with
$A$ in place of $C$.
Finally,
\[
 \bar U\Sigma=M Z_V,\qquad \bar V\Sigma=M^\top Z_U.
\]
The matrix envelope and \eqref{eq:recovery_canonical-Z}, followed by subtraction
of $U_\star\Sigma,V_\star\Sigma$, give \eqref{eq:recovery_weightedrows}.
Multiplication by $\Sigma^{-1}$ and the bounded inverse Grams give
\eqref{eq:recovery_unweightedrows}.

\textit{Sampling error.} In these coordinates, write
\[
 \Delta=A\Sigma V_\star^\top+U_\star\Sigma C^\top+A\Sigma C^\top.
\]
For $Z=Z_V$, expand $\TT(\Delta)Z$ into four terms:
\begin{align*}
 T_1&=\TT(A\Sigma V_\star^\top)V_\star,\\
 T_2&=\TT(A\Sigma V_\star^\top)(Z-V_\star),\\
 T_3&=\TT(U_\star\Sigma C^\top)Z,\\
 T_4&=\TT(A\Sigma C^\top)Z.
\end{align*}
By duality against $X\in\mathbb R^{m\times r}$ with $\norm{X}_\F=1$,
\eqref{eq:recovery_tangent-event} gives $\norm{T_1}_\F\le C\theta f$.
Apply \eqref{eq:recovery_bilinear-event} to the other three terms, respectively
using the following row/Frobenius factors:
\begin{align*}
 \norm{T_2}_\F&\le
 \Lambda\norm{A\Sigma}_{2,\infty}\norm{V_\star}_{2,\infty}
 \norm{Z-V_\star}_\F
 \le C\Lambda C_B\kappa b_Lb_R f,\\
 \norm{T_3}_\F&\le
 \Lambda\norm{U_\star}_{2,\infty}\norm{C\Sigma}_\F\norm{Z}_{2,\infty}
 \le C\Lambda C_B\kappa b_Lb_R f,\\
 \norm{T_4}_\F&\le
 \Lambda\norm{A}_{2,\infty}\norm{C\Sigma}_\F\norm{Z}_{2,\infty}
 \le C\Lambda C_B^2\kappa^2 b_Lb_R f.
\end{align*}
Here the test matrix $X$ contributes only its Frobenius norm, not a
rowwise bound.  Since $P_V=Z_V\bar V^\top$ and
$\norm{\bar V}_\op\le C$, these estimates prove the first right-projected
bound.  Transposing the calculation proves the left-projected bound.

For \eqref{eq:recovery_energy}, set
$\Delta_T=A\Sigma V_\star^\top+U_\star\Sigma C^\top$ and
$\Delta_Q=A\Sigma C^\top$.  Both have Frobenius norm at most $Cf$.
The tangent--tangent bilinear form is at most $C\theta f^2$.
The two mixed forms are bounded using, respectively,
\[
 \Lambda\norm{A\Sigma}_\F\norm{A\Sigma}_{2,\infty}
 \norm{V_\star}_{2,\infty}\norm{C}_\F,
 \qquad
 \Lambda\norm{U_\star}_{2,\infty}\norm{A}_\F
 \norm{C\Sigma}_{2,\infty}\norm{C\Sigma}_\F.
\]
Each is at most $C\Lambda C_B\kappa b_Lb_R f^2$.
For the quadratic--quadratic form, use
\[
 \Lambda\norm{A}_{2,\infty}\norm{A}_\F
 \norm{C\Sigma}_{2,\infty}\norm{C\Sigma}_\F
 \le C\Lambda C_B^2\kappa^2 b_Lb_R f^2.
\]
Sum these terms and use
$\inner{\TT\Delta}{\Delta}=p^{-1}\norm{\PP\Delta}_\F^2-f^2$.
\end{proof}

\subsection{One-step error bounds}
On the event of Lemma~\ref{lem:loss_for_true}, the required noise and scale
bounds are
\begin{align}
 \norm{p^{-1}\PP H}_\op&\le C h,\label{eq:recovery_noiseevent}\\
 \widehat\sigma(M)&\le C\sigma+
 C(mnp)^{-1/2}\norm{\PP(M-\M_0)}_\F
 \quad\hbox{simultaneously for all }M.\label{eq:recovery_scaleevent}
\end{align}
By the definition of the cap and \eqref{eq:recovery_energy}, for $a\le1$,
\begin{align}
 G_{\rm safe}(M)&=p^{-1}\PP\Delta+\Xi(M),\label{eq:recovery_decomp}\\
 \norm{\Xi(M)}_\op&\le C h+C f/\sqrt{np}.\label{eq:recovery_Xi}
\end{align}
Indeed $\Xi=\Clip_{\ell_Q}\Gamma-p^{-1}\PP H$, and
$\widehat\sigma(M)\sqrt{m/p}\le C h+C f/\sqrt{np}$.
These identities are deterministic on the two stated events.

\begin{slemma}[Scaled updates and row--column projection]
\label{lem:rsm_inexact_scaled_update}
Let $M$ have rank $r$, $\delta\le\varepsilon s_\star$, and balanced factors.
Use a matrix score $G=\Delta+W$ in
\eqref{eq:recovery_updateL}--\eqref{eq:recovery_updateR}, with $0<\eta\le2/3$.
Define
\begin{align*}
 w_o&=\max\{\norm{W P_V}_\op,\norm{P_U W}_\op\},\\
 w_f&=\max\{\norm{W P_V}_\F,\norm{P_U W}_\F\}.
\end{align*}
For $s=\sigma_r(M)\ge s_\star-\delta$,
\begin{align}
 \delta_{\rm pre} &\le
 \left\{1-\eta+C\eta\frac{\delta+\gamma}{s_\star}\right\}\delta
 +2\eta w_o+\frac{\eta^2}{s}(2\delta w_o+w_o^2),\label{eq:recovery_op-full}\\
 f_{\rm pre} &\le
 \left\{1-\eta+C\eta\frac{\delta+\gamma}{s_\star}\right\}f
 +2\eta w_f+\frac{\eta^2}{s}(2\delta+w_o)w_f.\label{eq:recovery_f-full}
\end{align}
The constant is numerical when $\varepsilon\le1/4$.

Independently of the update formula, the projection
$M^+=P_B(\widetilde M)$ with $B=C_B\sqrt{\mu r}s_1$ satisfies
\begin{align}
 \norm{M^+-\M_0}_\op&\le
 \norm{\widetilde M-\M_0}_\op+
 \frac{2}{C_B}\norm{\widetilde M-\M_0}_\F,\label{eq:recovery_projectionop}\\
 \norm{M^+-\M_0}_\F&\le
 (1+2/C_B)\norm{\widetilde M-\M_0}_\F.\label{eq:recovery_projectionF}
\end{align}
For rank-at-most-$r$ pre-projection matrices,
$\norm{M^+-\M_0}_\op\le(1+2\sqrt{2r}/C_B)
\norm{\widetilde M-\M_0}_\op$.
\end{slemma}
\begin{proof}[Proof.]
\textit{Scaled update.} With $M=U_M S V_M^\top$, put
\[
 P_U^\gamma=U_M S(S+\gamma I)^{-1}U_M^\top,\quad
 P_V^\gamma=V_M S(S+\gamma I)^{-1}V_M^\top,
 \quad K_\gamma=V_M S(S+\gamma I)^{-2}U_M^\top.
\]
The exact represented-matrix update is
\begin{equation}
 \widetilde M=M-\eta(P_U^\gamma G+G P_V^\gamma)
 +\eta^2 G K_\gamma G.\label{eq:recovery_exact-matrix}
\end{equation}
For $\gamma=0$, $K_0=M^\dagger$.
In all cases $\norm{K_\gamma}_\op\le1/s$, and
$\norm{P_U^\gamma-P_U}_\op\vee\norm{P_V^\gamma-P_V}_\op\le\gamma/s$.

To analyze the population part, decompose $\Delta$ in the true singular
spaces as $\left(\begin{smallmatrix}A_0&B_0\\C_0&D_0\end{smallmatrix}\right)$.
The rank identity yields
\[
 D_0=C_0(\Sigma+A_0)^{-1}B_0,\qquad
 \norm{D_0}_\op\le\frac{\delta^2}{s_\star-\delta},\quad
 \norm{D_0}_\F\le\frac{\delta f}{s_\star-\delta}.
\]
For the reflections $J_U=2P_{U_\star}-I$ and $J_V=2P_{V_\star}-I$,
\begin{align*}
 \Delta-\eta(P_{U_\star}\Delta+\Delta P_{V_\star})
 &=(1-3\eta/2)\Delta-(\eta/2)J_U\Delta J_V\\
 &\quad+2\eta P_{U_\star}^\perp\Delta P_{V_\star}^\perp.
\end{align*}
For either operator or Frobenius norm, the first two terms have norm at most
$(1-\eta)\norm{\Delta}$ when $\eta\le2/3$.
The projector differences between the current and true spaces have norm at
most $\delta/s_\star$.  The remaining normal block and the quadratic term
$\eta^2\Delta K_\gamma\Delta$ are bounded by
$C\eta\delta\norm{\Delta}/s_\star$.
Damping contributes at most $2\eta\gamma\norm{\Delta}/s$.
This proves the two population terms in the claimed bounds.

Finally, subtract the update with $G=\Delta$ from \eqref{eq:recovery_exact-matrix}.
The difference equals
\[
 -\eta(P_U^\gamma W+W P_V^\gamma)
 +\eta^2\{W K_\gamma\Delta+\Delta K_\gamma W+W K_\gamma W\}.
\]
The linear part has operator norm at most $2\eta w_o$ and Frobenius norm
at most $2\eta w_f$.  The quadratic part is controlled by the support of
$K_\gamma$ in the current singular spaces and $\norm{K_\gamma}_\op\le1/s$.
This gives exactly the last terms in \eqref{eq:recovery_op-full}--\eqref{eq:recovery_f-full}.

\textit{Row and column projection.} A clipped row of norm $x>B/\sqrt m$ has multiplier $d=B/(\sqrt m x)$.
Its true row has norm $y\le B/(C_B\sqrt m)$.  Therefore
\[
 (1-d)y\le(y/x)(x-y)\le C_B^{-1}
 \norm{\widetilde M_{i,:}-(\M_0)_{i,:}}_2.
\]
Summing squares gives
$\norm{(I-D_L)\M_0}_\F\le C_B^{-1}
\norm{\widetilde M-\M_0}_\F$.  The same argument applies to $D_R$.
Expand
\[
 D_L\widetilde M D_R-\M_0
 =D_L(\widetilde M-\M_0)D_R
 +(D_L-I)\M_0 D_R+\M_0(D_R-I).
\]
Both diagonal matrices are contractions.  This proves the bounds, with the
rank inequality used only in the final assertion.
\end{proof}
\end{revision}
% END MAJOR REVISION S05


% BEGIN MAJOR REVISION S06
\begin{revision}
\section{Proofs of the algorithmic results}
\label{sec:estimator_algorithms}

We analyze the factor and nuclear-norm algorithms with the prescribed
spectral weights and same-mask reference. The proofs combine the
initialization bounds, the spectral derivative, and the acceptance rules.

Throughout this section the reference $\widehat\lambda_i$ and the weights
$\omega_i$ are held fixed during optimization, with $\omega_i\ge0$ and
$\sum_i\omega_i=1$.  The rank/incoherence class is $\mathbb M_{r,\mu_1}$, as in Theorem~4.1.  The two objective functions are
\begin{align}
 \Phi_{\rm fac}(\M)&=\mathcal L(\M;\Y,\bOmega),
       &&\M\in\mathbb M_{r,\mu_1},\label{eq:of_factor_objective}\\
 \Phi_\lambda(\M)&=\mathcal L_{\rm cl}(\M;\Y,\bOmega)
                  +\lambda\|\M\|_*,
       &&\M\in\mathbb R^{m\times n}.\label{eq:of_nuclear_objective}
\end{align}
The working scale in $\mathcal L_{\rm cl}$ is defined by
\eqref{eq:nuclear_clip_definition}. Candidates with zero bulk scale have
infinite objective.

\subsection{Derivative of the spectral loss}
\label{sec:loss_derivative}
Let
\begin{align*}
 \bm E(\M)&=\Y-P_{\bOmega}(\M)
                =\U_E\operatorname{diag}(d_1,\ldots,d_n)\V_E',\\
 I&=\{\lceil n/3\rceil,\ldots,\lfloor2n/3\rfloor\},\qquad
 S=\sum_{i\in I}d_i,\qquad A_I=\sum_{i\in I}\widehat\lambda_i,\\
 b&=\sqrt m\,\widehat\sigma(\M)=S/A_I,
 \qquad e_i=d_i/b-\widehat\lambda_i,
 \qquad T_e=\sum_i\omega_i e_i d_i.
\end{align*}
Here $A_I>0$ on the reference event.

At a residual with positive, distinct singular values and $S>0$, define
\begin{align}
 a_j=\frac{2}{pb}
 \left\{\omega_j e_j-\frac{\mathbf1_{\{j\in I\}}}{S}T_e\right\}.
 \label{eq:of_spectral_derivative}
\end{align}
Then
\begin{align}
 \nabla\Phi_{\rm fac}(\M)
 =-P_{\bOmega}\{\U_E\operatorname{diag}(a_j)\V_E'\}.
 \label{eq:of_exact_matrix_gradient}
\end{align}
For $\mathcal L_{\rm cl}$, replace $b$ by
$b_c=\min\{b,4\sqrt m\,\widehat\sigma_{\rm pil}\}$ in $e_i$ and the
prefactor.  In the inactive branch retain the correction $T_e/S$; in the
strictly active branch omit it.

Indeed, since $\partial b/\partial d_j=\mathbf1_{\{j\in I\}}/A_I$,
\begin{align*}
 \frac{\partial}{\partial d_j}\frac1p\sum_i\omega_i e_i^2
 &=\frac2p\sum_i\omega_i e_i
   \left\{\frac{\mathbf1_{\{i=j\}}}{b}
     -\frac{d_i}{b^2}\frac{\mathbf1_{\{j\in I\}}}{A_I}\right\},
\end{align*}
which is \eqref{eq:of_spectral_derivative} because $A_Ib=S$.
The differential of a simple positive singular value is
$\mathrm d d_j=\langle u_jv_j',\mathrm d\bm E\rangle$.
Use $\mathrm d\bm E=-P_{\bOmega}(\mathrm d\M)$ and the
self-adjointness of $P_{\bOmega}$ to obtain
\eqref{eq:of_exact_matrix_gradient}.  If the upper cap is strictly active,
$b_c$ is constant and its derivative is zero.


The correction requires one scalar sum after the residual SVD and gives
$\sum_j a_jd_j=0$ in the inactive branch. At a spectral or clipping kink,
the algorithms accept a valid nonsmooth descent proposal or terminate.

\subsection{Initialization and linear convergence}
\label{sec:initialization}
Put
\begin{align*}
 h=\sigma\sqrt{m/p},\qquad \epsilon_*=\sqrt r\,h.
\end{align*}
Let $\M^{(0)}$ be the spectral starting matrix and $k_0$ the stopping
index specified below. The selected iterate $\M^{(k_0)}$ and its scale
$\widehat h$ have the following properties, for fixed constants $C_a,C_q$:
\begin{align}
 \operatorname{rank}(\M^{(k_0)})&\le r,\qquad
 \|\M^{(k_0)}-\M_0\|_F\le C_a\epsilon_*,\label{eq:of_anchor_frobenius}\\
 \|P_{\bOmega}(\M^{(k_0)}-\M_0)\|_F
       &\le C_q\sqrt p\,\epsilon_*,\label{eq:of_anchor_sampled}\\
 \frac12h\le\widehat h&\le\frac32h,
 \qquad \widehat\sigma_{\rm pil}=\widehat h\sqrt{p/m}.
 \label{eq:of_anchor_scale}
\end{align}
The index $k$ is continued without restarting after $k_0$, as in the main
paper. The matrix $\M^{(k_0)}$ is fixed when evaluating the subsequent
acceptance conditions.

\begin{slemma}[Spectral initialization and residual-spectral recovery]
\label{lem:of_anchor_stopping}
\label{lem:computable_initialization}
Assume the homogeneous model, the standing sparse spectral conditions,
$12r\le n$, and $h\le c_hs_\star$.  Let $K_0$ be sufficiently large and $t_0$ sufficiently small, and suppose
\begin{align}
 \mu_1&\ge C_\mu\mu_0\kappa^2,\qquad
 \bar\mu=\max\{\mu_0,\mu_1\},\label{eq:init_candidate_envelope}\\
 np&\ge C\left[\mu_0\kappa^2 r^2\log m+
                 \mu_0^2\kappa^4r^2+
        \kappa^2\{\bar\mu r\log m+\bar\mu^2r^2\}\right].
 \label{eq:init_default_sampling}
\end{align}
Write $Z=p^{-1}\Y$, $Z_r=\operatorname{Proj}_r(Z)$,
$\widehat s_i=s_i(Z)$ and
$\widehat h_{\rm raw}=\widehat\sigma(0)\sqrt{m/p}$. Set
$B=K_0\sqrt{\mu_0r}\,\widehat s_1$.  Form
$M_b=P_B(Z_r)$ with the simultaneous row/column projection in
Section~\ref{sec:recovery_theory}.
\begin{enumerate}[(i)]
\item\enspace If $\sqrt r\widehat h_{\rm raw}\ge t_0\widehat s_r$, set $k_0=0$ and
\begin{align}
 \M^{(0)}=\operatorname{HT}_{\widehat s_r/4}(M_b),\qquad
 \operatorname{HT}_t(A)=\sum_{s_i(A)\ge t}s_i(A)u_iv_i'.
 \label{eq:init_large_anchor}
\end{align}
\item\enspace Otherwise, set $\M^{(0)}=M_b$ and use the fixed-envelope
updates of Section~\ref{sec:recovery_theory}, with $\eta=1/2$ and $\gamma=0$.
Set $k_0=k+1$ at the first occurrence of
\begin{align}
 \|\M^{(k+1)}-\M^{(k)}\|_F
 \le C_{\rm stop}\sqrt r\,\widehat\sigma(\M^{(k)})\sqrt{m/p}.
 \label{eq:of_anchor_stop}
\end{align}
\end{enumerate}
In either branch set
$\widehat h=\widehat\sigma(\M^{(k_0)})\sqrt{m/p}$.
With probability at least $1-cn^{-3}$, the construction is well defined,
terminates, gives $\M^{(k_0)}\in\mathbb M_{r,\mu_1}$, and satisfies
\eqref{eq:of_anchor_frobenius}--\eqref{eq:of_anchor_scale}.
In case (ii), let $f_k=\|\M^{(k)}-\M_0\|_F$. There are fixed
$0<q\le3/4$ and $C_*>0$ such that
\begin{align}
 f_{k+1}&\le qf_k+C_*\sqrt r\,h,\label{eq:init_frob_recurrence}\\
 f_k&\le q^kf_0+\frac{C_*(1-q^k)}{1-q}\sqrt r\,h,
       \label{eq:recovery_explicit_geometric}
\end{align}
up to the stopping iterate. At each executed step for which
$f_k\ge2C_*\sqrt r h/(1-q)$,
\begin{align}
 f_{k+1}\le\frac{1+q}{2}f_k.
 \label{eq:recovery_above_floor_contraction}
\end{align}
The stopping test is reached within
$C\{1+\log_+(f_0/(\sqrt r h))\}$ updates, where
$\log_+(x)=\log\max\{1,x\}$.
\end{slemma}
\begin{proof}[Proof.]
\textit{Step 1: local contraction and stopping.} Consider the simultaneous updates in Section~\ref{sec:recovery_theory}
with fixed $C_B=K_0$ and $\eta=1/2$. For a local start satisfying
\begin{align}
 a_0+\sqrt{r/(np)}&\le c_s,\qquad
 a_0=C\left\{\sqrt{\frac{\mu_0r\log m}{np}}+
                   \frac{K_0^2\kappa^2\mu_0r}{\sqrt{np}}\right\},
 \label{eq:init_frob_sampling}\\
 f_0=\|\M^{(0)}-\M_0\|_F&\le c_0s_\star,
 \qquad \sqrt r\,h\le c_fs_\star,
 \label{eq:init_small_radius}
\end{align}
the radius $B$ may lie between fixed sufficiently large multiples
of $\sqrt{\mu_0r}s_1$.  
We first verify the local recurrence and stopping properties under these conditions.
The argument also permits
$0\le\gamma_k\le c_d\widehat\sigma(\M^{(k)})\sqrt{m/p}$ for fixed $c_d$.

Within $f\le\varepsilon s_\star$, one has $\delta\le f$.
Lemma~\ref{lem:recovery_sampling} with fixed $C_B$ gives sampled energy
comparable to $pf^2$.  The capped reference and \eqref{eq:working_scale_bound} give
$\|\Xi\|\le Ch+Cf/\sqrt{np}$.
For $W=T_\Omega\Delta+\Xi$, the quantities in
Lemma~\ref{lem:rsm_inexact_scaled_update} satisfy
\[
 w_o\le a_0f+Ch+Cf/\sqrt{np},\qquad
 w_f\le a_0f+C\sqrt r h+C\sqrt{r/(np)}f.
\]
Also $\gamma\le Cc_d(h+f/\sqrt{np})$.
Apply the Frobenius bound \eqref{eq:recovery_f-full}, including its
quadratic term.  The coefficient $(\delta+w_o)/s_\star$ is a small
constant under \eqref{eq:init_small_radius}; hence that term is absorbed
into $C\eta\varepsilon w_f$.  The pre-projection recurrence is
\[
 f_{\rm pre}\le[1-\eta+C\eta\{\varepsilon+a_0+
                 \sqrt{r/(np)}+c_dh/s_\star\}]f+C\eta\sqrt r h.
\]
By \eqref{eq:recovery_projectionF}, the Frobenius projection multiplier is $1+2/K_0$.  Choose $K_0$ large enough that the Frobenius projection multiplier is at
most $9/8$, allowing for the fixed comparison between observed and true
radii. The small constants can then make the pre-projection coefficient
at most $2/3$, so the post-projection coefficient is at most $3/4$.
This proves \eqref{eq:init_frob_recurrence} with a fixed factor and preserves
$f\le\varepsilon s_\star$ by choosing $c_0,c_f$ small enough.
Rank remains $r$ by Weyl's inequality. Summing the geometric series gives
\eqref{eq:recovery_explicit_geometric}. If
$f_k\ge2C_*\sqrt r h/(1-q)$, the additive term is at most
$(1-q)f_k/2$, proving \eqref{eq:recovery_above_floor_contraction}.


Let $\widehat h_k=\widehat\sigma(\M^{(k)})\sqrt{m/p}$.
The scale bound gives $\widehat h_k\le Ch+C f_k/\sqrt{np}$.
At a stopping step the triangle inequality and the recurrence imply
\[
 (1-q)f_k\le\|\M^{(k+1)}-\M^{(k)}\|_F+C\sqrt r h
 \le C\sqrt r h+C C_{\rm stop}\sqrt{r/(np)}f_k.
\]
Increasing the sampling constant absorbs the last term.
The recurrence then bounds $f_{k+1}$ as well, and the uniform sampled
energy bound supplies both sampled errors.  Conversely, writing $C_*$ for the additive constant in
\eqref{eq:init_frob_recurrence}, iteration gives
\[
 f_k\le q^k f_0+\frac{C_*}{1-q}\sqrt r\,h.
\]
For any fixed $C_{\rm tar}>C_* /(1-q)$, a sufficient bound is
\begin{align}
 k\ge
 \left\lceil
 \frac{\log_+\{f_0/[(C_{\rm tar}-C_* /(1-q))\sqrt r\,h]\}}
      {\log(1/q)}
 \right\rceil.
 \label{eq:initialization_iteration_complexity}
\end{align}
An additional step controls $f_k+f_{k+1}$ at the same order.
At that point the two-sided working-scale bound gives
$|\widehat\sigma(\M^{(k)})/\sigma-1|
\le C\{\log^{-1}m+\sqrt{r/(np)}\}<1/2$.
Choosing $C_{\rm stop}$ large enough ensures the stopping test.

\textit{Step 2: observed bulk and signal scales.} Let $Z=p^{-1}\Y$, $Z_r=\operatorname{Proj}_r(Z)$,
$\widehat s_1=s_1(Z)$, $\widehat s_r=s_r(Z)$, and
\begin{align}
 \widehat h_{\rm raw}
 =\frac{\sum_{i\in I}s_i(Z)}{\sum_{i\in I}\widehat\lambda_i/\sqrt p}
 =\widehat\sigma(0)\sqrt{m/p}.
 \label{eq:init_raw_bulk_scale}
\end{align}
Assume $12r\le n$.  On an event of probability at least $1-cn^{-4}$,
\begin{align}
 b_{\rm sig}:=\|(p^{-1}P_{\bOmega}-\operatorname{Id})\M_0\|
 &\le Cs_1\left\{\sqrt{\frac{\mu_0r\log m}{np}}+
            \frac{\mu_0r\log m}{p\sqrt{mn}}\right\},
 \label{eq:init_signal_concentration}\\
 |\widehat h_{\rm raw}-h|
 &\le C b_{\rm sig}+C\{\log^{-1}m+r/n\}h.
 \label{eq:init_raw_scale_error}
\end{align}
Fix a sufficiently small numerical $t_0>0$.  If
$b_{\rm sig}\le c_b s_\star/\sqrt r$ and $h\le c_hs_\star$, with
$c_b$ small relative to $t_0$ and $c_h$ sufficiently small, then
$\widehat s_1\asymp s_1$, $\widehat s_r\asymp s_\star$, and, for
all sufficiently large $n$,
\begin{align}
 \sqrt r\widehat h_{\rm raw}\ge t_0\widehat s_r
 &\quad\Longrightarrow\quad \sqrt r h\ge c t_0s_\star,
 \label{eq:init_large_branch}\\
 \sqrt r\widehat h_{\rm raw}<t_0\widehat s_r
 &\quad\Longrightarrow\quad \sqrt r h\le C t_0s_\star.
 \label{eq:init_small_branch}
\end{align}
A sufficient sampling condition for the stated signal bound is
$np\ge C\mu_0\kappa^2r^2\log m$.

For the fixed-signal summands, matrix Bernstein has variance bounded by
$Cs_1^2\mu_0r/(np)$ and summand norm bounded by
$Cs_1\mu_0r/(p\sqrt{mn})$.  This proves
\eqref{eq:init_signal_concentration}.  Write
$Z=\M_0+N+S_0$, where $N=p^{-1}P_{\bOmega}\bH$ and
$\|S_0\|=b_{\rm sig}$.
For $i\in I$, rank perturbation gives
$s_{i+r}(N)\le s_i(\M_0+N)\le s_{i-r}(N)$.
Summing over the consecutive middle-third indices cancels all but at most
$2r$ endpoint terms, so
\[
 \left|\sum_{i\in I}s_i(Z)-\sum_{i\in I}s_i(N)\right|
 \le |I|b_{\rm sig}+2r\|N\|.
\]
On the homogeneous spectral event $\|N\|\le Ch$ and
$|s_i(N)-h\widehat\lambda_i/\sqrt p|\le Ch\log^{-1}m$.
Divide by the reference middle sum, which is at least $cn$, to obtain
\eqref{eq:init_raw_scale_error}.  

Weyl's inequality and $b_{\rm sig}+Ch\le s_\star/10$ make both
observed signal scales lie between, for example, $0.9$ and $1.1$ times
their true values.  For the branch comparisons it is enough to use two-sided bulk bounds,
rather than to assume $r/n\to0$ in the preceding additive-error bound.
Because $12r\le n$, the shifted middle indices $i-r,i+r$ belong to
$[n/4,3n/4]$ up to integer rounding.  Coarse rigidity bounds the noise
singular values at these indices between fixed positive multiples of $h$.
Rank interlacing and Weyl's inequality therefore give
\[
 c h-Cb_{\rm sig}\le\widehat h_{\rm raw}
                        \le C h+Cb_{\rm sig}.
\]
Choose $c_b$ small relative to $t_0$ and these fixed bulk constants.
Substitution into the two observable comparisons, using
$\widehat s_r\asymp s_\star$, proves
\eqref{eq:init_large_branch}--\eqref{eq:init_small_branch}.
The displayed sample size makes the first signal term at most
$c_bs_\star/\sqrt r$; the second is smaller after increasing its constant.

\textit{Step 3: the computable initialization.} Choose the projection and contraction constants first, then $t_0$ small
enough for \eqref{eq:init_small_radius}, followed by $c_b,c_h$ and the
sampling constant.  The scale comparison and
$\|Z_r-\M_0\|\le2\|Z-\M_0\|$ give
\[
 \|Z_r-\M_0\|_F\le C\sqrt r(b_{\rm sig}+h).
\]
Lemma~\ref{lem:rsm_inexact_scaled_update} preserves this Frobenius bound with a
fixed multiplier.  Since $\widehat s_1\asymp s_1$, the observed $B$
is within fixed factors of the radii covered by the projection bound.
Consequently the projected spectral starting error is bounded by
$C\sqrt r(b_{\rm sig}+h)$, which is
\eqref{main-eq:main_spectral_start_accuracy}.

In branch (ii), \eqref{eq:init_small_branch} makes the initial Frobenius
error and $\sqrt r h$ sufficiently small multiples of $s_\star$.
Condition \eqref{eq:init_default_sampling} implies
\eqref{eq:init_frob_sampling}.  Step 1 gives
finite stopping and the two anchor error bounds.  Its final singular
values are bounded below by $s_\star/2$.  The projected row envelope
therefore implies incoherence at most $C\mu_0\kappa^2$.

In branch (i), \eqref{eq:init_large_branch} and
$b_{\rm sig}\le c_bs_\star/\sqrt r$ imply
$\|M_b-\M_0\|_F\le C\sqrt r h$.
Let $\tau=\widehat s_r/4$, so $s_\star/8\le\tau\le3s_\star/8$.
For a discarded singular value $d_i<\tau$, with $i\le r$,
\[
 d_i\le\frac{\tau}{s_\star-\tau}
                   |s_i(\M_0)-d_i|.
\]
Mirsky's inequality therefore gives
$\|\M^{(k_0)}-M_b\|_F\le
[\tau/(s_\star-\tau)]\|M_b-\M_0\|_F$.
Thus \eqref{eq:of_anchor_frobenius} holds also in this branch without
requiring $\sqrt r h<s_\star$.
Hard thresholding does not increase row or column Euclidean norms.
Every retained singular value is at least $\tau$, whence
\[
 \|U_{\M^{(k_0)}}\|_{2,\infty}^2
 \le\frac{B^2}{m\tau^2}\le\frac{C\mu_0\kappa^2 r}{m},
 \qquad
 \|V_{\M^{(k_0)}}\|_{2,\infty}^2
 \le\frac{C\mu_0\kappa^2 r}{n}.
\]
The assertion is trivial for rank zero.  Moreover
$\|\M^{(k_0)}\|\le\|M_b\|\le\|Z_r\|\le2s_1$.
Lemma~\ref{lem:global_relative_sampling}, which permits rank-deficient candidates
and does not require an operator-local error, gives
\[
 \|P_{\bOmega}(\M^{(k_0)}-\M_0)\|_F^2
 \le\tfrac32p\|\M^{(k_0)}-\M_0\|_F^2\le Cprh^2.
\]
This supplies the sampled error bound required by the nuclear acceptance test.

In both branches, the last sampled bound and
\eqref{eq:working_scale_bound} yield
\[
 |\widehat\sigma(\M^{(k_0)})/\sigma-1|
 \le C\{\log^{-1}m+\sqrt{r/(np)}\}\le1/2.
\]
This proves the scale assertion, positive bulk scale, and all membership
claims.  A union bound over the fixed uniform events finishes the proof;
no union bound over a random number of iterates is needed.
\end{proof}


The hard-thresholding operation can be computed by thin QR
factorizations and an $r\times r$ SVD. Recompute the pilot from the
returned matrix in both cases. A zero reference bulk, zero observed
$r$th singular value, or a resource limit reached before the stopping
test produces an initialization-failure status.

Fix a numerical $A_0>0$ and define the observable ball
\begin{align}
 R_a=A_0\sqrt r\,\widehat h,\qquad
 \mathcal B_a=\{\M:\|\M-\M^{(k_0)}\|_F\le R_a\}.
 \label{eq:of_statistical_ball}
\end{align}


\subsection{Statistical acceptance conditions}
\label{sec:acceptance_conditions}
For the factor branch define
\begin{align}
 \mathcal A_{\rm fac}
 =\left\{\M\in\mathbb M_{r,\mu_1}:
       \M\in\mathcal B_a\ \text{or}\ \mathcal L(\M)\le\omega_1\right\}.
 \label{eq:of_factor_acceptance}
\end{align}
Membership is checked from the current matrix, its thin singular spaces,
and its computed loss.  Once a descent sequence has loss at most
$\omega_1$, the ball condition can be dropped permanently.
In particular, before computing an anchor one may test a rank-constrained initial matrix.
If it belongs to $\mathbb M_{r,\mu_1}$ and has loss at most $\omega_1$, initialize
objective descent there and omit the anchor construction altogether.
The same test applies to an initial estimate from another solver.

For the nuclear branch let $\delta_m=\log^{-1}m$ and choose
\begin{align}
 \lambda=c_\lambda\delta_m^2/\widehat h,
 \qquad \bar\omega_r=\min_{i\le r}\omega_i\ge c_\omega\delta_m^2.
 \label{eq:of_nuclear_tuning}
\end{align}
Fix $a_0\in(0,1)$ and $C_J>0$.  Define the computable conditions
\begin{align}
 \Phi_\lambda(\M)&\le
       \lambda\|\M^{(k_0)}\|_*+C_Jr\delta_m^2,
       \label{eq:of_nuclear_level_test}\\
 \|P_{\bOmega}(\M-\M^{(k_0)})\|_F
       &\ge a_0\sqrt p\,\|\M-\M^{(k_0)}\|_F.
       \label{eq:of_nuclear_sampling_test}
\end{align}
Let $\mathcal N_a$ be the set satisfying both tests, and set
\begin{align}
 \mathcal A_{\rm nuc}=\mathcal B_a\cup\mathcal N_a.
 \label{eq:of_nuclear_acceptance}
\end{align}
The sampling test uses the computed $\M^{(k_0)}$ and permits
full-rank nuclear candidates.

\subsection{Descent and proximal updates}
\label{sec:descent_updates}
At a differentiable factor iterate $\M=\bL\R'$, let
$G=\nabla\mathcal L(\M)$ from
\eqref{eq:of_exact_matrix_gradient}.  Choose $\gamma>0$, for example a
fixed positive multiple of $\widehat h$, and set
\begin{align}
 D_L&=-G\R(\R'\R+\gamma I)^{-1},&
 D_R&=-G'\bL(\bL'\bL+\gamma I)^{-1}.
 \label{eq:of_factor_direction}
\end{align}
Define
\begin{align}
 d_{\mathrm{dir}}={}&
 \|G\R(\R'\R+\gamma I)^{-1/2}\|_F^2
 +\|G'\bL(\bL'\bL+\gamma I)^{-1/2}\|_F^2.
 \label{eq:of_factor_decrement}
\end{align}
Then the exact derivative of
$t\mapsto\mathcal L\{(\bL+tD_L)(\R+tD_R)'\}$ at zero is
$-d_{\mathrm{dir}}$.  With $0<c_A<1$, backtracking accepts only candidates
$Z(t)=(\bL+tD_L)(\R+tD_R)'$ satisfying
\begin{align}
 Z(t)&\in\mathcal A_{\rm fac},&
 \mathcal L\{Z(t)\}&\le\mathcal L(\M)-c_A td_{\mathrm{dir}}.
 \label{eq:of_factor_armijo}
\end{align}
Thin QR and a small SVD rebalance accepted factors. At a rank-deficient
initialization, factors may be padded with zero columns or zero-product
column pairs. Positive damping keeps the linear systems invertible.
Each proposal is checked for membership in $\mathbb M_{r,\mu_1}$.

For $A=U_A\operatorname{diag}(a_i)V_A'$ define
\begin{align}
 \operatorname{SVT}_t(A)
 =U_A\operatorname{diag}\{(a_i-t)_+\}V_A'
 =\operatorname*{argmin}_X
 \left\{\tfrac12\|X-A\|_F^2+t\|X\|_*\right\}.
 \label{eq:svt_definition}
\end{align}
For the nuclear branch put $g=\nabla\mathcal L_{\rm cl}(\M)$ and form
\begin{align}
 Z_0(t)=\operatorname{SVT}_{t\lambda}(\M-tg).
 \label{eq:of_plain_nuclear_prox}
\end{align}
The nuclear proximity identity is due to \cite{cai2010singular}.

If $\M\in\mathcal B_a$ but $Z_0(t)\notin\mathcal A_{\rm nuc}$, replace
the proposal by the ball-constrained proximal step
\begin{align}
 Z_a(t)=\operatorname*{argmin}_{Z\in\mathcal B_a}
 \left\{\frac1{2t}\|Z-(\M-tg)\|_F^2+\lambda\|Z\|_*\right\}.
 \label{eq:of_ball_nuclear_prox}
\end{align}
This is a convex subproblem even though the outer RSM objective is not
convex. If the current matrix is outside the ball, accept a plain
proposal passing the certificate or return the last accepted matrix.
For either type of proposal, check the proximal-model comparison
\begin{align}
 \langle g,Z-\M\rangle+\lambda(\|Z\|_*-\|\M\|_*)
       +\frac{1}{2t}\|Z-\M\|_F^2\le0.
 \label{eq:of_nuclear_model_test}
\end{align}
It holds automatically for the exact subproblems when the current matrix
is feasible. For an inexact proximal solve, impose the same model test.
The loss test is
\begin{align}
 \mathcal L_{\rm cl}(Z)
 \le\mathcal L_{\rm cl}(\M)+\langle g,Z-\M\rangle
       +\frac{1-c_A}{2t}\|Z-\M\|_F^2.
 \label{eq:of_nuclear_majorization}
\end{align}

Let $A=\M-tg$.  If $\operatorname{SVT}_{t\lambda}(A)$ is in
$\mathcal B_a$, it solves \eqref{eq:of_ball_nuclear_prox}.  Otherwise the
unique solution is
\begin{align}
 Z_\theta=
 \operatorname{SVT}_{t\lambda/(1+\theta)}
       \left(\frac{A+\theta\M^{(k_0)}}{1+\theta}\right),
 \qquad \theta>0,
 \qquad\|Z_\theta-\M^{(k_0)}\|_F=R_a.
 \label{eq:of_ball_prox_scalar}
\end{align}
The distance in this equation is continuous and nonincreasing in $\theta$
and tends to zero.  Bracketing and scalar bisection therefore determine a
solution.

To verify the formula, note that the subproblem is strongly convex and has a strictly feasible center.
For the multiplier $\theta/(2t)$ of
$\|Z-\M^{(k_0)}\|_F^2-R_a^2$, completing the square gives
\eqref{eq:of_ball_prox_scalar}.  For $\theta_2>\theta_1$, add the two
optimality inequalities for their penalized objectives to get
\begin{align*}
 (\theta_2-\theta_1)
 \{\|Z_{\theta_2}-\M^{(k_0)}\|_F^2
       -\|Z_{\theta_1}-\M^{(k_0)}\|_F^2\}\le0.
\end{align*}
Continuity follows from the displayed continuous SVT formula.  Comparing
the penalized objective with its value at $\M^{(k_0)}$ gives convergence
of the distance to zero.  The intermediate value theorem supplies a
multiplier, and its optimality and complementary slackness prove the claim.
For approximate bisection the final iterate must be chosen on the feasible
side and must also pass \eqref{eq:of_nuclear_model_test}.  If the model test
fails, refine the scalar search or reject the proposal.


\begin{algorithm}[!ht]\revisionfloatcolor{}
\caption{Nuclear-norm regularized residual spectral matching}
\label{alg:of_nuclear}
\begin{algorithmic}[1]
\State \textbf{Input:} $Y,\Omega,p$, weights $\omega_i$, reference draw count $T$,
regularization $\lambda>0$, a working rank/envelope for initialization, and
finite iteration and backtracking budgets.
\State Generate the same-mask reference once. Construct
$\M^{(0)},\ldots,\M^{(k_0)}$ by Section~\ref{sec:initialization}; set
$\widehat\sigma_{\rm pil}=\widehat h\sqrt{p/m}$ and
$R_a=A_0\sqrt r\,\widehat h$.
\State Construct $\mathcal A_{\rm nuc}$ using
\eqref{eq:of_nuclear_level_test}--\eqref{eq:of_nuclear_acceptance}.
\For{$k=k_0,\ldots,k_0+K-1$}
\State Compute $g_k=\nabla\mathcal L_{\rm cl}(\M^{(k)})$ using \eqref{eq:of_exact_matrix_gradient}; stop at a nonsmooth point if no valid descent step is available.
\State For trial $t>0$, set
$Z=\operatorname{SVT}_{t\lambda}(\M^{(k)}-tg_k)$.
\State If $\M^{(k)}\in\mathcal B_a$ but $Z\notin\mathcal A_{\rm nuc}$,
replace $Z$ by \eqref{eq:of_ball_nuclear_prox}, evaluated by \eqref{eq:of_ball_prox_scalar}.
\State Reduce $t$ until $Z\in\mathcal A_{\rm nuc}$ and both
\eqref{eq:of_nuclear_model_test} and \eqref{eq:of_nuclear_majorization} hold.
\State If no trial is accepted, stop and retain $\M^{(k)}$;
otherwise set $\M^{(k+1)}=Z$, without a rank truncation.
\State Stop when the prescribed tolerance is met.
\EndFor
\State \textbf{Output:} the last accepted matrix, its objective and certificate,
and the termination status.  Unsuccessful initialization is reported separately.
\end{algorithmic}
\end{algorithm}

\begin{proof}[Proof of Theorem~\ref{main-thm:track_interation} and Corollary~\ref{main-cor:iteration_results}.]
Work on the common event of Lemma~\ref{lem:computable_initialization},
Theorem~\ref{main-thm:upper_bound_for_rankr}, and the sparse spectral calibration.
Use the specified acceptance rules and nuclear tuning.

\textit{Statistical control of accepted candidates.} By Lemma~\ref{lem:computable_initialization}, \eqref{eq:of_anchor_frobenius}--\eqref{eq:of_anchor_scale} hold.
We first establish the uniform bounds
\begin{align}
 \sup_{\M\in\mathcal A_{\rm fac}}\|\M-\M_0\|_F
       &\le C\sigma\sqrt{mr/p},\label{eq:of_factor_rate}\\
 \sup_{\M\in\mathcal A_{\rm nuc}}\|\M-\M_0\|_F
       &\le C\sigma\sqrt{mr/p}.\label{eq:of_nuclear_rate}
\end{align}
The suprema are over finite-loss candidates.

Any $\M\in\mathcal B_a$ satisfies
\begin{align*}
 \|\M-\M_0\|_F\le R_a+C_a\epsilon_*
       \le(3A_0/2+C_a)\epsilon_*.
\end{align*}
For a factor candidate outside the ball, its low-loss test and
Theorem~\ref{main-thm:upper_bound_for_rankr} prove the result.

Consider a nuclear candidate in $\mathcal N_a$.  Since
$\operatorname{rank}(\M^{(k_0)}-\M_0)\le2r$,
\begin{align*}
 \big|\|\M^{(k_0)}\|_*-\|\M_0\|_*\big|
 &\le\|\M^{(k_0)}-\M_0\|_*\le C_a\sqrt{2r}\,\epsilon_*.
\end{align*}
Consequently \eqref{eq:of_nuclear_level_test} implies
\begin{align*}
 \Phi_\lambda(\M)
 \le\lambda\|\M_0\|_*+C r\delta_m^2
 \le\Phi_\lambda(\M_0)+C r\delta_m^2.
\end{align*}
Thus the optimization slack in
\eqref{eq:nuclear_approx_objective} can be chosen of order
$r\delta_m^2$, and $\ell_{\rm opt}\le C r\delta_m^2$.
The pilot-scale estimate also makes \eqref{eq:of_nuclear_tuning}
comparable to \eqref{eq:nuclear_optimized_tuning}.

Let $d=\|\M-\M_0\|_F$ and $x=\|\M-\M^{(k_0)}\|_F$.
If $d\le2(C_a+C_q/a_0)\epsilon_*$, the desired bound already holds.
Otherwise, by \eqref{eq:of_nuclear_sampling_test},
\begin{align*}
 \|P_{\bOmega}(\M-\M_0)\|_F
 &\ge\|P_{\bOmega}(\M-\M^{(k_0)})\|_F
       -\|P_{\bOmega}(\M^{(k_0)}-\M_0)\|_F\\
 &\ge\sqrt p\{a_0(d-C_a\epsilon_*)-C_q\epsilon_*\}
 \ge\tfrac12a_0\sqrt p\,d.
\end{align*}
This supplies the realized curvature with
$c_{\rm curv}=a_0^2/4$.  Apply
Theorem~\ref{main-thm:upper_bound_for_convexlambda}.  Its four terms are
bounded respectively by
\begin{align*}
 \frac{Cr\delta_m^2}{\lambda\sqrt r},\quad
 \sqrt r\,h,\quad
 h\sqrt{Cr\delta_m^2/\bar\omega_r},\quad
 h^2\lambda\sqrt r/\bar\omega_r,
\end{align*}
each of which is $O(\sqrt r\,h)$.  This proves
\eqref{eq:of_nuclear_rate}.  All arguments hold on common events and are
therefore valid for candidates selected using the same observations.

\textit{Objective descent.} The initialization belongs to $\mathcal B_a$,
and each accepted output belongs to its prescribed acceptance set.
Equations~\eqref{eq:of_factor_rate}--\eqref{eq:of_nuclear_rate} therefore give
\begin{align}
 \|\M^{(k_0+K)}-\M_0\|_F\le C\sigma\sqrt{mr/p}.
 \label{eq:of_computable_optimal_rate}
\end{align}
A directly certified factor start stays below $\omega_1$ by the Armijo test,
so the same assertion follows from Theorem~\ref{main-thm:upper_bound_for_rankr}.
The factor loss is nonincreasing for $k\ge k_0$ by \eqref{eq:of_factor_armijo}. For every
global factor estimator, the triangle inequality also yields
\begin{align}
 \|\M^{(k_0+K)}-\widetilde\M\|_F\le C\sigma\sqrt{mr/p},\qquad
 \widetilde\M\in\operatorname*{argmin}_{\M\in\mathbb M_{r,\mu_1}}\mathcal L(\M).
 \label{eq:of_estimator_distance}
\end{align}
For a nuclear step, comparison of the exact proximal subproblem with the
current feasible matrix gives
\[
 \langle g,Z-\M\rangle+\lambda(\|Z\|_*-\|\M\|_*)
 \le-\frac1{2t}\|Z-\M\|_F^2.
\]
The ball subproblem is used only when the current matrix is in the ball;
for an inexact solve, the same inequality is enforced by
\eqref{eq:of_nuclear_model_test}. Adding \eqref{eq:of_nuclear_majorization}
proves
\begin{align}
 \Phi_\lambda(\M)-\Phi_\lambda(Z)
 \ge\frac{c_A}{2t}\|Z-\M\|_F^2.
 \label{eq:of_nuclear_sufficient_decrease}
\end{align}
Rejected proposals leave the previous accepted matrix unchanged, so all
output bounds remain valid at termination.

\textit{The complete iterate sequence.} The matrices $\M^{(k)}$ are indexed
consecutively from $\M^{(0)}$. The direct spectral choice has $k_0=0$.
Rejected trials are not counted, and the terminal matrix is repeated after
termination. Lemma~\ref{lem:computable_initialization} gives, for $k<k_0$,
\[
 \|\M^{(k+1)}-M_0\|_F
 \le q\|\M^{(k)}-M_0\|_F+C_{\rm rec}\epsilon_*.
\]
For $k\ge k_0$, \eqref{eq:of_computable_optimal_rate} bounds each accepted
or repeated terminal matrix by $C_{\rm out}\epsilon_*$. With
$C_*\ge\max\{C_{\rm rec},C_{\rm out}\}$, nonnegativity of
$\|\M^{(k)}-M_0\|_F$ gives, for every $k\ge0$,
\begin{align}
 \|\M^{(k+1)}-M_0\|_F
 &\le q\|\M^{(k)}-M_0\|_F+C_*\epsilon_*,\label{eq:of_complete_recurrence}\\
 \|\M^{(k)}-M_0\|_F
 &\le q^k\|\M^{(0)}-M_0\|_F+
       \frac{C_*(1-q^k)}{1-q}\epsilon_*.
 \label{eq:of_complete_geometric}
\end{align}
The second inequality follows by summing the geometric series. The stopping
bound in Lemma~\ref{lem:computable_initialization} gives an update count of
$C\{1+\log_+(\|\M^{(0)}-M_0\|_F/\epsilon_*)\}$. Equivalently, solve
$q^k\|\M^{(0)}-M_0\|_F\le\epsilon_*$; termination after successful
initialization already returns a statistically accurate matrix.
These are \eqref{main-eq:main_uniform_contraction}--\eqref{main-eq:main_geometric_recovery}
and \eqref{main-eq:main_recovery_complexity}. When
$\|\M^{(k)}-M_0\|_F\ge2C_*\epsilon_* /(1-q)$, the additive term in
\eqref{eq:of_complete_recurrence} is at most
$(1-q)\|\M^{(k)}-M_0\|_F/2$, proving \eqref{main-eq:main_linear_above_floor}.

\end{proof}

\begin{scorollary}[Preservation of the factor estimator]
\label{cor:of_estimator_fidelity}
On the common factor event used above, for
all sufficiently large $m$,
\begin{align}
 \inf_{\M\in\mathcal A_{\rm fac}}\mathcal L(\M)
 &=\inf_{\M\in\mathbb M_{r,\mu_1}}\mathcal L(\M),\label{eq:of_same_infimum}\\
 \operatorname*{argmin}_{\M\in\mathcal A_{\rm fac}}\mathcal L(\M)
 &=\operatorname*{argmin}_{\M\in\mathbb M_{r,\mu_1}}\mathcal L(\M).
 \label{eq:of_same_argmin}
\end{align}
Moreover every candidate within $\omega_1/2$ of the latter infimum belongs
to $\mathcal A_{\rm fac}$ and obeys the optimal error bound.  For sufficiently
large fixed $A_0$, all such approximate minimizers also belong to
$\mathcal B_a$.
\end{scorollary}
\begin{proof}[Proof.]
Truth feasibility and Lemma~\ref{lem:loss_for_true} give
$\inf_{\mathbb M_{r,\mu_1}}\mathcal L\le\mathcal L(\M_0)
\le\omega_1/2$ for large $m$.  Hence every stated approximate minimizer
has loss at most $\omega_1$ and is retained by
\eqref{eq:of_factor_acceptance}.  This proves equality of infima and
argmin sets.  The low-loss certificate gives its distance at most
$C\epsilon_*$ from $\M_0$, and the triangle inequality with $\M^{(k_0)}$
proves ball inclusion when $A_0\ge2(C+C_a)$.  
\end{proof}

\subsection{Stationarity and computational cost}
\label{sec:stationarity_cost}
For local stationarity, suppose the nuclear sequence stays in
$\mathcal B_a$, uses the exact ball-constrained proximal step, and admits
sufficiently many backtracking reductions. Assume that its iterates and
trial segments lie in a compact subset of an open smooth region of the
finite-loss objective where $\nabla\mathcal L_{\rm cl}$ is
$L_g$-Lipschitz.
For a fixed maximum trial step $t_{\max}$ and a fixed backtracking ratio,
accepted steps are bounded below by a positive
$t_{\min}$ because \eqref{eq:of_nuclear_majorization} holds when
$t\le(1-c_A)/L_g$. Summing
\eqref{eq:of_nuclear_sufficient_decrease} over $K$ accepted updates after $k_0$ gives
\begin{align}
 \min_{k_0\le k<k_0+K}\left\|\frac{\M^{(k+1)}-\M^{(k)}}{t_k}\right\|_F^2
 \le\frac{2\{\Phi_\lambda(\M^{(k_0)})-\inf_{\mathcal B_a}\Phi_\lambda\}}
 {c_A t_{\min}K},\qquad K\ge1.
 \label{eq:of_stationarity_complexity}
\end{align}
The proximal optimality condition implies that every subsequential limit $\M^\ast$ satisfies
\begin{align*}
 0\in\nabla\mathcal L_{\rm cl}(\M^\ast)
       +\lambda\partial\|\M^\ast\|_*
       +N_{\mathcal B_a}(\M^\ast).
\end{align*}
At an interior limit the ball normal cone disappears, leaving the original
nuclear first-order condition.  The same smooth, strictly feasible
backtracking argument for the factor branch yields vanishing factor
stationarity residuals from \eqref{eq:of_factor_decrement} when step sizes
are bounded below and the damped Gram eigenvalues are bounded above and
away from zero.

Computationally, the exact gradient uses the residual SVD already needed
for the criterion and one extra scalar contraction.  The factor
preconditioner, rebalancing, and ball-membership calculation use small Gram
matrices at cost $O((m+n)r^2+r^3)$ beyond score construction.  The nuclear
proposal is the usual matrix SVT; only an active ball constraint requires
the scalar-multiplier search in \eqref{eq:of_ball_prox_scalar}.  It may
therefore require multiple SVTs.
The bulk-spectrum calculation and line-search loss evaluations remain part
of the total computational cost.  Inexact SVDs must preserve the stated feasibility and
objective and proximal-model tests, or carry explicit certified tolerances.
\end{revision}
% END MAJOR REVISION S06


% BEGIN MAJOR REVISION S07
\begin{revision}
\section{Oracle profile-reference analysis under heterogeneous noise}
\label{sec:heterogeneous_noise}
We consider coarse spectral calibration, fine calibration conditional on
the realized variance profile, and transfer to independently redrawn
profiles.  Each result specifies its calibration and sampling assumptions.

Let
\begin{align}
 \bH_{ij}=\sigma\sqrt{\nu_i\mu_j}\,\xi_{ij},\qquad
 \mathbb E\xi_{ij}=0,\quad\mathbb E\xi_{ij}^2=1,
 \quad \mathbb P(|\xi_{ij}|>x)\le2e^{-x/K}.
 \label{eq:vp_noise_model}
\end{align}
The displayed moments and tail bound hold conditionally on the profiles,
with a uniform constant $K$.  Conditional on the profiles the $\xi_{ij}$
are independent, and the
Bernoulli mask is independent of both profiles and standardized noise.
Thus $\operatorname{Var}(\bH_{ij}\mid\nu,\mu)=\sigma^2\nu_i\mu_j$.
The true matrix remains fixed and has the same rank/incoherence assumptions
as in the homogeneous model.


\begin{sassumption}[Variance profiles]
\label{assum:vp_regular}
For random profiles, the row and column coordinates are independent draws
from positive distributions with mean one and twelfth moments bounded by
a fixed constant.
There is a fixed $a>0$ such that
$\mathbb P(\nu\ge a)\wedge\mathbb P(\mu\ge a)\ge0.95$.
For a sufficiently small fixed $c_1$, let available deterministic
upper-tail envelopes satisfy
\begin{align}
 \mathbb P(\nu>U_\nu)\le c_1n^{-5},\qquad
 \mathbb P(\mu>U_\mu)\le c_1n^{-5},\qquad U_\nu,U_\mu\ge1,
 \label{eq:vp_U_definition}
\end{align}
and assume $U_\nu+U_\mu\le\log^{C_U}m$ for a fixed $C_U$.
The number of Monte Carlo reference draws satisfies $1\le T\le n$.
For deterministic profiles, replace the distributional clauses by the
profile bounds and the coordinate counts in \eqref{eq:vp_profile_event}
below.
\end{sassumption}

Write
\begin{align}
 H_{\nu,\mu}=\sqrt{U_\nu+U_\mu}=:H\ge1,
 \label{eq:vp_heterogeneity_scale}
\end{align}
and let $\mathcal E_{\rm vp}$ be the event that each data/reference
profile draw satisfies
\begin{align}
 \max_i\nu_i&\le U_\nu,\qquad \max_j\mu_j\le U_\mu,\nonumber\\
 m^{-1}\sum_i\nu_i&\le2,\qquad n^{-1}\sum_j\mu_j\le2,\nonumber\\
 |\{i:\nu_i\ge a\}|&\ge0.9m,\qquad |\{j:\mu_j\ge a\}|\ge0.9n.
 \label{eq:vp_profile_event}
\end{align}
The maxima have total exceptional probability at most $Cn^{-3}$ over
$T+1\le n+1$ draws because $m\asymp n$.
The twelfth-moment inequality for centered sums bounds failure of a mean
condition by $Cn^{-6}$ per draw; Bernoulli concentration handles the
coordinate counts.  Consequently $\mathbb P(\mathcal E_{\rm vp})
\ge1-Cn^{-3}$.  All conditional matrix estimates below are applied only
after fixing such profiles.
For one realization define
\begin{align}
 V_\nu=m^{-1}\sum_i\nu_i^2,\quad V_\mu=n^{-1}\sum_j\mu_j^2,
 \label{eq:vp_V_definition}\\
 h_{\nu,\mu}=\sqrt{U_\mu}+\sqrt{\rho U_\nu}
               +\rho^{1/4}(V_\nu V_\mu)^{1/4}\le CH.
 \label{eq:vp_edge_factor}
\end{align}
The quantity $h_{\nu,\mu}$ bounds the expected operator norm of the
covariance-matched Gaussian matrix.

For the covariance-matched realized-profile reference and the independently
redrawn reference, respectively, use
\begin{align}
 G_{ij}^{(t),\rm rp}&=\sqrt{\nu_i\mu_j}\,g_{ij}^{(t)},
 \label{eq:vp_gaussian_realized}\\
 G_{ij}^{(t),\rm dist}&=\sqrt{\bar\nu_i^{(t)}\bar\mu_j^{(t)}}\,g_{ij}^{(t)},
 \label{eq:vp_distribution_reference}\\
 \widehat\lambda_i^b&=T^{-1}\sum_{t=1}^T
        s_i\{P_{\bOmega}(G^{(t),b})/\sqrt m\},
        \quad b\in\{{\rm rp},{\rm dist}\}.
 \label{eq:vp_realized_reference}
\end{align}
The $g_{ij}^{(t)}$ are independent standard Gaussian variables, independent
of the observations.  The redrawn profiles are independent copies of the
profile laws.  For $I=\{\lceil n/3\rceil,\ldots,\lfloor2n/3\rfloor\}$,
define
\begin{align}
 \widehat\sigma_b(M)=
 \frac{\sum_{i\in I}s_i\{Y-P_{\bOmega}M\}}
      {\sqrt m\sum_{i\in I}\widehat\lambda_i^b}.
 \label{eq:vp_working_scale_general}
\end{align}
Both references match the variance of their specified profile; the chosen
profile and reference are held fixed during differentiation.


\subsection{Known-profile reference benchmark}
\label{sec:oracle_profile_benchmark}
For comparison with Theorem~\ref{main-thm:weighted_unknown_rate},
this subsection records the known-profile construction and bound. Its
reference spectrum is profile-specific, rather than homogeneous with
weighted residuals.

We construct the benchmark reference spectrum according to the available variance
information. When the realized multipliers are available, the Gaussian
reference uses the same profile. When only the profile laws are known,
new multipliers are sampled for each reference matrix. With independent
$g_{ij}^{(t)}\sim N(0,1)$, these two constructions are
\begin{align}
 G_{ij}^{(t),\mathrm{rp}}&=\sqrt{\nu_i\mu_j}\,g_{ij}^{(t)},&
 G_{ij}^{(t),\mathrm{dist}}
 &=\sqrt{\bar\nu_i^{(t)}\bar\mu_j^{(t)}}\,g_{ij}^{(t)},\nonumber\\
 \widehat\lambda_i^b
 &=\frac1T\sum_{t=1}^T
     \psi_i\{P_{\bOmega}(G^{(t),b})/\sqrt m\},&
 b&\in\{\mathrm{rp},\mathrm{dist}\},
 \label{eq:main_profile_reference}
\end{align}
where the barred profiles are independent copies of the profile laws.
The Gaussian draws are independent of the data conditional on the mask
and the profiles used. Each construction therefore matches the entry
variances specified by its chosen profile. The Monte Carlo reference is
held fixed during estimation. The working scale and loss become
\begin{align}
 \widehat\sigma_b(\M)&=
 \frac{\sum_{i\in I}\psi_i\{\Y-P_{\bOmega}(\M)\}}
      {\sqrt m\sum_{i\in I}\widehat\lambda_i^b},\nonumber\\
 \mathcal L_b(\M)&=p^{-1}\sum_i\omega_i
 \left\{\frac{\psi_i\{\Y-P_{\bOmega}(\M)\}}
              {\sqrt m\widehat\sigma_b(\M)}-\widehat\lambda_i^b\right\}^2.
 \label{eq:main_profile_loss}
\end{align}
Minimizing $\mathcal L_b$ over $\mathbb M_{r,\mu_1}$ gives
$\widetilde\M_{\mathrm{fac},b}$. For nuclear-norm regularization, use
$\widehat\sigma_{b,\mathrm{cl}}(\M)=
\min\{\widehat\sigma_b(\M),C_{\mathrm{clip}}
\widehat\sigma_{\mathrm{pil},b}\}$ and minimize
$\Phi_{\lambda,b}(\M)=\mathcal L_{b,\mathrm{cl}}(\M)+\lambda\|\M\|_*$,
where $\mathcal L_{b,\mathrm{cl}}$ is obtained by substituting the capped
scale in \eqref{eq:main_profile_loss}. The pilot is required to be
comparable to $\sigma$, and the fixed $C_{\mathrm{clip}}$ is chosen
sufficiently large for the cap to be inactive at the truth. A factor candidate satisfying the low-loss
condition below provides such a pilot through its working scale.

\begin{stheorem}
\label{thm:hetero_upper_bounds}
Under Assumptions~\ref{main-assum:target} and \ref{main-assum:hetero}, consider
independent Bernoulli$(p)$ observations. Suppose
\begin{align*}
 mp&\ge C H_{\nu,\mu}^{\,2}\log^{\alpha_{\mathrm{vp}}}m,
 \qquad\alpha_{\mathrm{vp}}>40+6C_U,\\
 np&\ge C\kappa^2\{\bar\mu r\log m+
                      H_{\nu,\mu}^{\,2}\bar\mu^2r^2\}.
\end{align*}
With probability at least $1-cn^{-3}$, every finite-loss
$\M\in\mathbb M_{r,\mu_1}$ satisfying
$\mathcal L_b(\M)\le C_LH_{\nu,\mu}^{\,2}\omega_1$ obeys
\begin{align*}
 \|\M-\M_0\|_F\le C H_{\nu,\mu}\sigma\sqrt{mr/p}.
\end{align*}
For either reference construction, this applies to every global minimizer
when $\mu_1\ge\mu_0$ and $\omega_1\ge c_\omega>0$.

For the nuclear objective, assume the calibrated pilot and inactive
truth cap described above. Let $\widehat{\M}$ satisfy the objective-comparison
and realized-curvature conditions
\eqref{main-eq:main_nuclear_objective_comparison}--\eqref{main-eq:main_nuclear_curvature},
with $\Phi_\lambda$ replaced by $\Phi_{\lambda,b}$.
Put $\ell_b=\mathcal L_{b,\mathrm{cl}}(\M_0)+\varepsilon_{\mathrm{opt}}$.
If
\begin{align*}
 \ell_b&\le C_\ell H_{\nu,\mu}^{\,2}r\bar\omega_r,\\
 c_-\frac{\ell_b}{H_{\nu,\mu}r}\sigma^{-1}\sqrt{p/m}
 &\le\lambda\le
 c_+H_{\nu,\mu}\bar\omega_r\sigma^{-1}\sqrt{p/m},
\end{align*}
then the same error bound holds for $\widehat{\M}$ on these events.
\end{stheorem}
The fixed-top-weight assertion uses only coarse spectral edge and bulk
control. For Gaussian standardized noise and the realized-profile
reference, a sharper calibration bound holds with probability at least
$1-cn^{-3}$:
\begin{align}
 \mathcal L_{\mathrm{rp}}(\M_0)
 \le C H_{\nu,\mu}^{\,2}\log^{-2}m.
 \label{eq:main_profile_fine_calibration}
\end{align}
This permits $\omega_1\ge c_\omega\log^{-2}m$ in the factor result.
For general non-Gaussian noise or a redrawn reference, the same smaller-weight
conclusion holds on the corresponding true-loss calibration event.

For a reference constructed from an estimated or resampled profile, the following identity describes how reference perturbations enter the loss. The weighted estimator in \eqref{main-eq:normalized_truncated_weights} instead uses a fixed homogeneous reference.

To describe the effect of fitting or resampling profiles, define
\begin{align}
 d_{\mathrm{prof},\omega}^2
 &=p^{-1}\sum_i\omega_i
     (\widehat\lambda_i^{\mathrm{alt}}-\widehat\lambda_i^{\mathrm{rp}})^2,
       \nonumber\\
 b_{\mathrm{prof}}
 &=\left|\frac{\sum_{i\in I}\widehat\lambda_i^{\mathrm{alt}}}
                  {\sum_{i\in I}\widehat\lambda_i^{\mathrm{rp}}}-1\right|.
 \label{eq:main_profile_discrepancies}
\end{align}
On the event \eqref{eq:main_profile_fine_calibration} and the coarse
edge/bulk event, suppose $b_{\mathrm{prof}}\le1/2$. The identity
$\widehat\sigma_{\mathrm{alt}}/\widehat\sigma_{\mathrm{rp}}=
(\sum_{i\in I}\widehat\lambda_i^{\mathrm{rp}})/
(\sum_{i\in I}\widehat\lambda_i^{\mathrm{alt}})$ then gives
\begin{align*}
 \mathcal L_{\mathrm{alt}}(\M_0)
 \le C\{H_{\nu,\mu}^{\,2}\log^{-2}m+
       d_{\mathrm{prof},\omega}^2+
       H_{\nu,\mu}^{\,2}b_{\mathrm{prof}}^2\}.
\end{align*}
The two discrepancy terms include uncertainty from the preliminary mean
estimate and the fitted variance profile. With the corresponding edge and
bulk bounds, an additional error of order at most
$H_{\nu,\mu}^{\,2}\omega_1$ preserves the factor low-loss guarantee;
the nuclear certificate uses the tolerance
$H_{\nu,\mu}^{\,2}r\bar\omega_r$ under the other premises of
Theorem~\ref{thm:hetero_upper_bounds}. The same derivative and proximal
formulas apply with the fitted reference held fixed.

The factor $H_{\nu,\mu}$ in Theorem~\ref{thm:hetero_upper_bounds}
quantifies the effect of the variance profile on the recovery bound.

\subsection{Coarse spectral calibration}
\begin{slemma}[Coarse data and reference calibration]
\label{lem:vp_distribution_coarse}
Under Assumption~\ref{assum:vp_regular}, a sufficient, unoptimized sparse
sampling condition is
\begin{align}
 mp\ge C H^2\log^{\alpha_{\rm vp}}m,
 \qquad \alpha_{\rm vp}>40+6C_U.
 \label{eq:vp_sparse_sampling}
\end{align}
With probability at least $1-cn^{-3}$, for either reference type,
\begin{align}
 \|P_{\bOmega}\bH\|/(\sigma\sqrt m)+\widehat\lambda_1^b
 &\le CH\sqrt p,\label{eq:vp_dist_coarse_edge}\\
 cn\sqrt p\le\sum_{i\in I}\widehat\lambda_i^b&\le Cn\sqrt p,
 \quad
 c\sigma n\sqrt{mp}\le\sum_{i\in I}s_i(P_{\bOmega}\bH)
       \le C\sigma n\sqrt{mp},\label{eq:vp_dist_coarse_bulk}\\
 c\sigma\le\widehat\sigma_b(M_0)&\le C\sigma,
 \qquad \mathcal L_b(M_0)\le CH^2.
 \label{eq:vp_dist_coarse_scale_loss}
\end{align}
The constants in the bulk and scale bounds depend on $a,\rho_-,K$ but
not on the growing upper envelope $H$.  Also, simultaneously for all $M$,
\begin{align}
 |\widehat\sigma_b(M)-\widehat\sigma_b(M_0)|
 \le\frac{C}{\sqrt{mnp}}\|P_{\bOmega}(M-M_0)\|_F.
 \label{eq:vp_working_scale_uniform}
\end{align}
\end{slemma}
\begin{proof}[Proof.]
Condition on $\mathcal E_{\rm vp}$.  We first establish a lower bound for the middle singular values.  Select the rows with $\nu_i\ge a$ and columns with $\mu_j\ge a$.
There are at least $0.9m$ and $0.9n$ of them.  On this submatrix, the two
diagonal multipliers have smallest singular value at least $\sqrt a$.
Therefore its singular values dominate $a$ times those of the standardized
masked submatrix.  If necessary transpose that submatrix.  Its minimum
dimension is at least $0.9n$ and its aspect ratio is bounded away from zero.
The homogeneous coarse rigidity lemma applies to its independent entries.
The index $\lfloor2n/3\rfloor$ is bounded away from its minimum dimension,
so the corresponding rooted Marchenko--Pastur quantile is bounded below
by a positive constant.  Singular-value interlacing for submatrices yields
\[
 s_{\lfloor2n/3\rfloor}(P_{\bOmega}\bH)
       \ge c\sigma\sqrt{mp}.
\]
This proves the data middle-sum lower bound and, with $\sigma=1$, each
reference draw's lower bound.

For the upper bulk bound, truncate standardized entries at $L=C_D\log n$.
The probability that any truncation is active in all $T+1$ matrices is at
most $Cn^{-D}$ after choosing $C_D$ large.  Conditional on a profile,
\[
 \mathbb E\sum_{ij}\Omega_{ij}\nu_i\mu_j\xi_{ij}^2
       =p\Big(\sum_i\nu_i\Big)\Big(\sum_j\mu_j\Big)\le4mnp.
\]
Truncated summands are bounded by $U_\nu U_\mu L^2$ and their total
variance is at most $C U_\nu U_\mu mnp$.  Scalar Bernstein therefore
bounds the squared Frobenius norm by $Cmnp$ with probability $1-n^{-D}$
when $mnp\ge C U_\nu U_\mu\log^3 n$, as is implied by
\eqref{eq:vp_sparse_sampling}.  Cauchy--Schwarz gives the upper middle
sum $Cn\sqrt{mp}$.  Averaging the reference bounds preserves them.

For the operator edge, apply the same logarithmic truncation, recentering,
and variance renormalization as in the homogeneous proof, followed by the
sparse bounded-summand comparison as in Section~\ref{sec:supp_sparse_rmt}.
For the matrix divided by $\sigma\sqrt{mp}$ its parameters are bounded by
\[
 \sigma_X\le CH,\qquad
 \sigma_{*,X}\le C\sqrt{U_\nu U_\mu/m},\qquad
 R_X\le C\sqrt{U_\nu U_\mu}\log n/\sqrt{mp}.
\]
The covariance-matched dense Gaussian comparator has expected norm at most
$Ch_{\nu,\mu}\le CH$ by
\citet{latala2005some}.  The comparison remainder with $t=C_D\log n$
is bounded by
\[
 C\{\sigma_{*,X}\sqrt t+R_X^{1/3}H^{2/3}t^{2/3}+R_Xt\}\le CH.
\]
The displayed sparse condition absorbs all profile logarithms, including
truncation errors.  This proves the edge bound for the data and for every
Gaussian reference draw.  A union bound over $T\le n$ is harmless by
first choosing $D>5$.
The ratios of bulk sums give the two-sided scale bound.  The edge bound
then gives $\mathcal L_b(M_0)\le CH^2$ because the weights sum to one.
Finally Mirsky's inequality and the reference bulk lower bound prove
\eqref{eq:vp_working_scale_uniform}, uniformly in the candidate.
\end{proof}

For $\nu,\mu\sim\chi_1^2$, these assumptions hold with a sufficiently
small fixed $a>0$. Indeed, the density is integrable at zero and has an
exponential upper tail, so all fixed moments are finite and
$U_\nu+U_\mu\le C\log n$. The bounds above therefore apply with
$H\le C\sqrt{\log n}$, including independently redrawn references.


\subsection{Fine calibration and reference transfer}
\begin{slemma}[Fine calibration for Gaussian noise with a realized profile]
\label{lem:vp_sparse_reference}
In addition to the coarse profile hypotheses, suppose the standardized
noise $\xi_{ij}$ is Gaussian and use the realized-profile reference.
Then on an event of probability at least $1-cn^{-3}$,
\begin{align}
 \max_i\left|\frac{s_i(P_{\bOmega}\bH)}{\sigma\sqrt m}
                      -\widehat\lambda_i^{\rm rp}\right|
 &\le C\sqrt{\frac{U_\nu U_\mu\log n}{m}}
 \le C\sqrt p\,\delta_m\le CH\sqrt p\,\delta_m,
 \quad\delta_m=\log^{-1}m,\label{eq:vp_realized_match}\\
 \left|\sum_{i\in I}\left\{
      \frac{s_i(P_{\bOmega}\bH)}{\sigma\sqrt m}
                      -\widehat\lambda_i^{\rm rp}\right\}\right|
 &\le Cn\sqrt p\,\delta_m,\label{eq:vp_bulk_average_match}\\
 \|P_{\bOmega}\bH\|/(\sigma\sqrt m)+\widehat\lambda_1^{\rm rp}
 &\le CH\sqrt p,\label{eq:vp_edge_explicit}\\
 cn\sqrt p\le\sum_{i\in I}\widehat\lambda_i^{\rm rp}
 &\le Cn\sqrt p,\qquad
 |\widehat\sigma_{\rm rp}(M_0)/\sigma-1|\le C\delta_m.
 \label{eq:vp_bulk_explicit}
\end{align}
On the same event,
\begin{align}
 \mathcal L_{\rm rp}(M_0)\le CH^2\delta_m^2.
 \label{eq:vp_loss_true}
\end{align}
\end{slemma}
\begin{proof}[Proof.]
Condition on the realized profile and mask.  The data matrix divided by
$\sigma\sqrt m$ and each reference draw have exactly the same
conditional Gaussian law.  The $i$th singular value of
$P_{\bOmega}(\operatorname{Diag}(\nu)^{1/2}G
\operatorname{Diag}(\mu)^{1/2})/\sqrt m$ is
$\sqrt{U_\nu U_\mu/m}$-Lipschitz in the Gaussian entries.
Its Monte Carlo average is
$\sqrt{U_\nu U_\mu/(mT)}$-Lipschitz.  Concentrate both around the same
conditional mean and take a union bound over $i\le n$.
The resulting first bound in \eqref{eq:vp_realized_match} is at most
$C\sqrt p\delta_m$ when $mp\ge C U_\nu U_\mu\log^3 m$;
this follows from \eqref{eq:vp_sparse_sampling}.
Sum over $I$ for the averaged comparison, and use the already proved
coarse edge/bulk event.  Dividing the averaged comparison by the reference
bulk sum proves the scale assertion.  

Put $\theta_i=s_i(P_{\bOmega}\bH)/(\sigma\sqrt m)$.
The averaged comparison and the bulk lower bound give
$|\widehat\sigma_{\rm rp}(M_0)/\sigma-1|\le C\delta_m$.
Expanding
$\sigma\theta_i/\widehat\sigma_{\rm rp}(M_0)-\widehat\lambda_i^{\rm rp}$
as a spectral discrepancy plus a common scale discrepancy gives
\[
 \mathcal L_{\rm rp}(M_0)
 \le2p^{-1}\max_i|\theta_i-\widehat\lambda_i^{\rm rp}|^2
 +2p^{-1}\theta_1^2
       |\sigma/\widehat\sigma_{\rm rp}(M_0)-1|^2
 \le CH^2\delta_m^2.
\]
The scale perturbation was proved in the coarse lemma and needs no fine
calibration or independence of the candidate.
\end{proof}


For non-Gaussian standardized noise, define $\mathcal E_{\rm cal}^{\rm rp}$
to be the coarse edge/bulk event together with the bounds
\eqref{eq:vp_realized_match} and \eqref{eq:vp_bulk_average_match}, using
the weaker permitted $CH\sqrt p\delta_m$ bound for the maximum discrepancy.
Fine-calibration statements for this model are conditional on
$\mathcal E_{\rm cal}^{\rm rp}$.

On this event, the same loss expansion in the preceding proof yields
\eqref{eq:vp_loss_true}, while the uniform scale bound remains
\eqref{eq:vp_working_scale_uniform}.

Let $\widehat{\bm\lambda}^{\mathrm{alt}}$ be either the distribution-redrawn
reference or any other nonnegative ordered reference, including one constructed
from estimated profiles.  Define $\widehat\sigma_{\mathrm{alt}}$ and
$\mathcal L_{\mathrm{alt}}$ by substituting this reference into the same
bulk-scale and loss formulas.  Define the discrepancies
\begin{align}
 d_{\rm prof,\omega}^2
 &=p^{-1}\sum_i\omega_i
       (\widehat\lambda_i^{\mathrm{alt}}-\widehat\lambda_i^{\rm rp})^2,
 \label{eq:vp_profile_weighted_mismatch}\\
 b_{\rm prof}
 &=\left|\frac{\sum_{i\in I}(\widehat\lambda_i^{\mathrm{alt}}
                          -\widehat\lambda_i^{\rm rp})}
                    {\sum_{i\in I}\widehat\lambda_i^{\rm rp}}\right|.
 \label{eq:vp_profile_bulk_mismatch}
\end{align}
\begin{slemma}[Transfer to a substituted reference]
\label{lem:vp_distribution_redraw}
On $\mathcal E_{\rm cal}^{\rm rp}$, if $b_{\rm prof}\le1/2$,
then
\begin{align}
 |\widehat\sigma_{\mathrm{alt}}(M_0)/\sigma-1|
 &\le C(\delta_m+b_{\rm prof}),
 \label{eq:vp_distribution_sigma}\\
 \mathcal L_{\mathrm{alt}}(M_0)
 &\le C\{H^2\delta_m^2+d_{\rm prof,\omega}^2
                           +H^2b_{\rm prof}^2\}.
 \label{eq:vp_distribution_loss}
\end{align}
\end{slemma}
\begin{proof}[Proof.]
Write $S_b=\sum_{i\in I}\widehat\lambda_i^b$.
The numerator of the two working scales is identical, so
$\widehat\sigma_{\mathrm{alt}}(M_0)/\widehat\sigma_{\rm rp}(M_0)
=S_{\rm rp}/S_{\dagger}$.
The bulk discrepancy controls this ratio and its reciprocal by fixed
constants, and bounds their difference from one by $Cb_{\rm prof}$.
For each normalized residual discrepancy add and subtract the realized-profile
reference, then apply $(x+y+z)^2\le3(x^2+y^2+z^2)$.
The three terms are the realized-profile loss, the weighted reference
mismatch, and the common scale mismatch times the squared edge bound.
This proves both assertions.  The argument is deterministic given the two references and applies equally
to a reference estimated from the observations.
\end{proof}

In particular, the fine loss bound is preserved whenever
$b_{\rm prof}\le1/2$ and
$d_{\rm prof,\omega}^2+H^2b_{\rm prof}^2
\le CH^2\delta_m^2$.

\subsection{Statistical error bounds}
We prove both parts of Theorem~\ref{thm:hetero_upper_bounds} on
the calibration and sampling events above. The argument also gives the
stated refinements when a fine-calibration or substituted-reference bound
is available.

\begin{proof}[Proof of Theorem~\ref{thm:hetero_upper_bounds}.]
\textit{Rank-constrained estimation.} On the coarse calibration event, suppose $12r\le n$ and
\begin{align}
 np\ge C\kappa^2\{\bar\mu r\log m+H^2\bar\mu^2r^2\}.
 \label{eq:vp_structural_sampling}
\end{align}
Then every finite-loss $M\in\mathbb M_{r,\mu_1}$ with
$\mathcal L_b(M)\le C_LH^2\omega_1$ satisfies
\begin{align}
 \|M-M_0\|_F\le CH\sigma\sqrt{mr/p}.
 \label{eq:vp_factor_rate}
\end{align}
Consequently, when $\mu_1\ge\mu_0$, this applies to every global minimizer, when it exists, if
$\mathcal L_b(M_0)\le C_LH^2\omega_1$.
For either reference this latter condition follows from coarse calibration
when $\omega_1\ge c_\omega>0$.
On the fine realized-profile event it suffices that
\begin{align}
 \omega_1\ge c_\omega\delta_m^2.
 \label{eq:vp_factor_weight_realized}
\end{align}
Fine distribution-redraw calibration additionally requires
\begin{align}
 b_{\rm prof}\le1/2,\qquad
 d_{\rm prof,\omega}^2+H^2b_{\rm prof}^2
       \le CH^2\omega_1.
 \label{eq:vp_factor_redraw_condition}
\end{align}
For general non-Gaussian profiles the weak-weight assertions are conditional
on the fine-calibration event; the fixed-top-weight coarse assertion is
unconditional under the stated model and sampling conditions.

Let $d=\|M-M_0\|_F$, $\delta=\|M-M_0\|$, $h=\sigma\sqrt{m/p}$,
$h_H=Hh$, and $\beta=\bar\mu r/\sqrt{np}$.
The same residual-operator and Mirsky argument as in
\eqref{eq:factor_energy_upper} gives
\begin{align}
 Q_\Omega(M-M_0)\le Cp\sqrt r h_Hd+C\frac{H^2r}{n}d^2.
 \label{eq:variance_profile_sampled_upper}
\end{align}
Indeed the low-loss bound makes the residual norm at most
$CH\widehat\sigma_b(M)\sqrt{mp}$, while
$\widehat\sigma_b(M)\le C\sigma+C\sqrt{Q_\Omega}/\sqrt{mnp}$;
Young's inequality produces the displayed $H^2r/n$ term.

For operator localization use \eqref{eq:sampling_endpoint_mask_bound}.
Rank interlacing on the residual bulk gives
$\widehat\sigma_b(M)\sqrt{m/p}\le C(Hh+
\|T_\Omega(M-M_0)\|)$.
The low-loss residual bound therefore implies
\[
 \delta\le CH^2h+CH\beta(s_1+\delta).
\]
The sampling condition makes $CH\beta\le1/2$, and hence
$\|M\|\le C(s_1+Hh_H)$.
If $h_H<s_1$, apply Lemma~\ref{lem:global_relative_sampling} with $C_M=CH$.
Its proof gives explicitly
$a_g\le C\kappa\{\theta+(1+C_M)\Lambda\bar\mu r/\sqrt{mn}\}$,
so \eqref{eq:vp_structural_sampling} makes $a_g\le1/2$.
Combining its lower energy bound with \eqref{eq:variance_profile_sampled_upper}
and absorbing $CH^2r/(np)$ gives $d\le C\sqrt r h_H$.

If $h_H\ge s_1$, the operator localization gives $\|M\|\le CHh_H$.
The endpoint incoherence bound is
$\|M-M_0\|_\infty\le C\bar\mu rHh_H/\sqrt{mn}$.
The coarse rank-$2r$ RSC and \eqref{eq:variance_profile_sampled_upper} imply,
after the same absorption,
\[
 d^2\le C\sqrt r h_Hd+
            C\frac{H^2\bar\mu^2r^3}{np}h_H^2
       \le C\sqrt r h_Hd+Crh_H^2.
\]
Solving this quadratic proves the rate in the remaining case.
The sufficient true-loss conditions follow from the coarse, fine, or
redraw loss bound, respectively.  

\textit{Pilot scale and nuclear-norm estimation.} Let $M_{\rm pil}$ satisfy the low-loss premise established in the factor part.
Set $\widehat\sigma_{\rm pil,vp}=\widehat\sigma_b(M_{\rm pil})$.
On the same event and for a sufficiently large constant in
\eqref{eq:vp_structural_sampling},
\begin{align}
 c_{\rm pil}\sigma\le\widehat\sigma_{\rm pil,vp}\le C_{\rm pil}\sigma,
 \qquad c_0\sigma\le\widehat\sigma_b(M_0)\le C_0\sigma.
 \label{eq:vp_pilot_event}
\end{align}
For $C_{\rm clip}>C_0/c_{\rm pil}$ define
\begin{align}
 \widehat\sigma_{b,\rm cl}(M)=
 \min\{\widehat\sigma_b(M),C_{\rm clip}\widehat\sigma_{\rm pil,vp}\}.
 \label{eq:vp_clipped_scale}
\end{align}
Clipping is inactive at $M_0$, and uniformly in $M$,
\begin{align}
 \widehat\sigma_{b,\rm cl}(M)\le C\sigma.
 \label{eq:vp_clipped_scale_uniform}
\end{align}

The factor rate and \eqref{eq:variance_profile_sampled_upper} give
$Q_\Omega(M_{\rm pil}-M_0)\le Cp r h_H^2=CH^2\sigma^2mr$.
Thus the uniform scale perturbation is at most
$CH\sigma\sqrt{r/(np)}$, a small fixed multiple of $\sigma$ by the
explicit structural sampling condition.  Combine this with the coarse
truth-scale bounds.  This proves the pilot and clipping assertions.

Use the clipped heterogeneous loss
\begin{align}
 \mathcal L_{b,\rm cl}(M)=\frac1p\sum_i\omega_i
 \left\{\frac{s_i(Y-P_{\bOmega}M)}{\sqrt m\widehat\sigma_{b,\rm cl}(M)}
                    -\widehat\lambda_i^b\right\}^2.
 \label{eq:vp_clipped_loss}
\end{align}

Let $\widehat{M}$ be a finite-loss candidate satisfying
\[
 \mathcal L_{b,\rm cl}(\widehat{M})+\lambda\|\widehat{M}\|_*
 \le\mathcal L_{b,\rm cl}(M_0)+\lambda\|M_0\|_*+\varepsilon_{\rm opt}.
\]
Put $\ell=\mathcal L_{b,\rm cl}(M_0)+\varepsilon_{\rm opt}$ and
$\bar\omega_r=\min_{i\le r}\omega_i>0$.  Assume the coarse edge/bulk
and pilot events, and
\begin{align}
 \|P_{\bOmega}(\widehat{M}-M_0)\|_F^2
 &\ge c_{\rm curv}p\|\widehat{M}-M_0\|_F^2,
 \label{eq:vp_nuclear_localization}\\
 \ell&\le C_\ell H^2r\bar\omega_r,
 \label{eq:vp_nuclear_balance}\\
 c_-\frac{\ell}{Hr}\sigma^{-1}\sqrt{p/m}
 \le\lambda&\le c_+H\bar\omega_r\sigma^{-1}\sqrt{p/m}.
 \label{eq:vp_nuclear_tuning}
\end{align}
Then
\begin{align}
 \|\widehat{M}-M_0\|_F\le CH\sigma\sqrt{mr/p}.
 \label{eq:vp_nuclear_rate}
\end{align}
For coarse calibration, $r\bar\omega_r\ge c>0$ and
$\varepsilon_{\rm opt}\le CH^2$ are sufficient for the balance condition.
On a fine realized-profile event they can be replaced by
$r\bar\omega_r\ge c\delta_m^2$ and a matching tolerance.
For fine redrawn calibration one additionally needs
\begin{align}
 d_{\rm prof,\omega}^2+H^2b_{\rm prof}^2
            \le CH^2r\bar\omega_r,\qquad b_{\rm prof}\le1/2.
 \label{eq:vp_nuclear_redraw_condition}
\end{align}

The nuclear cone and block decomposition in
Theorem~\ref{main-thm:upper_bound_for_convexlambda} are deterministic.
With the profile edge $Hh$ and the clipped scale bounded by $C\sigma$,
that proof yields the four terms
\[
 \frac{\ell}{\lambda\sqrt r},\qquad
 H\sqrt r h,\qquad h\sqrt{\ell/\bar\omega_r},\qquad
 h^2\lambda\sqrt r/\bar\omega_r.
\]
The lower tuning bound, the balance condition, and the upper tuning bound
control the first, third, and fourth terms, respectively, at scale
$H\sqrt r h$.  The available coarse or fine truth-loss bound supplies
the stated specializations.  
\end{proof}


\subsection{Estimated-profile reference as an alternative construction}
\label{sec:estimated_profiles}
The realized-profile reference can be implemented by substituting estimates
of its positive row and column multipliers.  For example, given a preliminary
mean estimate $\check M$, write $E=Y-P_{\bOmega}(\check M)$ and let
$d_i=\sum_j\Omega_{ij}$, $e_j=\sum_i\Omega_{ij}$.  When these degrees and
the total residual energy are positive, form
\begin{align*}
 q_i&=d_i^{-1}\sum_j\Omega_{ij}E_{ij}^2,&
 t_j&=e_j^{-1}\sum_i\Omega_{ij}E_{ij}^2,&
 q&=|\Omega|^{-1}\sum_{ij}\Omega_{ij}E_{ij}^2.
\end{align*}
One possible initialization of the variance multipliers is $q_i/q$ and
$t_j/q$, followed by positive truncation and normalization of their respective
coordinate means to one.  More generally, one may fit a nonnegative separable
variance model to the squared residuals.  The reference matrices are then
\begin{align}
 \widehat G^{(t)}_{ij}
 =\sqrt{\widehat\nu_i\widehat\mu_j}\,g_{ij}^{(t)},\qquad
 \widehat\lambda_i^{\mathrm{alt}}
 =T^{-1}\sum_{t=1}^{T}s_i\{P_{\bOmega}(\widehat G^{(t)})/\sqrt m\}.
 \label{eq:estimated_profile_reference}
\end{align}
The fitted multipliers and this reference are held fixed during optimization;
the working scale is still recomputed from the residual bulk.  Empirical
coordinate distributions may instead be resampled, but this is the distinct
redrawn-reference construction, with its own discrepancy terms.

The transfer bound in Lemma~\ref{lem:vp_distribution_redraw} gives,
on its stated baseline event,
\begin{align*}
 \mathcal L_{\mathrm{alt}}(M_0)
 \le C\{H^2\log^{-2}m+d_{\rm prof,\omega}^2
                         +H^2b_{\rm prof}^2\}.
\end{align*}
The two discrepancy terms quantify spectral and bulk-reference error.
For the factor
certificate, it is sufficient that their sum is at most $CH^2\omega_1$,
that $b_{\rm prof}\le1/2$, and that the edge and structural conditions
of Theorem~\ref{thm:hetero_upper_bounds} hold.  Nuclear certification uses the
corresponding tolerance $CH^2r\bar\omega_r$ and its other stated premises.


The derivative and proximal formulas in Section~\ref{sec:estimator_algorithms}
apply to any fixed profile-specific reference.  Their statistical acceptance
thresholds must use the profile-specific error and loss scales above.
\end{revision}
% END MAJOR REVISION S07



% BEGIN MAJOR REVISION S08
\begin{revision}
\section{Distribution-unknown normalized weighted recovery}
\label{sec:weighted_certification}
This section specifies the procedure and proves
Theorem~\ref{main-thm:weighted_unknown_rate}. Calibration uses only the
unweighted observations, the same-mask Gaussian reference, and the
regularity-class envelopes. Estimated weights may depend on all the
observations.

\subsection{Calibration and residual-scale fitting}
Let $H_n=(U_\nu+U_\mu)^{1/2}$. Valid deterministic envelopes are selected
from the regularity class, not from unobserved profile realizations.
For example, uniform sub-exponential upper tails for the multipliers
allow $U_\nu,U_\mu=C\log\{e(m+n)\}$.

For any matrix $M$, define the unweighted bulk scale
\begin{align*}
 s_h(M)=\frac{\sum_{i\in I}\psi_i\{Y-P_{\bOmega}(M)\}}
                  {\sqrt m\sum_{i\in I}\widehat\lambda_i}.
\end{align*}
Write $Z=p^{-1}Y$, $\widehat s_j=\psi_j(Z)$,
$h_0=s_h(0)\sqrt{m/p}$, and
$M_b=P_B\{\operatorname{Proj}_r(Z)\}$, where
$B=K_0\sqrt{\mu_0r}\,\widehat s_1$ and $P_B$ is
\eqref{main-eq:main_initial_projection}. Use sufficiently large fixed $K_0$ and
sufficiently small fixed $t_0$.
If
\begin{align}
 H_n\sqrt r\,h_0\ge t_0\widehat s_r,
 \label{eq:hetero_anchor_branch}
\end{align}
set $M_a=M_b$ and $\widehat s_a=s_h(0)$. Otherwise start at $M_b$ and
run the unweighted recovery score \eqref{main-eq:main_initial_score}, with
step $1/2$, zero damping, and projection $P_B$, until the first index $k$
satisfying
\begin{align}
 \|M^{(k+1)}-M^{(k)}\|_F
 \le C_{\rm stop}H_n\sqrt r\,s_h(M^{(k)})\sqrt{m/p}.
 \label{eq:hetero_anchor_stop}
\end{align}
Set $M_a=M^{(k+1)}$ and $\widehat s_a=s_h(M_a)$. In the direct branch it is
important to use $s_h(0)$, not to assume a scale bound at a possibly
rank-deficient projected matrix. The theoretical construction does not
impose a fixed cap on this preliminary stopping time. A resource limit
reached first is reported as an uncertified initialization, not as a
successful application of the theorem. For scale fitting one may take
$\check M=M_a$, or continue the unweighted numerical pilot; the theorem
below does not assume that this optional continuation improves $M_a$.


The relative-scale fit used to form a weighted proposal is as follows.
Let $\check M$ be the unweighted recovery output used for scale fitting,
and put $R_{ij}=Y_{ij}-\check M_{ij}$ on observed entries.  A two-way median
fit is specified by
\begin{align}
 \ell_{ij}&=\log\{|R_{ij}|\vee(\varepsilon_{\log}\widehat s_a)\},\nonumber\\
 c_j&=\operatorname{median}_{i:\Omega_{ij}=1}\ell_{ij},\qquad
 a_i=\operatorname{median}_{j:\Omega_{ij}=1}(\ell_{ij}-c_j),\nonumber\\
 u&=\operatorname{median}_{(i,j)\in\Omega}
                         (\ell_{ij}-a_i-c_j),\qquad
 S_{ij}=\exp(a_i+c_j+u).
 \label{eq:weighted_median_profile}
\end{align}
Here $\varepsilon_{\log}>0$ is a reported numerical tolerance and
$\widehat s_a>0$ is the calibration scale defined below.  Repeated
centering of these row and column effects leaves their fitted shape
unchanged. Empty rows or columns or a nonfinite scale fit trigger the
unweighted fallback. A common additive correction to all $\ell_{ij}$ is
unnecessary: its multiplicative effect on $S$ is removed by normalization.


Form the normalized, floored weights by
\eqref{main-eq:normalized_truncated_weights}; their quantile and median
conventions are fixed, monotone, and positively homogeneous, and use all
$mn$ fitted entries. The default quantile is $v=1/4$.

\subsection{Weighted proposals and optional refinement}
The heterogeneous implementation in Section~\ref{main-sec:4.4} uses a positive,
data-dependent matrix $W$, held fixed after its construction.  Write
\begin{align*}
 E_W(M)&=W\odot\{Y-P_{\bOmega}(M)\}
             =U_W\operatorname{diag}(d_i^W)V_W',\\
 s_W(M)&=\frac{\sum_{i\in I}d_i^W}
                   {\sqrt m\sum_{i\in I}\widehat\lambda_i},\qquad
 Q_W(M)=\sqrt m\,s_W(M)U_W
                \operatorname{diag}(\widehat\lambda_i)V_W'.
\end{align*}
A weighted recovery proposal uses the score
\begin{align}
 G_W^{\rm rec}(M)
 &=\frac1p W\odot\{-E_W(M)+\alpha_W(M)P_{\bOmega}Q_W(M)\},
 \label{eq:weighted_recovery_score}\\
 \alpha_W(M)
 &=1\wedge\frac{C_Qs_W(M)\sqrt{mp}}
                         {\|P_{\bOmega}Q_W(M)\|},
 \label{eq:weighted_pre_cap}
\end{align}
with $\alpha_W=1$ for a zero denominator.  The factor proposals use
\eqref{main-eq:main_factor_direction} with $G_W^{\rm rec}$ in place of $G$,
followed by the same row--column projection and balancing.  This score is
not asserted to be the full derivative of the candidate-dependent-scale
spectral loss. The cap in \eqref{eq:weighted_pre_cap} is explicitly applied
before the outer multiplication by $W$; applying it afterwards defines a
different proposal rule.

For the numerical proposal sequence, the initial step is
\begin{align}
 t_W=\frac12\min\left\{1,
    \frac{\|Y-P_{\bOmega}(M_{\rm start})\|}
         {\|W\odot\{Y-P_{\bOmega}(M_{\rm start})\}\|}\right\},
 \label{eq:weighted_auto_step}
\end{align}
where $M_{\rm start}$ is the data-based starting matrix of the chosen
proposal generator, kept fixed across paired weight comparisons. The norms here are
spectral norms, consistently with the notation of the paper. A zero
residual ends the proposal sequence. A fixed iteration budget specifies a
numerical output, not a convergence certificate for
\eqref{eq:weighted_auto_step}.

An optional refinement freezes $b_W=\sqrt m\,s_W(M_{\rm start})>0$ at its
starting matrix and minimizes
\begin{align*}
 \mathcal L_{W,b_W}(M)
 &=\frac1p\sum_i\omega_i
       \left\{\frac{d_i^W(M)}{b_W}-\widehat\lambda_i\right\}^{2}.
\end{align*}
At a differentiability point its matrix gradient is
\begin{align}
 \nabla\mathcal L_{W,b_W}(M)
 &=-P_{\bOmega}\!\left[W\odot
 U_W\operatorname{diag}\left\{
 \frac{2\omega_i}{pb_W}
 \left(\frac{d_i^W}{b_W}-\widehat\lambda_i\right)\right\}V_W'\right].
 \label{eq:weighted_fixed_gradient}
\end{align}
Factor backtracking uses this fixed-scale objective. Nuclear refinement
uses the same gradient and a singular-value-thresholding proposal for
$\mathcal L_{W,b_W}(M)+\lambda\|M\|_*$. These operations differ from
recomputing the working scale at every trial point in the homogeneous
objective above. Both types of numerical output are subjected to the
observable output rule in Section~\ref{main-sec:4.4}. Its statistical guarantee
also applies to alternative weighted proposal implementations and does
not require a contraction claim for the unguarded weighted trajectory.

\subsection{Model regularity and algorithm conventions}
Assumption~\ref{main-assum:hetero} concerns only the distributions.
The rank, incoherence, sampling, and tail-class conditions are imposed
on the statistical model. The calibration branches, residual-log
regularization, weight normalization, candidate budgets, and final
stabilization are parts of the procedure, not assumptions on fitted
random objects. The stopping construction is fully specified above;
its termination and statistical accuracy are proved below. Tail-class
constants are tuning inputs, not knowledge of a parametric law.


\subsection{Distributional calibration without a known profile law}
Put $H_n=(U_\nu+U_\mu)^{1/2}$, $h=\sigma\sqrt{m/p}$, and
$\epsilon_H=H_n\sqrt r\,h$. The deterministic envelopes $U_\nu,U_\mu$ are
upper-tail bounds for a class of unknown distributions. They are not the
unobserved sample maxima. Assumption~\ref{main-assum:hetero}, the moment
argument preceding \eqref{eq:vp_profile_event}, and a union bound imply
that the realized row and column profiles satisfy
\eqref{eq:vp_profile_event} with probability at least $1-cn^{-3}$.
Only the data profile is drawn in this construction; no multipliers are
redrawn for the homogeneous reference. In particular, bounded realized
noise entries are not an additional model assumption.

Let $\widehat\lambda$ be the homogeneous same-mask reference used by the
algorithm, and define
\begin{align}
 s_h(M)&=\frac{\sum_{i\in I}s_i\{Y-P_\Omega M\}}
                 {\sqrt m\sum_{i\in I}\widehat\lambda_i},&
 I&=\{\lceil n/3\rceil,\ldots,\lfloor2n/3\rfloor\}.
 \label{eq:wc_homogeneous_scale}
\end{align}

\begin{slemma}[Coarse calibration for an unknown profile]
\label{lem:wc_coarse}
Under the profile, noise, mask, and sparse-sampling conditions of
Theorem~\ref{main-thm:weighted_unknown_rate}, with probability at least
$1-cn^{-3}$ the following hold simultaneously:
\begin{align}
 \|p^{-1}P_\Omega H\|&\le C H_n h,\qquad
 \widehat\lambda_1\le C\sqrt p,\label{eq:wc_noise_edge}\\
 cn\sqrt p\le\sum_{i\in I}\widehat\lambda_i&\le Cn\sqrt p,
 \qquad c_0\sigma\le s_h(M_0)\le C_0\sigma,
 \label{eq:wc_bulk_scale}\\
 |s_h(M)-s_h(M_0)|
 &\le\frac{C\|P_\Omega(M-M_0)\|_F}{\sqrt{mnp}}
 \quad\text{for every }M.
 \label{eq:wc_scale_uniform}
\end{align}
Write
\begin{align*}
 Z=p^{-1}Y,
 \quad b_{\rm sig}=\|(p^{-1}P_\Omega-\operatorname{Id})M_0\|,
 \quad h_0=s_h(0)\sqrt{m/p}.
\end{align*}
On the same event,
\begin{align}
 ch-Cb_{\rm sig}\le h_0&\le Ch+Cb_{\rm sig},
 \label{eq:wc_raw_scale}\\
 b_{\rm sig}&\le Cs_1\left\{
 \sqrt{\frac{\mu_0r\log m}{np}}+
 \frac{\mu_0r\log m}{p\sqrt{mn}}\right\}.
 \label{eq:wc_signal_mask}
\end{align}
Constants in the two-sided bulk and scale bounds do not grow with $H_n$.
\end{slemma}
\begin{proof}
Condition first on the good profile event. The data-edge argument of
Lemma~\ref{lem:vp_distribution_coarse} gives
\eqref{eq:wc_noise_edge} for $H$. The reference has unit row and column
multipliers, so its edge and middle-sum bounds follow either from the
homogeneous coarse spectral lemma or from the same proof with both
multipliers equal to one. This does not require the data and the
reference to have the same variance profile.

For the data lower bulk, select the rows and columns with multipliers at
least the fixed constant $a$ from \eqref{eq:vp_profile_event}. There are
at least $0.9m$ and $0.9n$ such indices. On that rectangle the noise is a
left and right diagonal scaling of the standardized masked noise.
The singular-value product inequality and submatrix interlacing bound
$s_{\lfloor3n/4\rfloor}(p^{-1}P_\Omega H)$ below by $ch$.
The homogeneous coarse lemma applies to the standardized submatrix,
whose minimum dimension exceeds $0.9n$; the index $3n/4$ stays a fixed
distance from its lower edge. This is the same argument as in
Lemma~\ref{lem:vp_distribution_coarse}, with $3n/4$ in place of $2n/3$.
Rounding is absorbed for sufficiently large dimensions.

The proof of that lemma also gives
$\|p^{-1}P_\Omega H\|_F\le C\sqrt n\,h$; the empirical means of the
profile coordinates are bounded on the profile event. Consequently
sums over any $O(n)$ consecutive noise singular values are at most $Cnh$.
These facts give the data middle-sum bounds and then
\eqref{eq:wc_bulk_scale}. Mirsky's inequality gives
\eqref{eq:wc_scale_uniform} exactly as in
\eqref{eq:mirsky_middle_sum}. This last statement is deterministic on the
common reference event and is uniform in $M$.

Put $N=p^{-1}P_\Omega H$ and
$S_0=(p^{-1}P_\Omega-\operatorname{Id})M_0$, so
$Z=M_0+N+S_0$. Rank perturbation and Weyl's inequality imply, for $i\in I$,
\[
 s_{i+r}(N)-b_{\rm sig}\le s_i(Z)
                   \le s_{i-r}(N)+b_{\rm sig}.
\]
Since $12r\le n$, the shifted indices lie between $n/4$ and $3n/4$,
up to rounding. Sum and divide by
$\sum_{i\in I}\widehat\lambda_i/\sqrt p\asymp n$ to obtain
\eqref{eq:wc_raw_scale}. No fine ordered-spectrum calibration is used.
Finally, the fixed-signal matrix Bernstein calculation in
\eqref{eq:init_signal_concentration} proves \eqref{eq:wc_signal_mask}.
Intersecting the finite collection of events proves the lemma. There is
no conditioning on fitted weights at any step.
\end{proof}

\subsection{A data-based anchor at the heterogeneous statistical scale}
Use the projected spectral matrix $M_b$ and branch rule specified above. The calibration branch uses zero damping and step $\eta=1/2$.
These settings are independent of the auto step used later by the weighted
numerical procedure. The direct branch does not require hard thresholding
or an incoherence assertion for its output, because the output safeguard
below uses its Frobenius accuracy directly.

\begin{slemma}[Construction and accuracy of the anchor]
\label{lem:wc_anchor}
Under the hypotheses of Theorem~\ref{main-thm:weighted_unknown_rate}, the
construction in \eqref{eq:hetero_anchor_branch}--
\eqref{eq:hetero_anchor_stop} succeeds on an event of probability at
least $1-cn^{-3}$ and gives
\begin{align}
 \|M_a-M_0\|_F&\le C_a\epsilon_H,&
 c_a\sigma&\le\widehat s_a\le C_a'\sigma.
 \label{eq:wc_anchor_result}
\end{align}
In its iterative branch,
\begin{align}
 f_{k+1}&\le\rho_a f_k+C_R\epsilon_H,
               \qquad 0<\rho_a\le3/4,\label{eq:wc_rec}\\
 f_k&\le\rho_a^k f_0+
        \frac{C_R(1-\rho_a^k)}{1-\rho_a}\epsilon_H,
 \label{eq:wc_geometric}
\end{align}
up to the stopping point, and the update count is at most
\begin{align}
 C\{1+\log_+(f_0/\epsilon_H)\}.
 \label{eq:hetero_anchor_complexity}
\end{align}
\end{slemma}
\begin{proof}
Work on the common event of Lemma~\ref{lem:wc_coarse} and the uniform
mask bounds in Lemma~\ref{lem:mask_geometry}. All sampling inequalities
used below are uniform over the iterates.

\textit{1. Signal-mask error and the observable branch.}
The first term in \eqref{main-eq:hetero_anchor_samples} makes, for any
prescribed sufficiently small fixed $c_b>0$,
\begin{align}
 b_{\rm sig}\le\frac{c_bs_\star}{H_n\sqrt r}.
 \label{eq:wc_mask_small}
\end{align}
Indeed, multiplying the first square-root term of
\eqref{eq:wc_signal_mask} by $H_n\sqrt r/s_\star$ gives
$C H_n\kappa r\sqrt{\mu_0\log m/(np)}$.
The second term is smaller under the same condition, because $m\asymp n$
and $H_n,\mu_0,\kappa,r\ge1$.
Weyl's inequality gives
\begin{align}
 s_\star-b_{\rm sig}-CH_nh
 \le\widehat s_r
 \le s_\star+b_{\rm sig}+CH_nh.
 \label{eq:wc_observed_sr}
\end{align}
Importantly, this does not assume that $H_nh$ is small relative to
$s_\star$.

If the direct branch is selected, combine its comparison with the upper
bound in \eqref{eq:wc_raw_scale} and the lower bound in
\eqref{eq:wc_observed_sr}. This gives
\[
 C\epsilon_H+CH_n\sqrt r\,b_{\rm sig}
 \ge t_0(s_\star-b_{\rm sig}-CH_nh).
\]
Choosing $c_b$ sufficiently small relative to $t_0$ and using
$H_nh\le\epsilon_H$ yields
\begin{align}
 \epsilon_H\ge c_Dt_0s_\star,
 \qquad b_{\rm sig}\le c_D' h,
 \label{eq:wc_direct_comparison}
\end{align}
where $c_D'$ can be made arbitrarily small by reducing $c_b$.
Consequently \eqref{eq:wc_raw_scale} implies $h_0\asymp h$.
This proves the claimed scale bound in the direct branch using
$\widehat s_a=s_h(0)$.

If the iterative branch is selected, use the lower bound in
\eqref{eq:wc_raw_scale} and the upper bound in
\eqref{eq:wc_observed_sr} to obtain
\[
 c\epsilon_H-CH_n\sqrt r\,b_{\rm sig}
 <t_0(s_\star+b_{\rm sig}+CH_nh).
\]
Take $t_0$ small enough to absorb $Ct_0H_nh$ into the left side, and then
$c_b$ sufficiently small relative to $t_0$. We get
\begin{align}
 \epsilon_H\le C_It_0s_\star.
 \label{eq:wc_iterative_comparison}
\end{align}
Thus the required small-noise local regime is selected by an observable
test; it is not added as a global signal-to-noise assumption.

\textit{2. Accuracy of the direct projected spectral anchor.}
Rank approximation gives
\begin{align}
 \|\operatorname{Proj}_r(Z)-M_0\|_F
 \le 2\sqrt{2r}\,\|Z-M_0\|
 \le C\sqrt r(b_{\rm sig}+H_nh).
 \label{eq:wc_spectral_accuracy}
\end{align}
Consider first $H_nh\le c_1s_1$, for a sufficiently small fixed $c_1$.
Then \eqref{eq:wc_mask_small} and Weyl imply
$\widehat s_1\asymp s_1$. The row and column norms of the true matrix
are at most $s_1\sqrt{\mu_0r/m}$ and
$s_1\sqrt{\mu_0r/n}$, respectively. The projection estimate
\eqref{eq:recovery_projectionF}, with its radius constant adjusted by
these fixed comparison factors, preserves
\eqref{eq:wc_spectral_accuracy} up to a fixed multiplier. In the direct
branch \eqref{eq:wc_direct_comparison} then gives
$\|M_b-M_0\|_F\le C\epsilon_H$.

If $H_nh>c_1s_1$, no comparison between the observed and true leading
singular values is necessary. The diagonal factors in $P_B$ are
contractions, so
\begin{align*}
 \|M_b-M_0\|_F
 &\le\|\operatorname{Proj}_r(Z)\|_F+\|M_0\|_F\\
 &\le\sqrt r\{\|Z\|+s_1\}
 \le C\sqrt r(s_1+b_{\rm sig}+H_nh)
 \le C\epsilon_H.
\end{align*}
This proves the direct-branch Frobenius claim even at arbitrarily large
noise levels. In particular, no incoherence or sampled-error bound for a
rank-deficient direct anchor was used to estimate its scale.

\textit{3. Local start and the unweighted recovery recursion.}
In the iterative branch, \eqref{eq:wc_iterative_comparison} and
\eqref{eq:wc_mask_small} imply $\widehat s_1\asymp s_1$ and
\[
 f_0=\|M_b-M_0\|_F\le
 C\sqrt r(b_{\rm sig}+H_nh)\le c_0s_\star,
\]
with $c_0$ as small as needed after choosing $t_0,c_b$.
Within this neighborhood, Lemma~\ref{lem:recovery_sampling} applies to
the row--column envelope imposed by $P_B$. Its parameter is bounded by
\begin{align}
 a_s\le C\left\{
 \sqrt{\frac{\mu_0r\log m}{np}}+
 \frac{K_0^2\kappa^2\mu_0r}{\sqrt{np}}\right\}.
 \label{eq:wc_sampling_constant}
\end{align}
The sample-size condition makes $a_s+\sqrt{r/(np)}$ a sufficiently small
fixed constant.

At an unweighted recovery iterate the score decomposes as
\begin{align}
 G_h^{\rm rec}(M)&=p^{-1}P_\Omega(M-M_0)+\Xi(M),\nonumber\\
 \|\Xi(M)\|&\le CH_nh+C\|M-M_0\|_F/\sqrt{np}.
 \label{eq:wc_score_decomposition}
\end{align}
The first part of the remainder is $-p^{-1}P_\Omega H$.
The second is the operator-capped reference term; its norm is at most
$C_Qs_h(M)\sqrt{m/p}$. By \eqref{eq:wc_scale_uniform} and the local sampled
energy bound, this is at most $Ch+Cf/\sqrt{np}$.
Together with \eqref{eq:wc_noise_edge} these prove
\eqref{eq:wc_score_decomposition}.

For the remainder $W$ in Lemma~\ref{lem:rsm_inexact_scaled_update}, use
\eqref{eq:wc_sampling_constant} and \eqref{eq:wc_score_decomposition} to
bound its projected norms by
\begin{align*}
 w_o&\le a_sf+CH_nh+Cf/\sqrt{np},\\
 w_f&\le a_sf+C\epsilon_H+C\sqrt{r/(np)}f.
\end{align*}
The Frobenius one-step bound \eqref{eq:recovery_f-full}, with $\gamma=0$
and $\eta=1/2$, gives the explicit pre-projection inequality
\begin{align*}
 f_{\rm pre}\le
 \left[\frac12+C\left\{c_0+t_0+a_s+\sqrt{\frac r{np}}\right\}\right]f
 +C\epsilon_H.
\end{align*}
Indeed $f\le c_0s_\star$, $\epsilon_H\le C_It_0s_\star$, and the
projected bounds imply $(2f+w_o)/(s_\star-f)\le C(c_0+t_0+a_s)$.
Thus the quadratic term in \eqref{eq:recovery_f-full} is bounded by this
small multiplier times $w_f$, giving the displayed inequality. Choose $K_0$ sufficiently large
that the Frobenius multiplier in \eqref{eq:recovery_projectionF} is at most
$9/8$, allowing for the fixed observed-radius comparison. Choosing
$t_0,c_b,a_s$ smaller then makes the post-projection factor at most $3/4$.
This proves \eqref{eq:wc_rec}. Choose the initial radius to be at most
half of the radius of the local neighborhood and
$C_R\epsilon_H/(1-\rho_a)$ at most its other half; these are consequences
of the choices of $t_0,c_b$, not additional random assumptions. Induction
then preserves that neighborhood, and Weyl's inequality preserves rank $r$. Summing the geometric
series proves \eqref{eq:wc_geometric}.

\textit{4. The observable stop, its converse, and final scale.}
The uniform scale upper bound and local sampled-energy comparison give
\[
 s_h(M^{(k)})\sqrt{m/p}\le Ch+Cf_k/\sqrt{np}.
\]
At a stopping update, the triangle inequality and \eqref{eq:wc_rec} imply
\begin{align*}
 (1-\rho_a)f_k
 &\le\|M^{(k+1)}-M^{(k)}\|_F+C_R\epsilon_H\\
 &\le C\epsilon_H+
 C C_{\rm stop}H_n\sqrt{r/(np)}f_k.
\end{align*}
The last term is absorbed by $np\ge C H_n^2r$, with its constant allowed
to depend on the fixed $C_{\rm stop}$. Hence both the previous and the
returned errors are at most $C\epsilon_H$.
Conversely, \eqref{eq:wc_geometric} enters this error scale within
$C\{1+\log_+(f_0/\epsilon_H)\}$ steps. There the uniform scale bound,
local sampled-energy comparison, and \eqref{eq:wc_bulk_scale} give
\[
 c_0\sigma-C H_n\sigma\sqrt{r/(np)}
 \le s_h(M^{(k)})
 \le C_0\sigma+C H_n\sigma\sqrt{r/(np)}.
\]
The lower bound is a positive multiple of $\sigma$ under the sample-size
condition. One more update bounds
$\|M^{(k+1)}-M^{(k)}\|_F\le C\epsilon_H$ and triggers the stop when
$C_{\rm stop}$ is sufficiently large. The same comparison at the returned
matrix proves $\widehat s_a=s_h(M_a)\asymp\sigma$.
All events were uniform. No union bound over the random number of
iterations, no fitted-weight independence, and no lower bound on
$\sigma$ were introduced.
\end{proof}

\subsection{Proof of the distribution-unknown output theorem}
\begin{proof}[Proof of Theorem~\ref{main-thm:weighted_unknown_rate}]
Lemma~\ref{lem:wc_anchor} proves the initialization and scale assertions,
and its recurrence and stopping complexity.
On its common event, let $\widetilde M$ be any candidate, without a rank,
stationarity, independence, or objective-comparison requirement. If it
passes the output rule, then
\begin{align*}
 \|\mathfrak C_a(\widetilde M)-M_0\|_F
 &\le\|\widetilde M-M_a\|_F+\|M_a-M_0\|_F\\
 &\le A_0H_n\widehat s_a\sqrt{mr/p}+C_a\epsilon_H\\
 &\le(A_0C_a'+C_a)\epsilon_H.
\end{align*}
If it fails or is nonfinite, the output is $M_a$, whose error is already
bounded. This is a deterministic statement for all candidates on the
same event, so it covers data-dependent weights, repeated tuning,
arbitrary weighted recovery budgets, and nuclear-norm refinement.
The candidate may use either the estimated profile or the true profile.
Both returned matrices have the same $\epsilon_H$ order. In fact their
mutual Frobenius distance is at most $2R_a$, but no smaller-order
oracle equivalence follows from this bound.

If $\widetilde M_{\rm raw}$ denotes an unguarded result, a useful precise
corollary is
\begin{align}
 \mathbb P\bigl(\|\widetilde M_{\rm raw}-M_0\|_F>C\epsilon_H,
      \ \|\widetilde M_{\rm raw}-M_a\|_F\le R_a\bigr)
 \le cn^{-3}.
 \label{eq:wc_raw_joint_certificate}
\end{align}
This is a joint event bound. It must not be described as a conditional
coverage statement given acceptance without accounting for the
acceptance probability. Nor does it bound the error of rejected raw
outputs. The bound has order $H_n\sigma\sqrt{mr/p}$ for both estimated and true-profile weights under the stated sampling conditions.
\end{proof}

\subsection{A standard nonparametric tail class}
\begin{scorollary}[Unknown sub-exponential multiplier laws]
\label{cor:wc_standard_tail_class}
Retain the positivity, mean-one, lower-mass, independence, and
standardized-error conditions of Assumption~\ref{main-assum:hetero}.
Suppose the two multiplier laws are otherwise arbitrary and satisfy
\[
 \mathbb P(\nu>x)\vee\mathbb P(\mu>x)\le2e^{-x/K_V},\qquad x>0,
\]
for a fixed tail-class constant $K_V$. Write $L=\log\{e(m+n)\}$.
For any fixed $\alpha>52$, sufficient sampling conditions for
Theorem~\ref{main-thm:weighted_unknown_rate} are
\begin{align*}
 mp&\ge C L^{\alpha+1},\\
 np&\ge C\{\mu_0\kappa^2r^2L^2+
                  \mu_0^2\kappa^4r^2+rL\}.
\end{align*}
Under these conditions, the stabilized estimated-weight and true-profile
versions both satisfy
\[
 \|\widehat M-M_0\|_F\le C\sigma\sqrt{mrL/p}
\]
with probability at least $1-cn^{-3}$. The constants do not depend on the
specific multiplier laws within this class.
\end{scorollary}
\begin{proof}
The tail bound implies a uniform twelfth-moment bound by tail
integration. Choose $U_\nu=U_\mu=C_VL$, where $C_V$ is sufficiently large
relative to $K_V$ and the fixed exceptional-probability constant. Then
$\mathbb P(\nu>U_\nu)\vee\mathbb P(\mu>U_\mu)\le c_1n^{-5}$.
For sufficiently large dimensions, $U_\nu+U_\mu\le\log^2m$ because the
aspect ratio is bounded away from zero. Thus one may take $C_U=2$,
$H_n^2=2C_VL$, and $\alpha_{\rm vp}=\alpha>52$ in the main theorem.
Substitution proves the sampling conditions and error bound.
\end{proof}
The logarithm in this corollary is the common tail-class envelope for
unbounded multipliers. It is not an extra error term due to estimating
the weights. No Gaussian or other parametric law for those multipliers
is specified. The exponents are sufficient and are not asserted optimal.

\subsection{Shape normalization and same-data weights}
\begin{slemma}[Scale invariance and pathwise weight stability]
\label{lem:wc_weight_stability}
For a positive matrix $S$, define $T_v(S)=S\vee Q_v(S)$ entrywise and
$\mathcal W_v(S)=\operatorname{median}\{T_v(S)\}/T_v(S)$.
If the quantile and median conventions are monotone and positively
homogeneous, then, for every $c>0$,
\begin{align}
 \mathcal W_v(cS)=\mathcal W_v(S).
 \label{eq:wc_exact_shape_invariance}
\end{align}
If, for some $c>0$,
\begin{align*}
 \max_{ij}\bigl|\log T_v(\widehat S)_{ij}
                    -\log\{cT_v(S^*)_{ij}\}\bigr|\le\varepsilon,
\end{align*}
then
\begin{align}
 \max_{ij}\left|\log\frac{\widehat W_{ij}}{W^*_{ij}}\right|
 &\le2\varepsilon,\label{eq:wc_weight_log_error}\\
 e^{-4\varepsilon}\|W^*\odot P_\Omega\Delta\|_F^2
 \le\|\widehat W\odot P_\Omega\Delta\|_F^2
 &\le e^{4\varepsilon}\|W^*\odot P_\Omega\Delta\|_F^2
 \quad\text{for every }\Delta.
 \label{eq:wc_weight_energy_transfer}
\end{align}
Both statements permit weights fitted from the same observations.
\end{slemma}
\begin{proof}
Positive homogeneity gives $T_v(cS)=cT_v(S)$ and its median also multiplies
by $c$, proving \eqref{eq:wc_exact_shape_invariance}. The log-error premise
puts $T_v(\widehat S)$ between $ce^{-\varepsilon}T_v(S^*)$ and
$ce^{\varepsilon}T_v(S^*)$ entrywise. Monotonicity gives the same bounds
for their medians. Dividing yields relative weights between
$e^{-2\varepsilon}$ and $e^{2\varepsilon}$. Squaring and summing the
nonnegative entries of the weighted sampled energy proves the result.
\end{proof}

The normalization removes any common scale constant from the median-based weights. Lemma~\ref{lem:wc_weight_stability} is a pathwise stability statement. The stabilized error bound does not require consistency of the fitted weights.

For the two median passes in \eqref{eq:weighted_median_profile},
let $c_j^0=\operatorname{median}_i\ell_{ij}$ and
$a_i^0=\operatorname{median}_j(\ell_{ij}-c_j^0)$.
Subtracting a common centering constant $b$ from all column effects adds
$b$ to every subsequent raw row effect. Subtracting a common constant
from all row effects is compensated by the final overall median.
Consequently the repeated recentering formula in the implementation
summary gives the same fitted shape as these two passes. This identity applies to the two-pass fit in \eqref{eq:weighted_median_profile}.

\subsection{Weighted derivatives and recovery proposals}
The linear observation operator associated with fixed $W$ is
$\mathcal A_W(M)=W\odot P_\Omega M$, with adjoint
$\mathcal A_W^*(B)=P_\Omega(W\odot B)$. The singular-value differential
therefore proves \eqref{eq:weighted_fixed_gradient}. If the scale is
instead recomputed at every candidate, the full derivative has the
additional correction
\begin{align*}
 -P_\Omega\!\left[W\odot U_W\operatorname{diag}\left\{
 \frac{2}{pb_W}\left(\omega_j e_j-
 \frac{\mathbf1_{\{j\in I\}}}{S_W}
 \sum_i\omega_i e_i d_i^W\right)\right\}V_W'\right],
\end{align*}
where $S_W=\sum_{i\in I}d_i^W$, $b_W=S_W/\sum_{i\in I}\widehat\lambda_i$,
and $e_i=d_i^W/b_W-\widehat\lambda_i$.
Neither formula should be substituted for the recovery score without
changing the numerical algorithm.

For the pre-outer-weight cap defined above,
$E_{cW}=cE_W$, $Q_{cW}=cQ_W$, and $\alpha_{cW}=\alpha_W$.
Thus $G_{cW}^{\rm rec}=c^2G_W^{\rm rec}$. The uncapped part of the auto
step only scales as $c^{-1}$. This explains why invariance of a profiled
spectral objective alone does not imply invariance of the score
trajectory. Normalization selects a unique scale representative before
this trajectory is run. The scaling calculation uses the cap on $P_\Omega Q_W$ specified above.

With $\Delta=M-M_0$, the weighted score has the exact decomposition
\begin{align}
 G_W^{\rm rec}(M)
 =p^{-1}P_\Omega(W^2\odot\Delta)
  -p^{-1}P_\Omega(W^2\odot H)
  +p^{-1}\alpha_W W\odot P_\Omega Q_W(M).
 \label{eq:wc_weighted_score_decomp}
\end{align}
The first operator is not the identity in expectation, even for known
$W$. Moreover, general entrywise multiplication does not obey
$\|W\odot A\|\le\|W\|_\infty\|A\|$ in spectral norm. The recovery guarantee follows from the calibration recurrence and the observable output rule.

For clarity, with an exact separable scale $\tau_{ij}$ and an oracle
floor $c_v$, set $s_v=\operatorname{median}(\tau\vee c_v)$. Its normalized
weight gives conditional variances
\begin{align}
 \operatorname{Var}(W_{ij}^*H_{ij}\mid\tau)
 =\sigma^2s_v^2\left\{\frac{\tau_{ij}}{\tau_{ij}\vee c_v}\right\}^2.
 \label{eq:wc_residual_variance_after_floor}
\end{align}
The floor is thus a regularization of whitening, not an exact
homoskedastic transformation. It can also destroy separability of the
weight matrix. The output theorem accommodates the resulting nonconstant conditional variances.
\end{revision}
% END MAJOR REVISION S08


\section{Detailed numerical results}
\label{sec:supp_numerical}
% BEGIN MAJOR REVISION S09
\begin{revision}
The simulations use the settings and relative error definitions of Section~\ref{main-sec:numerical}. Unless specified otherwise, each table entry is the mean of 100 repetitions. The factor and nuclear-norm implementations are denoted by $\widetilde{\M}_{\mathrm{fac}}$ and $\widetilde{\M}_{\mathrm{nuc}}$, and their least-squares baselines by $\widehat{\M}_{\mathrm{fac}}$ and $\widehat{\M}_{\mathrm{nuc}}$. The perturbed-truth initialization, validation stopping rule and oracle selection of the proposed nuclear penalty are specified in Section~\ref{main-sec:simulation_implementation}. The tabulated comparisons are conditional on these choices.
All errors are displayed to four decimal places, so $0.0000$ denotes rounding rather than a claim of exact recovery. The factor and nuclear implementations have finite reported errors in all configurations included here.

\end{revision}
% END MAJOR REVISION S09

\subsection{Matrix dimensions}
Tables~\ref{tab:change_matrix_size_sp} and \ref{tab:change_matrix_size_max} report the spectral and maximum errors for the dimension experiment. Figure~\ref{fig:change_matrix_size_spandmax} plots the mean errors for the same configurations. The errors decrease with matrix dimension, and the two proposed implementations have similar accuracy.
\begin{table}[!htbp]
\centering
\caption{Mean relative spectral errors $\mathcal E_{sp}$ for $\sigma=1$, $p=0.2$ and $n/m=0.5$. Each entry averages 100 repetitions.}
\label{tab:change_matrix_size_sp}
\begingroup
\small
\setstretch{1.05}
\setlength{\tabcolsep}{3.2pt}
\resizebox{\linewidth}{!}{%
\begin{tabular}{@{}llrrrrrrrrr@{}}
\toprule
$r$ & Estimator & $m=200$ & $m=300$ & $m=400$ & $m=500$ & $m=600$ & $m=700$ & $m=800$ & $m=900$ & $m=1000$ \\
\midrule
\multirow{4}{*}{$5$} & $\widetilde{\M}_{\mathrm{fac}}$ & 0.0899 & 0.0675 & 0.0560 & 0.0487 & 0.0438 & 0.0398 & 0.0369 & 0.0346 & 0.0324 \\
 & $\widehat{\M}_{\mathrm{fac}}$ & 0.1019 & 0.0739 & 0.0599 & 0.0525 & 0.0468 & 0.0420 & 0.0386 & 0.0358 & 0.0334 \\
 & $\widetilde{\M}_{\mathrm{nuc}}$ & 0.0895 & 0.0673 & 0.0558 & 0.0486 & 0.0437 & 0.0397 & 0.0368 & 0.0346 & 0.0324 \\
 & $\widehat{\M}_{\mathrm{nuc}}$ & 0.0973 & 0.0735 & 0.0596 & 0.0521 & 0.0471 & 0.0429 & 0.0397 & 0.0373 & 0.0351 \\
\midrule
\multirow{4}{*}{$10$} & $\widetilde{\M}_{\mathrm{fac}}$ & 0.1047 & 0.0743 & 0.0596 & 0.0508 & 0.0452 & 0.0408 & 0.0374 & 0.0348 & 0.0327 \\
 & $\widehat{\M}_{\mathrm{fac}}$ & 0.1323 & 0.0838 & 0.0653 & 0.0545 & 0.0479 & 0.0436 & 0.0396 & 0.0365 & 0.0341 \\
 & $\widetilde{\M}_{\mathrm{nuc}}$ & 0.1043 & 0.0742 & 0.0594 & 0.0508 & 0.0451 & 0.0407 & 0.0374 & 0.0347 & 0.0327 \\
 & $\widehat{\M}_{\mathrm{nuc}}$ & 0.1199 & 0.0819 & 0.0646 & 0.0553 & 0.0493 & 0.0435 & 0.0400 & 0.0373 & 0.0351 \\
\midrule
\multirow{4}{*}{$20$} & $\widetilde{\M}_{\mathrm{fac}}$ & 0.1378 & 0.0933 & 0.0711 & 0.0590 & 0.0508 & 0.0452 & 0.0410 & 0.0376 & 0.0350 \\
 & $\widehat{\M}_{\mathrm{fac}}$ & 0.1512 & 0.1137 & 0.0838 & 0.0658 & 0.0556 & 0.0488 & 0.0436 & 0.0397 & 0.0367 \\
 & $\widetilde{\M}_{\mathrm{nuc}}$ & 0.1367 & 0.0932 & 0.0712 & 0.0589 & 0.0508 & 0.0452 & 0.0409 & 0.0376 & 0.0350 \\
 & $\widehat{\M}_{\mathrm{nuc}}$ & 0.1430 & 0.1064 & 0.0820 & 0.0667 & 0.0563 & 0.0503 & 0.0445 & 0.0410 & 0.0382 \\
\bottomrule
\end{tabular}
}
\endgroup
\end{table}
\begin{table}[!htbp]
\centering
\caption{Mean relative maximum errors $\mathcal E_\infty$ for $\sigma=1$, $p=0.2$ and $n/m=0.5$. Each entry averages 100 repetitions.}
\label{tab:change_matrix_size_max}
\begingroup
\small
\setstretch{1.05}
\setlength{\tabcolsep}{3.2pt}
\resizebox{\linewidth}{!}{%
\begin{tabular}{@{}llrrrrrrrrr@{}}
\toprule
$r$ & Estimator & $m=200$ & $m=300$ & $m=400$ & $m=500$ & $m=600$ & $m=700$ & $m=800$ & $m=900$ & $m=1000$ \\
\midrule
\multirow{4}{*}{$5$} & $\widetilde{\M}_{\mathrm{fac}}$ & 0.1220 & 0.0889 & 0.0741 & 0.0640 & 0.0575 & 0.0522 & 0.0485 & 0.0449 & 0.0432 \\
 & $\widehat{\M}_{\mathrm{fac}}$ & 0.1433 & 0.0988 & 0.0800 & 0.0687 & 0.0609 & 0.0549 & 0.0498 & 0.0461 & 0.0437 \\
 & $\widetilde{\M}_{\mathrm{nuc}}$ & 0.1211 & 0.0886 & 0.0740 & 0.0640 & 0.0575 & 0.0521 & 0.0484 & 0.0449 & 0.0432 \\
 & $\widehat{\M}_{\mathrm{nuc}}$ & 0.1333 & 0.0984 & 0.0799 & 0.0679 & 0.0616 & 0.0567 & 0.0519 & 0.0493 & 0.0462 \\
\midrule
\multirow{4}{*}{$10$} & $\widetilde{\M}_{\mathrm{fac}}$ & 0.1424 & 0.1007 & 0.0804 & 0.0685 & 0.0605 & 0.0546 & 0.0499 & 0.0471 & 0.0445 \\
 & $\widehat{\M}_{\mathrm{fac}}$ & 0.1840 & 0.1143 & 0.0876 & 0.0724 & 0.0636 & 0.0585 & 0.0520 & 0.0490 & 0.0459 \\
 & $\widetilde{\M}_{\mathrm{nuc}}$ & 0.1413 & 0.1002 & 0.0805 & 0.0686 & 0.0605 & 0.0546 & 0.0498 & 0.0471 & 0.0445 \\
 & $\widehat{\M}_{\mathrm{nuc}}$ & 0.1636 & 0.1107 & 0.0859 & 0.0740 & 0.0671 & 0.0584 & 0.0526 & 0.0507 & 0.0476 \\
\midrule
\multirow{4}{*}{$20$} & $\widetilde{\M}_{\mathrm{fac}}$ & 0.1778 & 0.1268 & 0.0971 & 0.0807 & 0.0688 & 0.0621 & 0.0545 & 0.0516 & 0.0477 \\
 & $\widehat{\M}_{\mathrm{fac}}$ & 0.1910 & 0.1535 & 0.1148 & 0.0910 & 0.0745 & 0.0669 & 0.0580 & 0.0545 & 0.0499 \\
 & $\widetilde{\M}_{\mathrm{nuc}}$ & 0.1748 & 0.1248 & 0.0968 & 0.0806 & 0.0688 & 0.0621 & 0.0545 & 0.0516 & 0.0476 \\
 & $\widehat{\M}_{\mathrm{nuc}}$ & 0.1824 & 0.1429 & 0.1122 & 0.0924 & 0.0758 & 0.0696 & 0.0596 & 0.0570 & 0.0528 \\
\bottomrule
\end{tabular}
}
\endgroup
\end{table}
\begin{figure}[!htbp]
\centering
\subfigure{\includegraphics[width=0.31\linewidth]{plot/l2_loss_5.png}}
\subfigure{\includegraphics[width=0.31\linewidth]{plot/l2_loss_10.png}}
\subfigure{\includegraphics[width=0.31\linewidth]{plot/l2_loss_20.png}}
\par\medskip
\subfigure{\includegraphics[width=0.31\linewidth]{plot/spectral_loss_5.png}}
\subfigure{\includegraphics[width=0.31\linewidth]{plot/spectral_loss_10.png}}
\subfigure{\includegraphics[width=0.31\linewidth]{plot/spectral_loss_20.png}}
\par\medskip
\subfigure{\includegraphics[width=0.31\linewidth]{plot/maximum_loss_5.png}}
\subfigure{\includegraphics[width=0.31\linewidth]{plot/maximum_loss_10.png}}
\subfigure{\includegraphics[width=0.31\linewidth]{plot/maximum_loss_20.png}}
\caption{Mean relative Frobenius, spectral and maximum errors as a function of matrix dimension, shown in the first, second and third rows, respectively. The columns correspond to $r=5,10,20$. The observation probability is $p=0.2$, the noise level is $\sigma=1$ and $n/m=0.5$. Each point averages 100 repetitions.}
\label{fig:change_matrix_size_spandmax}
\end{figure}

% BEGIN MAJOR REVISION S10
\begin{revision}
\subsection{Homogeneous noise levels}
We fix $m=500$, $n=250$ and $p=0.2$, and vary $\sigma$ over $\{0,0.01,0.02,0.05,0.1,0.2,0.5,1,2,5\}$ for $r=5,10,20$. Tables~\ref{tab:supp:different_noise_fnorm_addition}, \ref{tab:supp:different_noise_specnorm_addition} and \ref{tab:supp:different_noise_maxnorm_addition} report the three errors. The advantage over least squares is most pronounced at moderate and high noise. At very small noise and $r=20$, some recovery runs reach the 1500-update cap and retain a visible numerical error. These finite-budget results do not establish exact noiseless recovery of the numerical procedure.

\end{revision}
% END MAJOR REVISION S10
\begin{table}[!htbp]
\centering
\caption{Mean relative Frobenius errors $\mathcal E_F$ under homogeneous noise for $m=500$, $n=250$ and $p=0.2$. Each entry averages 100 repetitions.}
\label{tab:supp:different_noise_fnorm_addition}
\begingroup
\small
\setstretch{1.05}
\setlength{\tabcolsep}{3.2pt}
\resizebox{\linewidth}{!}{%
\begin{tabular}{@{}llrrrrrrrrrr@{}}
\toprule
$r$ & Estimator & $\sigma=0$ & $\sigma=0.01$ & $\sigma=0.02$ & $\sigma=0.05$ & $\sigma=0.1$ & $\sigma=0.2$ & $\sigma=0.5$ & $\sigma=1$ & $\sigma=2$ & $\sigma=5$ \\
\midrule
\multirow{4}{*}{$5$} & $\widetilde{\M}_{\mathrm{fac}}$ & 0.0000 & 0.0008 & 0.0015 & 0.0038 & 0.0077 & 0.0154 & 0.0377 & 0.0692 & 0.1089 & 0.1470 \\
 & $\widehat{\M}_{\mathrm{fac}}$ & 0.0000 & 0.0008 & 0.0016 & 0.0040 & 0.0078 & 0.0154 & 0.0377 & 0.0745 & 0.1445 & 0.3714 \\
 & $\widetilde{\M}_{\mathrm{nuc}}$ & 0.0000 & 0.0008 & 0.0016 & 0.0039 & 0.0078 & 0.0155 & 0.0377 & 0.0691 & 0.1087 & 0.1466 \\
 & $\widehat{\M}_{\mathrm{nuc}}$ & 0.0000 & 0.0008 & 0.0016 & 0.0039 & 0.0077 & 0.0154 & 0.0379 & 0.0741 & 0.1483 & 0.3809 \\
\midrule
\multirow{4}{*}{$10$} & $\widetilde{\M}_{\mathrm{fac}}$ & 0.0001 & 0.0008 & 0.0016 & 0.0040 & 0.0081 & 0.0161 & 0.0391 & 0.0703 & 0.1069 & 0.1412 \\
 & $\widehat{\M}_{\mathrm{fac}}$ & 0.0000 & 0.0009 & 0.0017 & 0.0042 & 0.0084 & 0.0164 & 0.0392 & 0.0754 & 0.1498 & 0.3844 \\
 & $\widetilde{\M}_{\mathrm{nuc}}$ & 0.0001 & 0.0008 & 0.0016 & 0.0041 & 0.0081 & 0.0162 & 0.0392 & 0.0702 & 0.1067 & 0.1407 \\
 & $\widehat{\M}_{\mathrm{nuc}}$ & 0.0000 & 0.0009 & 0.0017 & 0.0042 & 0.0082 & 0.0163 & 0.0393 & 0.0766 & 0.1470 & 0.3739 \\
\midrule
\multirow{4}{*}{$20$} & $\widetilde{\M}_{\mathrm{fac}}$ & 0.0042 & 0.0043 & 0.0046 & 0.0064 & 0.0106 & 0.0200 & 0.0462 & 0.0766 & 0.1085 & 0.1381 \\
 & $\widehat{\M}_{\mathrm{fac}}$ & 0.0000 & 0.0014 & 0.0027 & 0.0061 & 0.0113 & 0.0211 & 0.0465 & 0.0853 & 0.1613 & 0.4182 \\
 & $\widetilde{\M}_{\mathrm{nuc}}$ & 0.0039 & 0.0041 & 0.0045 & 0.0064 & 0.0107 & 0.0201 & 0.0465 & 0.0765 & 0.1083 & 0.1376 \\
 & $\widehat{\M}_{\mathrm{nuc}}$ & 0.0000 & 0.0014 & 0.0026 & 0.0060 & 0.0112 & 0.0208 & 0.0464 & 0.0868 & 0.1654 & 0.4198 \\
\bottomrule
\end{tabular}
}
\endgroup
\end{table}
\begin{table}[!htbp]
\centering
\caption{Mean relative spectral errors $\mathcal E_{sp}$ under homogeneous noise for $m=500$, $n=250$ and $p=0.2$. Each entry averages 100 repetitions.}
\label{tab:supp:different_noise_specnorm_addition}
\begingroup
\small
\setstretch{1.05}
\setlength{\tabcolsep}{3.2pt}
\resizebox{\linewidth}{!}{%
\begin{tabular}{@{}llrrrrrrrrrr@{}}
\toprule
$r$ & Estimator & $\sigma=0$ & $\sigma=0.01$ & $\sigma=0.02$ & $\sigma=0.05$ & $\sigma=0.1$ & $\sigma=0.2$ & $\sigma=0.5$ & $\sigma=1$ & $\sigma=2$ & $\sigma=5$ \\
\midrule
\multirow{4}{*}{$5$} & $\widetilde{\M}_{\mathrm{fac}}$ & 0.0000 & 0.0005 & 0.0011 & 0.0027 & 0.0054 & 0.0108 & 0.0265 & 0.0487 & 0.0765 & 0.1035 \\
 & $\widehat{\M}_{\mathrm{fac}}$ & 0.0000 & 0.0006 & 0.0012 & 0.0029 & 0.0056 & 0.0109 & 0.0266 & 0.0525 & 0.1023 & 0.2679 \\
 & $\widetilde{\M}_{\mathrm{nuc}}$ & 0.0000 & 0.0006 & 0.0011 & 0.0027 & 0.0055 & 0.0109 & 0.0265 & 0.0486 & 0.0764 & 0.1030 \\
 & $\widehat{\M}_{\mathrm{nuc}}$ & 0.0000 & 0.0006 & 0.0011 & 0.0028 & 0.0055 & 0.0108 & 0.0267 & 0.0521 & 0.1046 & 0.2747 \\
\midrule
\multirow{4}{*}{$10$} & $\widetilde{\M}_{\mathrm{fac}}$ & 0.0002 & 0.0006 & 0.0012 & 0.0029 & 0.0058 & 0.0117 & 0.0283 & 0.0508 & 0.0774 & 0.1023 \\
 & $\widehat{\M}_{\mathrm{fac}}$ & 0.0000 & 0.0007 & 0.0014 & 0.0032 & 0.0062 & 0.0120 & 0.0284 & 0.0545 & 0.1088 & 0.2844 \\
 & $\widetilde{\M}_{\mathrm{nuc}}$ & 0.0002 & 0.0006 & 0.0012 & 0.0030 & 0.0059 & 0.0117 & 0.0283 & 0.0508 & 0.0772 & 0.1020 \\
 & $\widehat{\M}_{\mathrm{nuc}}$ & 0.0000 & 0.0007 & 0.0013 & 0.0031 & 0.0060 & 0.0119 & 0.0285 & 0.0553 & 0.1063 & 0.2754 \\
\midrule
\multirow{4}{*}{$20$} & $\widetilde{\M}_{\mathrm{fac}}$ & 0.0046 & 0.0047 & 0.0048 & 0.0058 & 0.0086 & 0.0157 & 0.0358 & 0.0590 & 0.0833 & 0.1055 \\
 & $\widehat{\M}_{\mathrm{fac}}$ & 0.0000 & 0.0019 & 0.0029 & 0.0055 & 0.0094 & 0.0167 & 0.0360 & 0.0658 & 0.1249 & 0.3301 \\
 & $\widetilde{\M}_{\mathrm{nuc}}$ & 0.0043 & 0.0045 & 0.0047 & 0.0058 & 0.0087 & 0.0158 & 0.0362 & 0.0589 & 0.0831 & 0.1053 \\
 & $\widehat{\M}_{\mathrm{nuc}}$ & 0.0000 & 0.0018 & 0.0028 & 0.0053 & 0.0092 & 0.0164 & 0.0359 & 0.0667 & 0.1271 & 0.3275 \\
\bottomrule
\end{tabular}
}
\endgroup
\end{table}
\begin{table}[!htbp]
\centering
\caption{Mean relative maximum errors $\mathcal E_\infty$ under homogeneous noise for $m=500$, $n=250$ and $p=0.2$. Each entry averages 100 repetitions.}
\label{tab:supp:different_noise_maxnorm_addition}
\begingroup
\small
\setstretch{1.05}
\setlength{\tabcolsep}{3.2pt}
\resizebox{\linewidth}{!}{%
\begin{tabular}{@{}llrrrrrrrrrr@{}}
\toprule
$r$ & Estimator & $\sigma=0$ & $\sigma=0.01$ & $\sigma=0.02$ & $\sigma=0.05$ & $\sigma=0.1$ & $\sigma=0.2$ & $\sigma=0.5$ & $\sigma=1$ & $\sigma=2$ & $\sigma=5$ \\
\midrule
\multirow{4}{*}{$5$} & $\widetilde{\M}_{\mathrm{fac}}$ & 0.0001 & 0.0007 & 0.0015 & 0.0037 & 0.0074 & 0.0147 & 0.0358 & 0.0640 & 0.0986 & 0.1297 \\
 & $\widehat{\M}_{\mathrm{fac}}$ & 0.0000 & 0.0013 & 0.0024 & 0.0050 & 0.0085 & 0.0154 & 0.0359 & 0.0687 & 0.1310 & 0.3397 \\
 & $\widetilde{\M}_{\mathrm{nuc}}$ & 0.0001 & 0.0008 & 0.0015 & 0.0039 & 0.0077 & 0.0149 & 0.0357 & 0.0640 & 0.0988 & 0.1290 \\
 & $\widehat{\M}_{\mathrm{nuc}}$ & 0.0000 & 0.0012 & 0.0022 & 0.0044 & 0.0080 & 0.0151 & 0.0356 & 0.0679 & 0.1356 & 0.3512 \\
\midrule
\multirow{4}{*}{$10$} & $\widetilde{\M}_{\mathrm{fac}}$ & 0.0009 & 0.0013 & 0.0019 & 0.0042 & 0.0084 & 0.0166 & 0.0395 & 0.0685 & 0.0990 & 0.1272 \\
 & $\widehat{\M}_{\mathrm{fac}}$ & 0.0000 & 0.0021 & 0.0034 & 0.0064 & 0.0107 & 0.0185 & 0.0395 & 0.0724 & 0.1416 & 0.3633 \\
 & $\widetilde{\M}_{\mathrm{nuc}}$ & 0.0009 & 0.0013 & 0.0019 & 0.0043 & 0.0084 & 0.0167 & 0.0395 & 0.0686 & 0.0988 & 0.1267 \\
 & $\widehat{\M}_{\mathrm{nuc}}$ & 0.0000 & 0.0020 & 0.0032 & 0.0058 & 0.0101 & 0.0181 & 0.0395 & 0.0740 & 0.1372 & 0.3455 \\
\midrule
\multirow{4}{*}{$20$} & $\widetilde{\M}_{\mathrm{fac}}$ & 0.0140 & 0.0140 & 0.0142 & 0.0148 & 0.0171 & 0.0253 & 0.0528 & 0.0807 & 0.1080 & 0.1317 \\
 & $\widehat{\M}_{\mathrm{fac}}$ & 0.0000 & 0.0064 & 0.0088 & 0.0137 & 0.0196 & 0.0292 & 0.0531 & 0.0910 & 0.1665 & 0.4315 \\
 & $\widetilde{\M}_{\mathrm{nuc}}$ & 0.0131 & 0.0135 & 0.0139 & 0.0148 & 0.0171 & 0.0254 & 0.0528 & 0.0806 & 0.1078 & 0.1310 \\
 & $\widehat{\M}_{\mathrm{nuc}}$ & 0.0000 & 0.0062 & 0.0086 & 0.0134 & 0.0189 & 0.0282 & 0.0528 & 0.0924 & 0.1705 & 0.4285 \\
\bottomrule
\end{tabular}
}
\endgroup
\end{table}

% BEGIN MAJOR REVISION S11
\begin{revision}
\subsection{Heterogeneous variance profiles}
The four scenarios are defined in Section~\ref{main-sec:heterogeneous_simulation}. Scenario numbers $1,2,3,4$ refer to row-specific two-point, row-specific $\chi_1^2$, row-and-column two-point and row-and-column $\chi_1^2$ profiles, respectively. In every scenario $m=500$, $n=250$, $p=0.2$ and $s=r=5$.
The subscripts $W$ identify the proposed estimators with weights fitted from residuals. The superscript $\mathrm{or}$ identifies weights constructed from the true profile. Both weight constructions use lower-quartile flooring and median normalization. Tables~\ref{tab:supp:heterogeneous_fnorm}, \ref{tab:supp:different_noise_specnorm} and \ref{tab:supp:different_noise_maxnorm} give the full grid, including the zero-noise diagnostic. The errors are those of the recovery and refinement fits; Theorem~\ref{main-thm:weighted_unknown_rate} concerns their stabilized outputs.
The estimated-profile and true-profile fits have similar error magnitudes for positive noise, but their ordering depends on the setting and norm. With estimated weights, the zero-noise Frobenius errors range from $0.0006$ to $0.0008$ and the maximum errors are larger. The zero-noise rows are numerical diagnostics outside the positive-noise assumptions of the heterogeneous theorem.

\end{revision}
% END MAJOR REVISION S11
\begin{table}[!htbp]
\centering
\caption{Mean relative Frobenius errors $\mathcal E_F$ under the four heterogeneous variance profiles, with $m=500$, $n=250$, $p=0.2$ and $s=r=5$. The weighted estimators use estimated profiles; the superscript $\mathrm{or}$ denotes true-profile weights. Each entry averages 100 repetitions.}
\label{tab:supp:heterogeneous_fnorm}
\begingroup
\small
\setstretch{1.05}
\setlength{\tabcolsep}{3.2pt}
\resizebox{\linewidth}{!}{%
\begin{tabular}{@{}llrrrrrrrrrr@{}}
\toprule
Scenario & Estimator & $\sigma=0$ & $\sigma=0.01$ & $\sigma=0.02$ & $\sigma=0.05$ & $\sigma=0.1$ & $\sigma=0.2$ & $\sigma=0.5$ & $\sigma=1$ & $\sigma=2$ & $\sigma=5$ \\
\midrule
\multirow{6}{*}{1} & $\widetilde{\M}_{\mathrm{fac},W}$ & 0.0006 & 0.0007 & 0.0013 & 0.0033 & 0.0066 & 0.0131 & 0.0305 & 0.0521 & 0.0819 & 0.1239 \\
 & $\widetilde{\M}_{\mathrm{fac},W}^{\mathrm{or}}$ & 0.0000 & 0.0006 & 0.0013 & 0.0032 & 0.0065 & 0.0129 & 0.0298 & 0.0502 & 0.0802 & 0.1547 \\
 & $\widehat{\M}_{\mathrm{fac}}$ & 0.0000 & 0.0008 & 0.0015 & 0.0038 & 0.0076 & 0.0153 & 0.0382 & 0.0768 & 0.1589 & 0.5235 \\
 & $\widetilde{\M}_{\mathrm{nuc},W}$ & 0.0006 & 0.0007 & 0.0013 & 0.0033 & 0.0066 & 0.0131 & 0.0306 & 0.0521 & 0.0819 & 0.1237 \\
 & $\widetilde{\M}_{\mathrm{nuc},W}^{\mathrm{or}}$ & 0.0000 & 0.0006 & 0.0013 & 0.0032 & 0.0065 & 0.0129 & 0.0298 & 0.0502 & 0.0802 & 0.1546 \\
 & $\widehat{\M}_{\mathrm{nuc}}$ & 0.0000 & 0.0008 & 0.0015 & 0.0038 & 0.0076 & 0.0153 & 0.0382 & 0.0769 & 0.1594 & 0.5254 \\
\midrule
\multirow{6}{*}{2} & $\widetilde{\M}_{\mathrm{fac},W}$ & 0.0007 & 0.0007 & 0.0014 & 0.0034 & 0.0069 & 0.0137 & 0.0329 & 0.0596 & 0.0948 & 0.1345 \\
 & $\widetilde{\M}_{\mathrm{fac},W}^{\mathrm{or}}$ & 0.0000 & 0.0007 & 0.0013 & 0.0034 & 0.0067 & 0.0133 & 0.0319 & 0.0573 & 0.0952 & 0.1811 \\
 & $\widehat{\M}_{\mathrm{fac}}$ & 0.0000 & 0.0008 & 0.0015 & 0.0039 & 0.0077 & 0.0154 & 0.0385 & 0.0773 & 0.1585 & 0.5180 \\
 & $\widetilde{\M}_{\mathrm{nuc},W}$ & 0.0007 & 0.0007 & 0.0014 & 0.0035 & 0.0069 & 0.0137 & 0.0330 & 0.0598 & 0.0952 & 0.1358 \\
 & $\widetilde{\M}_{\mathrm{nuc},W}^{\mathrm{or}}$ & 0.0000 & 0.0007 & 0.0013 & 0.0034 & 0.0067 & 0.0134 & 0.0320 & 0.0574 & 0.0953 & 0.1811 \\
 & $\widehat{\M}_{\mathrm{nuc}}$ & 0.0000 & 0.0008 & 0.0015 & 0.0039 & 0.0077 & 0.0154 & 0.0385 & 0.0774 & 0.1589 & 0.5207 \\
\midrule
\multirow{6}{*}{3} & $\widetilde{\M}_{\mathrm{fac},W}$ & 0.0006 & 0.0004 & 0.0007 & 0.0019 & 0.0038 & 0.0075 & 0.0186 & 0.0353 & 0.0616 & 0.0992 \\
 & $\widetilde{\M}_{\mathrm{fac},W}^{\mathrm{or}}$ & 0.0000 & 0.0003 & 0.0007 & 0.0017 & 0.0033 & 0.0067 & 0.0166 & 0.0314 & 0.0528 & 0.0962 \\
 & $\widehat{\M}_{\mathrm{fac}}$ & 0.0000 & 0.0008 & 0.0015 & 0.0038 & 0.0076 & 0.0152 & 0.0382 & 0.0778 & 0.1703 & 0.9152 \\
 & $\widetilde{\M}_{\mathrm{nuc},W}$ & 0.0006 & 0.0004 & 0.0008 & 0.0019 & 0.0038 & 0.0076 & 0.0187 & 0.0360 & 0.0654 & 0.1008 \\
 & $\widetilde{\M}_{\mathrm{nuc},W}^{\mathrm{or}}$ & 0.0000 & 0.0003 & 0.0007 & 0.0017 & 0.0033 & 0.0067 & 0.0166 & 0.0314 & 0.0528 & 0.0960 \\
 & $\widehat{\M}_{\mathrm{nuc}}$ & 0.0000 & 0.0008 & 0.0015 & 0.0038 & 0.0076 & 0.0152 & 0.0382 & 0.0779 & 0.1709 & 0.9008 \\
\midrule
\multirow{6}{*}{4} & $\widetilde{\M}_{\mathrm{fac},W}$ & 0.0008 & 0.0005 & 0.0009 & 0.0024 & 0.0047 & 0.0094 & 0.0231 & 0.0444 & 0.0780 & 0.1185 \\
 & $\widetilde{\M}_{\mathrm{fac},W}^{\mathrm{or}}$ & 0.0000 & 0.0004 & 0.0007 & 0.0017 & 0.0035 & 0.0069 & 0.0171 & 0.0328 & 0.0588 & 0.1220 \\
 & $\widehat{\M}_{\mathrm{fac}}$ & 0.0000 & 0.0008 & 0.0015 & 0.0038 & 0.0077 & 0.0154 & 0.0385 & 0.0779 & 0.1658 & 0.8036 \\
 & $\widetilde{\M}_{\mathrm{nuc},W}$ & 0.0008 & 0.0005 & 0.0010 & 0.0024 & 0.0048 & 0.0097 & 0.0239 & 0.0459 & 0.0836 & 0.1319 \\
 & $\widetilde{\M}_{\mathrm{nuc},W}^{\mathrm{or}}$ & 0.0000 & 0.0004 & 0.0007 & 0.0017 & 0.0035 & 0.0069 & 0.0170 & 0.0328 & 0.0585 & 0.1205 \\
 & $\widehat{\M}_{\mathrm{nuc}}$ & 0.0000 & 0.0008 & 0.0015 & 0.0038 & 0.0077 & 0.0154 & 0.0385 & 0.0780 & 0.1663 & 0.7769 \\
\bottomrule
\end{tabular}
}
\endgroup
\end{table}
\begin{table}[!htbp]
\centering
\caption{Mean relative spectral errors $\mathcal E_{sp}$ under the four heterogeneous variance profiles, with $m=500$, $n=250$, $p=0.2$ and $s=r=5$. The weighted estimators use estimated profiles; the superscript $\mathrm{or}$ denotes true-profile weights. Each entry averages 100 repetitions.}
\label{tab:supp:different_noise_specnorm}
\begingroup
\small
\setstretch{1.05}
\setlength{\tabcolsep}{3.2pt}
\resizebox{\linewidth}{!}{%
\begin{tabular}{@{}llrrrrrrrrrr@{}}
\toprule
Scenario & Estimator & $\sigma=0$ & $\sigma=0.01$ & $\sigma=0.02$ & $\sigma=0.05$ & $\sigma=0.1$ & $\sigma=0.2$ & $\sigma=0.5$ & $\sigma=1$ & $\sigma=2$ & $\sigma=5$ \\
\midrule
\multirow{6}{*}{1} & $\widetilde{\M}_{\mathrm{fac},W}$ & 0.0010 & 0.0006 & 0.0011 & 0.0029 & 0.0057 & 0.0113 & 0.0261 & 0.0429 & 0.0634 & 0.0889 \\
 & $\widetilde{\M}_{\mathrm{fac},W}^{\mathrm{or}}$ & 0.0000 & 0.0006 & 0.0011 & 0.0028 & 0.0056 & 0.0112 & 0.0257 & 0.0416 & 0.0621 & 0.1114 \\
 & $\widehat{\M}_{\mathrm{fac}}$ & 0.0000 & 0.0006 & 0.0011 & 0.0029 & 0.0057 & 0.0114 & 0.0285 & 0.0574 & 0.1192 & 0.4281 \\
 & $\widetilde{\M}_{\mathrm{nuc},W}$ & 0.0010 & 0.0006 & 0.0011 & 0.0029 & 0.0057 & 0.0113 & 0.0261 & 0.0429 & 0.0633 & 0.0886 \\
 & $\widetilde{\M}_{\mathrm{nuc},W}^{\mathrm{or}}$ & 0.0000 & 0.0006 & 0.0011 & 0.0028 & 0.0057 & 0.0112 & 0.0257 & 0.0416 & 0.0621 & 0.1114 \\
 & $\widehat{\M}_{\mathrm{nuc}}$ & 0.0000 & 0.0006 & 0.0011 & 0.0029 & 0.0057 & 0.0114 & 0.0285 & 0.0575 & 0.1198 & 0.4326 \\
\midrule
\multirow{6}{*}{2} & $\widetilde{\M}_{\mathrm{fac},W}$ & 0.0013 & 0.0006 & 0.0011 & 0.0029 & 0.0057 & 0.0113 & 0.0266 & 0.0467 & 0.0705 & 0.0956 \\
 & $\widetilde{\M}_{\mathrm{fac},W}^{\mathrm{or}}$ & 0.0000 & 0.0006 & 0.0011 & 0.0028 & 0.0056 & 0.0112 & 0.0262 & 0.0454 & 0.0715 & 0.1305 \\
 & $\widehat{\M}_{\mathrm{fac}}$ & 0.0000 & 0.0006 & 0.0011 & 0.0029 & 0.0057 & 0.0114 & 0.0285 & 0.0573 & 0.1185 & 0.4567 \\
 & $\widetilde{\M}_{\mathrm{nuc},W}$ & 0.0012 & 0.0006 & 0.0011 & 0.0029 & 0.0057 & 0.0113 & 0.0267 & 0.0467 & 0.0706 & 0.0958 \\
 & $\widetilde{\M}_{\mathrm{nuc},W}^{\mathrm{or}}$ & 0.0000 & 0.0006 & 0.0011 & 0.0028 & 0.0056 & 0.0112 & 0.0262 & 0.0454 & 0.0715 & 0.1305 \\
 & $\widehat{\M}_{\mathrm{nuc}}$ & 0.0000 & 0.0006 & 0.0011 & 0.0029 & 0.0057 & 0.0114 & 0.0286 & 0.0574 & 0.1190 & 0.4635 \\
\midrule
\multirow{6}{*}{3} & $\widetilde{\M}_{\mathrm{fac},W}$ & 0.0011 & 0.0003 & 0.0006 & 0.0015 & 0.0030 & 0.0060 & 0.0148 & 0.0275 & 0.0468 & 0.0726 \\
 & $\widetilde{\M}_{\mathrm{fac},W}^{\mathrm{or}}$ & 0.0000 & 0.0003 & 0.0005 & 0.0013 & 0.0026 & 0.0053 & 0.0130 & 0.0242 & 0.0399 & 0.0709 \\
 & $\widehat{\M}_{\mathrm{fac}}$ & 0.0000 & 0.0006 & 0.0012 & 0.0030 & 0.0060 & 0.0121 & 0.0303 & 0.0618 & 0.1358 & 0.9912 \\
 & $\widetilde{\M}_{\mathrm{nuc},W}$ & 0.0011 & 0.0003 & 0.0006 & 0.0015 & 0.0030 & 0.0060 & 0.0149 & 0.0275 & 0.0473 & 0.0728 \\
 & $\widetilde{\M}_{\mathrm{nuc},W}^{\mathrm{or}}$ & 0.0000 & 0.0003 & 0.0005 & 0.0013 & 0.0026 & 0.0053 & 0.0130 & 0.0241 & 0.0398 & 0.0708 \\
 & $\widehat{\M}_{\mathrm{nuc}}$ & 0.0000 & 0.0006 & 0.0012 & 0.0030 & 0.0060 & 0.0121 & 0.0303 & 0.0619 & 0.1364 & 0.9580 \\
\midrule
\multirow{6}{*}{4} & $\widetilde{\M}_{\mathrm{fac},W}$ & 0.0013 & 0.0004 & 0.0008 & 0.0019 & 0.0038 & 0.0075 & 0.0183 & 0.0343 & 0.0582 & 0.0856 \\
 & $\widetilde{\M}_{\mathrm{fac},W}^{\mathrm{or}}$ & 0.0000 & 0.0003 & 0.0006 & 0.0014 & 0.0029 & 0.0057 & 0.0139 & 0.0261 & 0.0452 & 0.0924 \\
 & $\widehat{\M}_{\mathrm{fac}}$ & 0.0000 & 0.0006 & 0.0012 & 0.0030 & 0.0059 & 0.0119 & 0.0297 & 0.0603 & 0.1298 & 0.9106 \\
 & $\widetilde{\M}_{\mathrm{nuc},W}$ & 0.0013 & 0.0004 & 0.0008 & 0.0019 & 0.0038 & 0.0076 & 0.0184 & 0.0346 & 0.0599 & 0.0874 \\
 & $\widetilde{\M}_{\mathrm{nuc},W}^{\mathrm{or}}$ & 0.0000 & 0.0003 & 0.0006 & 0.0014 & 0.0029 & 0.0057 & 0.0141 & 0.0261 & 0.0450 & 0.0895 \\
 & $\widehat{\M}_{\mathrm{nuc}}$ & 0.0000 & 0.0006 & 0.0012 & 0.0030 & 0.0059 & 0.0119 & 0.0298 & 0.0604 & 0.1302 & 0.8561 \\
\bottomrule
\end{tabular}
}
\endgroup
\end{table}
\begin{table}[!htbp]
\centering
\caption{Mean relative maximum errors $\mathcal E_\infty$ under the four heterogeneous variance profiles, with $m=500$, $n=250$, $p=0.2$ and $s=r=5$. The weighted estimators use estimated profiles; the superscript $\mathrm{or}$ denotes true-profile weights. Each entry averages 100 repetitions.}
\label{tab:supp:different_noise_maxnorm}
\begingroup
\small
\setstretch{1.05}
\setlength{\tabcolsep}{3.2pt}
\resizebox{\linewidth}{!}{%
\begin{tabular}{@{}llrrrrrrrrrr@{}}
\toprule
Scenario & Estimator & $\sigma=0$ & $\sigma=0.01$ & $\sigma=0.02$ & $\sigma=0.05$ & $\sigma=0.1$ & $\sigma=0.2$ & $\sigma=0.5$ & $\sigma=1$ & $\sigma=2$ & $\sigma=5$ \\
\midrule
\multirow{6}{*}{1} & $\widetilde{\M}_{\mathrm{fac},W}$ & 0.0056 & 0.0012 & 0.0025 & 0.0062 & 0.0124 & 0.0243 & 0.0533 & 0.0824 & 0.1134 & 0.1302 \\
 & $\widetilde{\M}_{\mathrm{fac},W}^{\mathrm{or}}$ & 0.0000 & 0.0012 & 0.0024 & 0.0061 & 0.0122 & 0.0240 & 0.0526 & 0.0805 & 0.1100 & 0.1623 \\
 & $\widehat{\M}_{\mathrm{fac}}$ & 0.0000 & 0.0012 & 0.0025 & 0.0062 & 0.0123 & 0.0246 & 0.0615 & 0.1233 & 0.2548 & 0.9101 \\
 & $\widetilde{\M}_{\mathrm{nuc},W}$ & 0.0056 & 0.0012 & 0.0025 & 0.0062 & 0.0124 & 0.0243 & 0.0533 & 0.0824 & 0.1135 & 0.1309 \\
 & $\widetilde{\M}_{\mathrm{nuc},W}^{\mathrm{or}}$ & 0.0000 & 0.0012 & 0.0024 & 0.0061 & 0.0122 & 0.0240 & 0.0525 & 0.0805 & 0.1099 & 0.1621 \\
 & $\widehat{\M}_{\mathrm{nuc}}$ & 0.0000 & 0.0012 & 0.0025 & 0.0062 & 0.0124 & 0.0247 & 0.0616 & 0.1237 & 0.2562 & 0.9181 \\
\midrule
\multirow{6}{*}{2} & $\widetilde{\M}_{\mathrm{fac},W}$ & 0.0076 & 0.0013 & 0.0026 & 0.0064 & 0.0127 & 0.0244 & 0.0521 & 0.0822 & 0.1091 & 0.1274 \\
 & $\widetilde{\M}_{\mathrm{fac},W}^{\mathrm{or}}$ & 0.0000 & 0.0013 & 0.0026 & 0.0063 & 0.0123 & 0.0237 & 0.0505 & 0.0777 & 0.1092 & 0.1762 \\
 & $\widehat{\M}_{\mathrm{fac}}$ & 0.0000 & 0.0013 & 0.0025 & 0.0063 & 0.0125 & 0.0250 & 0.0624 & 0.1257 & 0.2672 & 1.5306 \\
 & $\widetilde{\M}_{\mathrm{nuc},W}$ & 0.0076 & 0.0013 & 0.0026 & 0.0064 & 0.0127 & 0.0244 & 0.0521 & 0.0822 & 0.1092 & 0.1289 \\
 & $\widetilde{\M}_{\mathrm{nuc},W}^{\mathrm{or}}$ & 0.0000 & 0.0013 & 0.0026 & 0.0063 & 0.0123 & 0.0237 & 0.0505 & 0.0776 & 0.1091 & 0.1748 \\
 & $\widehat{\M}_{\mathrm{nuc}}$ & 0.0000 & 0.0013 & 0.0025 & 0.0063 & 0.0125 & 0.0250 & 0.0625 & 0.1261 & 0.2700 & 1.5627 \\
\midrule
\multirow{6}{*}{3} & $\widetilde{\M}_{\mathrm{fac},W}$ & 0.0062 & 0.0007 & 0.0014 & 0.0035 & 0.0070 & 0.0139 & 0.0333 & 0.0583 & 0.0916 & 0.1262 \\
 & $\widetilde{\M}_{\mathrm{fac},W}^{\mathrm{or}}$ & 0.0000 & 0.0006 & 0.0013 & 0.0031 & 0.0061 & 0.0120 & 0.0290 & 0.0514 & 0.0797 & 0.1248 \\
 & $\widehat{\M}_{\mathrm{fac}}$ & 0.0000 & 0.0013 & 0.0027 & 0.0067 & 0.0133 & 0.0267 & 0.0670 & 0.1376 & 0.3123 & 4.0331 \\
 & $\widetilde{\M}_{\mathrm{nuc},W}$ & 0.0062 & 0.0007 & 0.0014 & 0.0035 & 0.0070 & 0.0139 & 0.0334 & 0.0593 & 0.1231 & 0.1310 \\
 & $\widetilde{\M}_{\mathrm{nuc},W}^{\mathrm{or}}$ & 0.0000 & 0.0006 & 0.0013 & 0.0031 & 0.0061 & 0.0120 & 0.0290 & 0.0513 & 0.0796 & 0.1235 \\
 & $\widehat{\M}_{\mathrm{nuc}}$ & 0.0000 & 0.0013 & 0.0027 & 0.0067 & 0.0133 & 0.0267 & 0.0671 & 0.1377 & 0.3140 & 3.6400 \\
\midrule
\multirow{6}{*}{4} & $\widetilde{\M}_{\mathrm{fac},W}$ & 0.0076 & 0.0009 & 0.0019 & 0.0046 & 0.0088 & 0.0172 & 0.0389 & 0.0646 & 0.0972 & 0.1211 \\
 & $\widetilde{\M}_{\mathrm{fac},W}^{\mathrm{or}}$ & 0.0001 & 0.0008 & 0.0014 & 0.0035 & 0.0069 & 0.0133 & 0.0299 & 0.0512 & 0.0776 & 0.1804 \\
 & $\widehat{\M}_{\mathrm{fac}}$ & 0.0000 & 0.0013 & 0.0027 & 0.0067 & 0.0134 & 0.0269 & 0.0674 & 0.1378 & 0.3139 & 5.7521 \\
 & $\widetilde{\M}_{\mathrm{nuc},W}$ & 0.0076 & 0.0009 & 0.0019 & 0.0046 & 0.0088 & 0.0172 & 0.0391 & 0.0662 & 0.1212 & 0.2337 \\
 & $\widetilde{\M}_{\mathrm{nuc},W}^{\mathrm{or}}$ & 0.0001 & 0.0008 & 0.0014 & 0.0035 & 0.0069 & 0.0133 & 0.0294 & 0.0509 & 0.0773 & 0.1553 \\
 & $\widehat{\M}_{\mathrm{nuc}}$ & 0.0000 & 0.0013 & 0.0027 & 0.0067 & 0.0135 & 0.0269 & 0.0675 & 0.1380 & 0.3156 & 5.0720 \\
\bottomrule
\end{tabular}
}
\endgroup
\end{table}

Table~\ref{tab:supp:weight_comparison} compares the weighted and unweighted implementations using the same starting inputs and observation masks. Estimated weights reduce the average Frobenius error over the displayed noise levels. The comparison assesses the effect of weighting within the RSM procedure.
\begin{table}[!htbp]
\centering
\caption{Mean relative Frobenius errors averaged over the four heterogeneous scenarios. Each scenario has 100 repetitions.}
\label{tab:supp:weight_comparison}
\begingroup
\small
\setstretch{1.05}
\setlength{\tabcolsep}{3.2pt}
\begin{tabular}{@{}llrrrrr@{}}
\toprule
Implementation & Estimator & $\sigma=0.01$ & $\sigma=0.1$ & $\sigma=1$ & $\sigma=2$ & $\sigma=5$ \\
\midrule
\multirow{3}{*}{Factor} & Unweighted RSM & 0.0008 & 0.0077 & 0.0699 & 0.1102 & 0.1498 \\
 & Estimated weights & 0.0005 & 0.0055 & 0.0478 & 0.0791 & 0.1190 \\
 & True-profile weights & 0.0005 & 0.0050 & 0.0429 & 0.0718 & 0.1385 \\
\midrule
\multirow{3}{*}{Nuclear} & Unweighted RSM & 0.0010 & 0.0098 & 0.0938 & 0.1634 & 0.3649 \\
 & Estimated weights & 0.0006 & 0.0055 & 0.0485 & 0.0815 & 0.1230 \\
 & True-profile weights & 0.0005 & 0.0050 & 0.0429 & 0.0717 & 0.1380 \\
\bottomrule
\end{tabular}
\endgroup
\end{table}

% BEGIN MAJOR REVISION S12
\begin{revision}
\subsection{Real-data experiments}
\label{sec:real_data_implementation}
The real-data experiments use the preprocessing and data partitions described in Sections~\ref{main-sec:netflix_data} and \ref{main-sec:amazon_data}. Rank and penalty parameters are selected by validation RMSE, and test entries are used only for evaluation. Test RMSE on the original rating scale is defined in \eqref{main-eq:real_data_rmse}. Predictions clipped to $[1,5]$ are evaluated on the same test entries without retuning.

\paragraph{Weighted RSM.}
The scale profile is estimated from the residuals of an RSM pilot fitted at the rank selected by factor RSM. Two-way median polish of the logarithms of the absolute training residuals gives an exponentiated scale fit $\widehat\tau_{ij}^{\mathrm{raw}}$. We use
\begin{align}
 c_{1/4}&=Q_{1/4}\bigl(\{\widehat\tau_{ij}^{\mathrm{raw}}\}_{i,j}\bigr),\nonumber\\
 \widehat T_{ij}&=\widehat\tau_{ij}^{\mathrm{raw}}\vee c_{1/4},
 &W_{ij}&=\widehat T_{ij}^{-1},
 \label{eq:real_data_inverse_scale_weights}
\end{align}
and form the weighted residual $W\odot\{Y-P_{\Omega_{\mathrm{tr}}}(M)\}$, where $\Omega_{\mathrm{tr}}$ is the training index set. The weights remain fixed during fitting. The factor and nuclear-norm implementations use the rank selected by factor RSM in their recovery initialization; the weighted nuclear penalty is selected separately by validation RMSE.

\end{revision}
% END MAJOR REVISION S12

% BEGIN MAJOR REVISION S13
\begin{revision}
\subsection{Netflix rating data}
Movies with more than 25,000 ratings are retained. Scenario 1 contains 3,036 users with more than 700 ratings and 1,034 movies. Scenarios 2 and 3 each contain 3,000 users sampled from those with more than 200 and 75 ratings, respectively, with the same 1,034 movies.

Table~\ref{tab:netflix_absolute} compares the test RMSE of RSM and weighted RSM. Weighting slightly reduces the factor error in the three reported scenarios. For nuclear-norm RSM, the weighted fit has smaller errors in scenarios 1 and 2 and a larger error in scenario 3.

\end{revision}
% END MAJOR REVISION S13


\begin{table}[!htbp]
\centering
\caption{Test RMSE for RSM and weighted RSM on Netflix rating data, using the data partition in Table~\ref{main-tab:netflix_data}. Unclipped and clipped predictions are shown side by side; clipping restricts predictions to $[1,5]$ without retuning. Boldface marks the smallest displayed value in each row, including ties.}
\label{tab:netflix_absolute}
\begingroup
\small
\setstretch{1.05}
\setlength{\tabcolsep}{4pt}
\begin{tabular}{@{}crrrrrrrr@{}}
\toprule
 & \multicolumn{4}{c}{RSM} & \multicolumn{4}{c}{Weighted RSM} \\
\cmidrule(lr){2-5}\cmidrule(lr){6-9}
 & \multicolumn{2}{c}{Factor} & \multicolumn{2}{c}{Nuclear} & \multicolumn{2}{c}{Factor} & \multicolumn{2}{c}{Nuclear} \\
\cmidrule(lr){2-3}\cmidrule(lr){4-5}\cmidrule(lr){6-7}\cmidrule(lr){8-9}
Scenario & Raw & Clipped & Raw & Clipped & Raw & Clipped & Raw & Clipped \\
\midrule
1 & 0.7692 & 0.7675 & 0.7619 & 0.7617 & 0.7676 & 0.7660 & 0.7603 & \textbf{0.7600} \\
2 & 0.8101 & 0.8089 & 0.8042 & 0.8037 & 0.8087 & 0.8075 & 0.8037 & \textbf{0.8031} \\
3 & 0.8426 & 0.8409 & 0.8354 & \textbf{0.8351} & 0.8407 & 0.8392 & 0.8380 & 0.8368 \\
\bottomrule
\end{tabular}
\endgroup
\end{table}

% BEGIN MAJOR REVISION S14
\begin{revision}
\subsection{Amazon review data}
The three scenarios correspond to movies, books and electronics reviews. We retain the 1,000 most reviewed items in each category and sample 3,000 users from those with more than $14,12,12$ reviews in the respective scenarios.

All RSM factor fits select $s=1$ from the candidate ranks $1,\ldots,10$, and the weighted implementations use this rank. Table~\ref{tab:amazon_absolute} reports their test RMSE on the original rating scale. Weighted RSM has larger test errors than RSM in the three reported scenarios for both implementations, before and after clipping.

\end{revision}
% END MAJOR REVISION S14


\begin{table}[!htbp]
\centering
\caption{Test RMSE for RSM and weighted RSM on Amazon rating data, using the data partition in Table~\ref{main-tab:amazon_data}. Unclipped and clipped predictions are shown side by side; clipping restricts predictions to $[1,5]$ without retuning. Boldface marks the smallest displayed value in each row, including ties.}
\label{tab:amazon_absolute}
\begingroup
\small
\setstretch{1.05}
\setlength{\tabcolsep}{4pt}
\begin{tabular}{@{}crrrrrrrr@{}}
\toprule
 & \multicolumn{4}{c}{RSM} & \multicolumn{4}{c}{Weighted RSM} \\
\cmidrule(lr){2-5}\cmidrule(lr){6-9}
 & \multicolumn{2}{c}{Factor} & \multicolumn{2}{c}{Nuclear} & \multicolumn{2}{c}{Factor} & \multicolumn{2}{c}{Nuclear} \\
\cmidrule(lr){2-3}\cmidrule(lr){4-5}\cmidrule(lr){6-7}\cmidrule(lr){8-9}
Scenario & Raw & Clipped & Raw & Clipped & Raw & Clipped & Raw & Clipped \\
\midrule
1 & 0.9219 & 0.9132 & 0.9210 & \textbf{0.9128} & 0.9377 & 0.9313 & 0.9372 & 0.9309 \\
2 & 0.8513 & \textbf{0.8408} & 0.8510 & \textbf{0.8408} & 0.8872 & 0.8736 & 0.8866 & 0.8729 \\
3 & 0.9817 & 0.9710 & 0.9802 & \textbf{0.9695} & 1.0024 & 0.9934 & 1.0018 & 0.9942 \\
\bottomrule
\end{tabular}
\endgroup
\end{table}

\section{Additional simulation experiments}
The following experiments assess sensitivity to the noise structure, signal conditioning and working rank. The unweighted implementations are used for unstructured variance perturbations and dependent noise. Each table reports the three norm criteria under the simulated-data initialization and tuning protocol in Section~\ref{main-sec:simulation_implementation}.

\subsection{Heterogeneous noise without structure}
We generate
\begin{align*}
 \bH=\bm E\odot\bH_{\mathrm{homo}},\qquad
 E_{ij}\stackrel{\mathrm{i.i.d.}}{\sim}\operatorname{Uniform}(1-h,1+h),
 \quad 1\le i\le m,\quad1\le j\le n,
\end{align*}
where $\bH_{\mathrm{homo}}$ has homogeneous noise level $\sigma=1$ and is independent of $\bm E$. Conditional variances lie in $[(1-h)^2,(1+h)^2]$. We use $h\in\{0.1,0.2,0.3\}$, $r=5$, $p=0.2$ and $n/m=0.5$. The parameter $h$ measures heterogeneity and is distinct from the aspect ratio. Tables~\ref{tab:supp:hete_noise_fnorm}, \ref{tab:supp:hete_noise_specnorm} and \ref{tab:supp:hete_noise_maxnorm} show that the unweighted proposed estimators retain smaller mean errors than their baselines over these perturbations without estimating a variance profile.
\begin{table}[!htbp]
\centering
\caption{Mean relative Frobenius errors $\mathcal E_F$ under unstructured variance perturbations for $\sigma=1$, $p=0.2$, $r=5$ and $n/m=0.5$. Each entry averages 100 repetitions.}
\label{tab:supp:hete_noise_fnorm}
\begingroup
\small
\setstretch{1.05}
\setlength{\tabcolsep}{3.2pt}
\resizebox{\linewidth}{!}{%
\begin{tabular}{@{}llrrrrrrrrr@{}}
\toprule
$h$ & Estimator & $m=200$ & $m=300$ & $m=400$ & $m=500$ & $m=600$ & $m=700$ & $m=800$ & $m=900$ & $m=1000$ \\
\midrule
\multirow{4}{*}{$0.1$} & $\widetilde{\M}_{\mathrm{fac}}$ & 0.1214 & 0.0940 & 0.0789 & 0.0693 & 0.0628 & 0.0576 & 0.0535 & 0.0503 & 0.0475 \\
 & $\widehat{\M}_{\mathrm{fac}}$ & 0.1365 & 0.1023 & 0.0842 & 0.0746 & 0.0668 & 0.0605 & 0.0557 & 0.0519 & 0.0486 \\
 & $\widetilde{\M}_{\mathrm{nuc}}$ & 0.1209 & 0.0937 & 0.0788 & 0.0692 & 0.0627 & 0.0575 & 0.0534 & 0.0502 & 0.0475 \\
 & $\widehat{\M}_{\mathrm{nuc}}$ & 0.1310 & 0.1021 & 0.0840 & 0.0743 & 0.0674 & 0.0619 & 0.0577 & 0.0543 & 0.0513 \\
\midrule
\multirow{4}{*}{$0.2$} & $\widetilde{\M}_{\mathrm{fac}}$ & 0.1219 & 0.0943 & 0.0792 & 0.0696 & 0.0630 & 0.0578 & 0.0537 & 0.0505 & 0.0477 \\
 & $\widehat{\M}_{\mathrm{fac}}$ & 0.1373 & 0.1028 & 0.0846 & 0.0750 & 0.0671 & 0.0608 & 0.0560 & 0.0521 & 0.0488 \\
 & $\widetilde{\M}_{\mathrm{nuc}}$ & 0.1214 & 0.0941 & 0.0791 & 0.0695 & 0.0630 & 0.0577 & 0.0537 & 0.0504 & 0.0477 \\
 & $\widehat{\M}_{\mathrm{nuc}}$ & 0.1319 & 0.1026 & 0.0844 & 0.0746 & 0.0678 & 0.0622 & 0.0580 & 0.0546 & 0.0516 \\
\midrule
\multirow{4}{*}{$0.3$} & $\widetilde{\M}_{\mathrm{fac}}$ & 0.1226 & 0.0949 & 0.0797 & 0.0700 & 0.0634 & 0.0581 & 0.0541 & 0.0508 & 0.0480 \\
 & $\widehat{\M}_{\mathrm{fac}}$ & 0.1385 & 0.1037 & 0.0852 & 0.0756 & 0.0677 & 0.0613 & 0.0564 & 0.0525 & 0.0492 \\
 & $\widetilde{\M}_{\mathrm{nuc}}$ & 0.1222 & 0.0947 & 0.0796 & 0.0699 & 0.0634 & 0.0581 & 0.0540 & 0.0508 & 0.0480 \\
 & $\widehat{\M}_{\mathrm{nuc}}$ & 0.1332 & 0.1035 & 0.0851 & 0.0752 & 0.0683 & 0.0627 & 0.0585 & 0.0550 & 0.0520 \\
\bottomrule
\end{tabular}
}
\endgroup
\end{table}
\begin{table}[!htbp]
\centering
\caption{Mean relative spectral errors $\mathcal E_{sp}$ under unstructured variance perturbations for $\sigma=1$, $p=0.2$, $r=5$ and $n/m=0.5$. Each entry averages 100 repetitions.}
\label{tab:supp:hete_noise_specnorm}
\begingroup
\small
\setstretch{1.05}
\setlength{\tabcolsep}{3.2pt}
\resizebox{\linewidth}{!}{%
\begin{tabular}{@{}llrrrrrrrrr@{}}
\toprule
$h$ & Estimator & $m=200$ & $m=300$ & $m=400$ & $m=500$ & $m=600$ & $m=700$ & $m=800$ & $m=900$ & $m=1000$ \\
\midrule
\multirow{4}{*}{$0.1$} & $\widetilde{\M}_{\mathrm{fac}}$ & 0.0900 & 0.0676 & 0.0561 & 0.0487 & 0.0439 & 0.0398 & 0.0369 & 0.0346 & 0.0325 \\
 & $\widehat{\M}_{\mathrm{fac}}$ & 0.1021 & 0.0740 & 0.0600 & 0.0526 & 0.0469 & 0.0421 & 0.0387 & 0.0359 & 0.0335 \\
 & $\widetilde{\M}_{\mathrm{nuc}}$ & 0.0897 & 0.0675 & 0.0559 & 0.0486 & 0.0438 & 0.0397 & 0.0368 & 0.0346 & 0.0325 \\
 & $\widehat{\M}_{\mathrm{nuc}}$ & 0.0975 & 0.0736 & 0.0597 & 0.0522 & 0.0472 & 0.0429 & 0.0398 & 0.0373 & 0.0351 \\
\midrule
\multirow{4}{*}{$0.2$} & $\widetilde{\M}_{\mathrm{fac}}$ & 0.0904 & 0.0679 & 0.0563 & 0.0489 & 0.0440 & 0.0400 & 0.0370 & 0.0348 & 0.0326 \\
 & $\widehat{\M}_{\mathrm{fac}}$ & 0.1026 & 0.0744 & 0.0603 & 0.0528 & 0.0471 & 0.0423 & 0.0389 & 0.0360 & 0.0336 \\
 & $\widetilde{\M}_{\mathrm{nuc}}$ & 0.0900 & 0.0678 & 0.0562 & 0.0488 & 0.0440 & 0.0399 & 0.0370 & 0.0348 & 0.0326 \\
 & $\widehat{\M}_{\mathrm{nuc}}$ & 0.0981 & 0.0740 & 0.0600 & 0.0524 & 0.0474 & 0.0431 & 0.0400 & 0.0375 & 0.0353 \\
\midrule
\multirow{4}{*}{$0.3$} & $\widetilde{\M}_{\mathrm{fac}}$ & 0.0909 & 0.0683 & 0.0567 & 0.0492 & 0.0443 & 0.0402 & 0.0373 & 0.0350 & 0.0329 \\
 & $\widehat{\M}_{\mathrm{fac}}$ & 0.1036 & 0.0750 & 0.0608 & 0.0533 & 0.0475 & 0.0426 & 0.0392 & 0.0363 & 0.0339 \\
 & $\widetilde{\M}_{\mathrm{nuc}}$ & 0.0906 & 0.0682 & 0.0566 & 0.0492 & 0.0442 & 0.0401 & 0.0372 & 0.0350 & 0.0328 \\
 & $\widehat{\M}_{\mathrm{nuc}}$ & 0.0991 & 0.0747 & 0.0605 & 0.0528 & 0.0478 & 0.0435 & 0.0403 & 0.0378 & 0.0356 \\
\bottomrule
\end{tabular}
}
\endgroup
\end{table}
\begin{table}[!htbp]
\centering
\caption{Mean relative maximum errors $\mathcal E_\infty$ under unstructured variance perturbations for $\sigma=1$, $p=0.2$, $r=5$ and $n/m=0.5$. Each entry averages 100 repetitions.}
\label{tab:supp:hete_noise_maxnorm}
\begingroup
\small
\setstretch{1.05}
\setlength{\tabcolsep}{3.2pt}
\resizebox{\linewidth}{!}{%
\begin{tabular}{@{}llrrrrrrrrr@{}}
\toprule
$h$ & Estimator & $m=200$ & $m=300$ & $m=400$ & $m=500$ & $m=600$ & $m=700$ & $m=800$ & $m=900$ & $m=1000$ \\
\midrule
\multirow{4}{*}{$0.1$} & $\widetilde{\M}_{\mathrm{fac}}$ & 0.1224 & 0.0890 & 0.0741 & 0.0641 & 0.0576 & 0.0524 & 0.0486 & 0.0450 & 0.0432 \\
 & $\widehat{\M}_{\mathrm{fac}}$ & 0.1439 & 0.0992 & 0.0801 & 0.0688 & 0.0611 & 0.0551 & 0.0498 & 0.0462 & 0.0436 \\
 & $\widetilde{\M}_{\mathrm{nuc}}$ & 0.1216 & 0.0889 & 0.0741 & 0.0641 & 0.0576 & 0.0523 & 0.0485 & 0.0450 & 0.0431 \\
 & $\widehat{\M}_{\mathrm{nuc}}$ & 0.1341 & 0.0989 & 0.0800 & 0.0680 & 0.0618 & 0.0569 & 0.0519 & 0.0493 & 0.0462 \\
\midrule
\multirow{4}{*}{$0.2$} & $\widetilde{\M}_{\mathrm{fac}}$ & 0.1232 & 0.0895 & 0.0744 & 0.0644 & 0.0579 & 0.0527 & 0.0489 & 0.0453 & 0.0433 \\
 & $\widehat{\M}_{\mathrm{fac}}$ & 0.1454 & 0.1000 & 0.0806 & 0.0691 & 0.0614 & 0.0555 & 0.0501 & 0.0464 & 0.0437 \\
 & $\widetilde{\M}_{\mathrm{nuc}}$ & 0.1225 & 0.0894 & 0.0745 & 0.0643 & 0.0578 & 0.0527 & 0.0488 & 0.0453 & 0.0432 \\
 & $\widehat{\M}_{\mathrm{nuc}}$ & 0.1356 & 0.0997 & 0.0803 & 0.0683 & 0.0622 & 0.0572 & 0.0521 & 0.0495 & 0.0463 \\
\midrule
\multirow{4}{*}{$0.3$} & $\widetilde{\M}_{\mathrm{fac}}$ & 0.1242 & 0.0903 & 0.0749 & 0.0650 & 0.0583 & 0.0533 & 0.0492 & 0.0456 & 0.0435 \\
 & $\widehat{\M}_{\mathrm{fac}}$ & 0.1471 & 0.1011 & 0.0814 & 0.0698 & 0.0621 & 0.0561 & 0.0505 & 0.0468 & 0.0440 \\
 & $\widetilde{\M}_{\mathrm{nuc}}$ & 0.1237 & 0.0904 & 0.0750 & 0.0648 & 0.0583 & 0.0532 & 0.0491 & 0.0456 & 0.0434 \\
 & $\widehat{\M}_{\mathrm{nuc}}$ & 0.1376 & 0.1008 & 0.0812 & 0.0690 & 0.0628 & 0.0578 & 0.0526 & 0.0499 & 0.0466 \\
\bottomrule
\end{tabular}
}
\endgroup
\end{table}

\subsection{Weakly dependent noise}
We consider $\bH=\bm\Sigma_\varepsilon^{1/2}\bH_{\mathrm{homo}}$, with the two dependence structures
\begin{align*}
 \text{Case 1: }&[\bm\Sigma_\varepsilon^{1/2}]_{ij}=\gamma^{|i-j|},\\
 \text{Case 2: }&[\bm\Sigma_\varepsilon^{1/2}]_{ij}=\mathbf1_{\{|i-j|\le M\}}.
\end{align*}
These constructions represent exponentially decaying and finite-range dependence, as considered in the factor-model literature \citep{fan2024can,choi2025inference}. We set $m=500$, $n=250$, $p=0.2$ and the input noise scale to $\sigma=1$. The values are $\gamma\in\{0.2,0.4,0.6,0.8\}$ in Case 1 and $M\in\{2,3,5,10\}$ in Case 2.
Tables~\ref{tab:supp:dependent_fnorm}, \ref{tab:supp:dependent_specnorm} and \ref{tab:supp:dependent_maxnorm} show stable errors over these dependence settings. The proposed estimators have smaller mean errors than the corresponding baselines in the displayed configurations. These experiments examine sensitivity beyond the independent-noise theory.
\begin{table}[!htbp]
\centering
\caption{Mean relative Frobenius errors $\mathcal E_F$ for the two dependent-noise constructions with $m=500$, $n=250$ and $p=0.2$. Each entry averages 100 repetitions.}
\label{tab:supp:dependent_fnorm}
\begingroup
\small
\setstretch{1.05}
\setlength{\tabcolsep}{3.2pt}
\resizebox{\linewidth}{!}{%
\begin{tabular}{@{}llrrrrrrrr@{}}
\toprule
$r$ & Estimator & $\gamma=0.2$ & $\gamma=0.4$ & $\gamma=0.6$ & $\gamma=0.8$ & $M=2$ & $M=3$ & $M=5$ & $M=10$ \\
\midrule
\multirow{4}{*}{$5$} & $\widetilde{\M}_{\mathrm{fac}}$ & 0.0692 & 0.0693 & 0.0693 & 0.0692 & 0.0693 & 0.0694 & 0.0693 & 0.0692 \\
 & $\widehat{\M}_{\mathrm{fac}}$ & 0.0745 & 0.0746 & 0.0745 & 0.0745 & 0.0746 & 0.0746 & 0.0746 & 0.0745 \\
 & $\widetilde{\M}_{\mathrm{nuc}}$ & 0.0692 & 0.0693 & 0.0695 & 0.0706 & 0.0694 & 0.0696 & 0.0700 & 0.0713 \\
 & $\widehat{\M}_{\mathrm{nuc}}$ & 0.0742 & 0.0742 & 0.0742 & 0.0741 & 0.0742 & 0.0742 & 0.0742 & 0.0741 \\
\midrule
\multirow{4}{*}{$10$} & $\widetilde{\M}_{\mathrm{fac}}$ & 0.0703 & 0.0703 & 0.0703 & 0.0703 & 0.0704 & 0.0704 & 0.0703 & 0.0703 \\
 & $\widehat{\M}_{\mathrm{fac}}$ & 0.0754 & 0.0754 & 0.0754 & 0.0753 & 0.0754 & 0.0754 & 0.0754 & 0.0753 \\
 & $\widetilde{\M}_{\mathrm{nuc}}$ & 0.0702 & 0.0703 & 0.0705 & 0.0709 & 0.0704 & 0.0705 & 0.0707 & 0.0712 \\
 & $\widehat{\M}_{\mathrm{nuc}}$ & 0.0766 & 0.0766 & 0.0766 & 0.0766 & 0.0767 & 0.0766 & 0.0766 & 0.0765 \\
\midrule
\multirow{4}{*}{$20$} & $\widetilde{\M}_{\mathrm{fac}}$ & 0.0766 & 0.0766 & 0.0766 & 0.0765 & 0.0765 & 0.0765 & 0.0765 & 0.0765 \\
 & $\widehat{\M}_{\mathrm{fac}}$ & 0.0853 & 0.0853 & 0.0852 & 0.0852 & 0.0852 & 0.0852 & 0.0852 & 0.0851 \\
 & $\widetilde{\M}_{\mathrm{nuc}}$ & 0.0765 & 0.0766 & 0.0766 & 0.0767 & 0.0766 & 0.0766 & 0.0767 & 0.0768 \\
 & $\widehat{\M}_{\mathrm{nuc}}$ & 0.0868 & 0.0868 & 0.0868 & 0.0867 & 0.0868 & 0.0867 & 0.0868 & 0.0867 \\
\bottomrule
\end{tabular}
}
\endgroup
\end{table}
\begin{table}[!htbp]
\centering
\caption{Mean relative spectral errors $\mathcal E_{sp}$ for the two dependent-noise constructions with $m=500$, $n=250$ and $p=0.2$. Each entry averages 100 repetitions.}
\label{tab:supp:dependent_specnorm}
\begingroup
\small
\setstretch{1.05}
\setlength{\tabcolsep}{3.2pt}
\resizebox{\linewidth}{!}{%
\begin{tabular}{@{}llrrrrrrrr@{}}
\toprule
$r$ & Estimator & $\gamma=0.2$ & $\gamma=0.4$ & $\gamma=0.6$ & $\gamma=0.8$ & $M=2$ & $M=3$ & $M=5$ & $M=10$ \\
\midrule
\multirow{4}{*}{$5$} & $\widetilde{\M}_{\mathrm{fac}}$ & 0.0487 & 0.0486 & 0.0487 & 0.0490 & 0.0487 & 0.0488 & 0.0489 & 0.0492 \\
 & $\widehat{\M}_{\mathrm{fac}}$ & 0.0525 & 0.0526 & 0.0526 & 0.0530 & 0.0526 & 0.0527 & 0.0530 & 0.0532 \\
 & $\widetilde{\M}_{\mathrm{nuc}}$ & 0.0486 & 0.0486 & 0.0487 & 0.0493 & 0.0487 & 0.0488 & 0.0490 & 0.0498 \\
 & $\widehat{\M}_{\mathrm{nuc}}$ & 0.0521 & 0.0522 & 0.0523 & 0.0526 & 0.0522 & 0.0523 & 0.0526 & 0.0528 \\
\midrule
\multirow{4}{*}{$10$} & $\widetilde{\M}_{\mathrm{fac}}$ & 0.0508 & 0.0509 & 0.0510 & 0.0514 & 0.0509 & 0.0510 & 0.0512 & 0.0516 \\
 & $\widehat{\M}_{\mathrm{fac}}$ & 0.0546 & 0.0548 & 0.0550 & 0.0554 & 0.0549 & 0.0549 & 0.0552 & 0.0557 \\
 & $\widetilde{\M}_{\mathrm{nuc}}$ & 0.0508 & 0.0508 & 0.0510 & 0.0516 & 0.0509 & 0.0510 & 0.0513 & 0.0520 \\
 & $\widehat{\M}_{\mathrm{nuc}}$ & 0.0554 & 0.0555 & 0.0557 & 0.0562 & 0.0556 & 0.0556 & 0.0560 & 0.0565 \\
\midrule
\multirow{4}{*}{$20$} & $\widetilde{\M}_{\mathrm{fac}}$ & 0.0589 & 0.0589 & 0.0589 & 0.0593 & 0.0589 & 0.0590 & 0.0590 & 0.0595 \\
 & $\widehat{\M}_{\mathrm{fac}}$ & 0.0657 & 0.0657 & 0.0658 & 0.0663 & 0.0658 & 0.0659 & 0.0660 & 0.0666 \\
 & $\widetilde{\M}_{\mathrm{nuc}}$ & 0.0589 & 0.0588 & 0.0589 & 0.0593 & 0.0588 & 0.0589 & 0.0590 & 0.0596 \\
 & $\widehat{\M}_{\mathrm{nuc}}$ & 0.0666 & 0.0666 & 0.0668 & 0.0673 & 0.0667 & 0.0668 & 0.0670 & 0.0676 \\
\bottomrule
\end{tabular}
}
\endgroup
\end{table}
\begin{table}[!htbp]
\centering
\caption{Mean relative maximum errors $\mathcal E_\infty$ for the two dependent-noise constructions with $m=500$, $n=250$ and $p=0.2$. Each entry averages 100 repetitions.}
\label{tab:supp:dependent_maxnorm}
\begingroup
\small
\setstretch{1.05}
\setlength{\tabcolsep}{3.2pt}
\resizebox{\linewidth}{!}{%
\begin{tabular}{@{}llrrrrrrrr@{}}
\toprule
$r$ & Estimator & $\gamma=0.2$ & $\gamma=0.4$ & $\gamma=0.6$ & $\gamma=0.8$ & $M=2$ & $M=3$ & $M=5$ & $M=10$ \\
\midrule
\multirow{4}{*}{$5$} & $\widetilde{\M}_{\mathrm{fac}}$ & 0.0645 & 0.0648 & 0.0647 & 0.0652 & 0.0646 & 0.0643 & 0.0643 & 0.0651 \\
 & $\widehat{\M}_{\mathrm{fac}}$ & 0.0686 & 0.0687 & 0.0694 & 0.0699 & 0.0690 & 0.0692 & 0.0702 & 0.0707 \\
 & $\widetilde{\M}_{\mathrm{nuc}}$ & 0.0643 & 0.0648 & 0.0647 & 0.0653 & 0.0645 & 0.0643 & 0.0643 & 0.0656 \\
 & $\widehat{\M}_{\mathrm{nuc}}$ & 0.0677 & 0.0678 & 0.0684 & 0.0690 & 0.0684 & 0.0685 & 0.0691 & 0.0699 \\
\midrule
\multirow{4}{*}{$10$} & $\widetilde{\M}_{\mathrm{fac}}$ & 0.0685 & 0.0686 & 0.0690 & 0.0691 & 0.0686 & 0.0688 & 0.0690 & 0.0688 \\
 & $\widehat{\M}_{\mathrm{fac}}$ & 0.0732 & 0.0738 & 0.0742 & 0.0739 & 0.0734 & 0.0739 & 0.0733 & 0.0740 \\
 & $\widetilde{\M}_{\mathrm{nuc}}$ & 0.0685 & 0.0684 & 0.0690 & 0.0691 & 0.0686 & 0.0687 & 0.0690 & 0.0690 \\
 & $\widehat{\M}_{\mathrm{nuc}}$ & 0.0749 & 0.0757 & 0.0761 & 0.0759 & 0.0752 & 0.0756 & 0.0752 & 0.0759 \\
\midrule
\multirow{4}{*}{$20$} & $\widetilde{\M}_{\mathrm{fac}}$ & 0.0808 & 0.0806 & 0.0806 & 0.0805 & 0.0819 & 0.0811 & 0.0815 & 0.0813 \\
 & $\widehat{\M}_{\mathrm{fac}}$ & 0.0904 & 0.0904 & 0.0910 & 0.0916 & 0.0914 & 0.0916 & 0.0913 & 0.0915 \\
 & $\widetilde{\M}_{\mathrm{nuc}}$ & 0.0805 & 0.0806 & 0.0806 & 0.0805 & 0.0818 & 0.0811 & 0.0815 & 0.0813 \\
 & $\widehat{\M}_{\mathrm{nuc}}$ & 0.0919 & 0.0921 & 0.0926 & 0.0934 & 0.0932 & 0.0932 & 0.0935 & 0.0935 \\
\bottomrule
\end{tabular}
}
\endgroup
\end{table}

% BEGIN MAJOR REVISION S15
\begin{revision}
\subsection{Low sampling probability and ill-conditioned signals}
We set $m=500$, $n=250$, $p=0.05$ and $\sigma=1$, and consider $r=5,10,20$. The condition number takes the values $\kappa\in\{2,5,10,100,1000,10000\}$, with linearly spaced signal singular values, as in \cite{kummerle2021scalable}. The common informative initialization is used in this experiment.
Tables~\ref{tab:supp:ill_fnorm}, \ref{tab:supp:ill_specnorm} and \ref{tab:supp:ill_maxnorm} report the ordinary relative errors defined in Section~\ref{main-sec:numerical}, without additional rescaling by the smallest singular value. Since the signal norm increases with $\kappa$ while the noise scale is fixed, decreasing relative errors across columns should not be interpreted as evidence that a larger condition number makes the recovery problem easier. At each fixed configuration, the factor implementation improves on its least-squares baseline, while the nuclear-norm errors are often close to those of its baseline.

\end{revision}
% END MAJOR REVISION S15
\begin{table}[!htbp]
\centering
\caption{Mean ordinary relative Frobenius errors $\mathcal E_F$ for ill-conditioned signals with $m=500$, $n=250$, $p=0.05$ and $\sigma=1$. Each entry averages 100 repetitions.}
\label{tab:supp:ill_fnorm}
\begingroup
\small
\setstretch{1.05}
\setlength{\tabcolsep}{3.2pt}
\begin{tabular}{@{}llrrrrrr@{}}
\toprule
$r$ & Estimator & $\kappa=2$ & $\kappa=5$ & $\kappa=10$ & $\kappa=100$ & $\kappa=1000$ & $\kappa=10000$ \\
\midrule
\multirow{4}{*}{$5$} & $\widetilde{\M}_{\mathrm{fac}}$ & 0.1093 & 0.0506 & 0.0264 & 0.0027 & 0.0003 & 0.0000 \\
 & $\widehat{\M}_{\mathrm{fac}}$ & 0.1530 & 0.0683 & 0.0360 & 0.0038 & 0.0004 & 0.0000 \\
 & $\widetilde{\M}_{\mathrm{nuc}}$ & 0.1090 & 0.0506 & 0.0264 & 0.0027 & 0.0003 & 0.0000 \\
 & $\widehat{\M}_{\mathrm{nuc}}$ & 0.1260 & 0.0528 & 0.0265 & 0.0033 & 0.0004 & 0.0000 \\
\midrule
\multirow{4}{*}{$10$} & $\widetilde{\M}_{\mathrm{fac}}$ & 0.1164 & 0.0546 & 0.0287 & 0.0030 & 0.0003 & 0.0000 \\
 & $\widehat{\M}_{\mathrm{fac}}$ & 0.1568 & 0.0724 & 0.0387 & 0.0042 & 0.0004 & 0.0000 \\
 & $\widetilde{\M}_{\mathrm{nuc}}$ & 0.1160 & 0.0545 & 0.0286 & 0.0030 & 0.0003 & 0.0000 \\
 & $\widehat{\M}_{\mathrm{nuc}}$ & 0.1314 & 0.0570 & 0.0289 & 0.0032 & 0.0004 & 0.0000 \\
\midrule
\multirow{4}{*}{$20$} & $\widetilde{\M}_{\mathrm{fac}}$ & 0.1252 & 0.0590 & 0.0310 & 0.0032 & 0.0003 & 0.0000 \\
 & $\widehat{\M}_{\mathrm{fac}}$ & 0.1329 & 0.0644 & 0.0347 & 0.0038 & 0.0004 & 0.0000 \\
 & $\widetilde{\M}_{\mathrm{nuc}}$ & 0.1245 & 0.0589 & 0.0310 & 0.0032 & 0.0003 & 0.0000 \\
 & $\widehat{\M}_{\mathrm{nuc}}$ & 0.1286 & 0.0597 & 0.0311 & 0.0033 & 0.0004 & 0.0000 \\
\bottomrule
\end{tabular}
\endgroup
\end{table}
\begin{table}[!htbp]
\centering
\caption{Mean ordinary relative spectral errors $\mathcal E_{sp}$ for ill-conditioned signals with $m=500$, $n=250$, $p=0.05$ and $\sigma=1$. Each entry averages 100 repetitions.}
\label{tab:supp:ill_specnorm}
\begingroup
\small
\setstretch{1.05}
\setlength{\tabcolsep}{3.2pt}
\begin{tabular}{@{}llrrrrrr@{}}
\toprule
$r$ & Estimator & $\kappa=2$ & $\kappa=5$ & $\kappa=10$ & $\kappa=100$ & $\kappa=1000$ & $\kappa=10000$ \\
\midrule
\multirow{4}{*}{$5$} & $\widetilde{\M}_{\mathrm{fac}}$ & 0.0750 & 0.0299 & 0.0149 & 0.0015 & 0.0001 & 0.0000 \\
 & $\widehat{\M}_{\mathrm{fac}}$ & 0.1074 & 0.0430 & 0.0221 & 0.0022 & 0.0002 & 0.0000 \\
 & $\widetilde{\M}_{\mathrm{nuc}}$ & 0.0748 & 0.0299 & 0.0149 & 0.0015 & 0.0001 & 0.0000 \\
 & $\widehat{\M}_{\mathrm{nuc}}$ & 0.0863 & 0.0311 & 0.0150 & 0.0018 & 0.0002 & 0.0000 \\
\midrule
\multirow{4}{*}{$10$} & $\widetilde{\M}_{\mathrm{fac}}$ & 0.0835 & 0.0333 & 0.0166 & 0.0017 & 0.0002 & 0.0000 \\
 & $\widehat{\M}_{\mathrm{fac}}$ & 0.1147 & 0.0467 & 0.0241 & 0.0026 & 0.0003 & 0.0000 \\
 & $\widetilde{\M}_{\mathrm{nuc}}$ & 0.0834 & 0.0332 & 0.0166 & 0.0017 & 0.0002 & 0.0000 \\
 & $\widehat{\M}_{\mathrm{nuc}}$ & 0.0938 & 0.0346 & 0.0167 & 0.0018 & 0.0002 & 0.0000 \\
\midrule
\multirow{4}{*}{$20$} & $\widetilde{\M}_{\mathrm{fac}}$ & 0.0954 & 0.0380 & 0.0190 & 0.0019 & 0.0002 & 0.0000 \\
 & $\widehat{\M}_{\mathrm{fac}}$ & 0.1022 & 0.0426 & 0.0220 & 0.0023 & 0.0002 & 0.0000 \\
 & $\widetilde{\M}_{\mathrm{nuc}}$ & 0.0950 & 0.0380 & 0.0190 & 0.0019 & 0.0002 & 0.0000 \\
 & $\widehat{\M}_{\mathrm{nuc}}$ & 0.0975 & 0.0383 & 0.0189 & 0.0019 & 0.0002 & 0.0000 \\
\bottomrule
\end{tabular}
\endgroup
\end{table}
\begin{table}[!htbp]
\centering
\caption{Mean ordinary relative maximum errors $\mathcal E_\infty$ for ill-conditioned signals with $m=500$, $n=250$, $p=0.05$ and $\sigma=1$. Each entry averages 100 repetitions.}
\label{tab:supp:ill_maxnorm}
\begingroup
\small
\setstretch{1.05}
\setlength{\tabcolsep}{3.2pt}
\begin{tabular}{@{}llrrrrrr@{}}
\toprule
$r$ & Estimator & $\kappa=2$ & $\kappa=5$ & $\kappa=10$ & $\kappa=100$ & $\kappa=1000$ & $\kappa=10000$ \\
\midrule
\multirow{4}{*}{$5$} & $\widetilde{\M}_{\mathrm{fac}}$ & 0.1011 & 0.0435 & 0.0221 & 0.0023 & 0.0002 & 0.0000 \\
 & $\widehat{\M}_{\mathrm{fac}}$ & 0.1421 & 0.0597 & 0.0314 & 0.0033 & 0.0003 & 0.0000 \\
 & $\widetilde{\M}_{\mathrm{nuc}}$ & 0.1009 & 0.0433 & 0.0222 & 0.0022 & 0.0002 & 0.0000 \\
 & $\widehat{\M}_{\mathrm{nuc}}$ & 0.1111 & 0.0440 & 0.0220 & 0.0026 & 0.0003 & 0.0000 \\
\midrule
\multirow{4}{*}{$10$} & $\widetilde{\M}_{\mathrm{fac}}$ & 0.1108 & 0.0494 & 0.0254 & 0.0026 & 0.0003 & 0.0000 \\
 & $\widehat{\M}_{\mathrm{fac}}$ & 0.1472 & 0.0650 & 0.0345 & 0.0038 & 0.0004 & 0.0000 \\
 & $\widetilde{\M}_{\mathrm{nuc}}$ & 0.1102 & 0.0493 & 0.0254 & 0.0026 & 0.0003 & 0.0000 \\
 & $\widehat{\M}_{\mathrm{nuc}}$ & 0.1207 & 0.0504 & 0.0254 & 0.0027 & 0.0003 & 0.0000 \\
\midrule
\multirow{4}{*}{$20$} & $\widetilde{\M}_{\mathrm{fac}}$ & 0.1233 & 0.0556 & 0.0288 & 0.0030 & 0.0003 & 0.0000 \\
 & $\widehat{\M}_{\mathrm{fac}}$ & 0.1287 & 0.0599 & 0.0319 & 0.0034 & 0.0003 & 0.0000 \\
 & $\widetilde{\M}_{\mathrm{nuc}}$ & 0.1225 & 0.0553 & 0.0287 & 0.0030 & 0.0003 & 0.0000 \\
 & $\widehat{\M}_{\mathrm{nuc}}$ & 0.1251 & 0.0559 & 0.0287 & 0.0029 & 0.0003 & 0.0000 \\
\bottomrule
\end{tabular}
\endgroup
\end{table}

\subsection{Gaussian factor signals}
We generate the signal through the factor representation
\begin{align*}
 \M_0=\bm\Lambda'\bm F,\qquad
 \bm\Lambda\in\mathbb R^{r\times m},\quad
 \bm F\in\mathbb R^{r\times n},
\end{align*}
with independent standard Gaussian entries in both factor matrices. We set $r=5$, $m=500$, $n=250$ and $p=0.2$ and vary the homogeneous noise scale.
% BEGIN MAJOR REVISION S16
\begin{revision}
Tables~\ref{tab:supp:factor_fnorm}, \ref{tab:supp:factor_specnorm} and \ref{tab:supp:factor_maxnorm} report the errors. The proposed factor and nuclear-norm implementations have similar accuracy. Their errors decrease with the noise scale, subject to the numerical stopping budget. The nuclear-norm baseline retains appreciable shrinkage error at small noise under its reported regularization settings; consequently, the nuclear comparison reflects regularization as well as the fitting criterion. At $\sigma=2$, the proposed factor error is $0.2701$, compared with $0.3375$ for the factor baseline.

\end{revision}
% END MAJOR REVISION S16
\begin{table}[!htbp]
\centering
\caption{Mean relative Frobenius errors $\mathcal E_F$ for the Gaussian factor signal with $m=500$, $n=250$, $p=0.2$ and $r=5$. Each entry averages 100 repetitions.}
\label{tab:supp:factor_fnorm}
\begingroup
\small
\setstretch{1.05}
\setlength{\tabcolsep}{3.2pt}
\resizebox{\linewidth}{!}{%
\begin{tabular}{@{}llrrrrrrrr@{}}
\toprule
$r$ & Estimator & $\sigma=0$ & $\sigma=0.001$ & $\sigma=0.01$ & $\sigma=0.05$ & $\sigma=0.1$ & $\sigma=0.5$ & $\sigma=1$ & $\sigma=2$ \\
\midrule
\multirow{4}{*}{$5$} & $\widetilde{\M}_{\mathrm{fac}}$ & 0.0000 & 0.0002 & 0.0019 & 0.0094 & 0.0187 & 0.0921 & 0.1704 & 0.2701 \\
 & $\widehat{\M}_{\mathrm{fac}}$ & 0.0000 & 0.0002 & 0.0022 & 0.0103 & 0.0200 & 0.0924 & 0.1753 & 0.3375 \\
 & $\widetilde{\M}_{\mathrm{nuc}}$ & 0.0000 & 0.0002 & 0.0019 & 0.0095 & 0.0189 & 0.0923 & 0.1692 & 0.2677 \\
 & $\widehat{\M}_{\mathrm{nuc}}$ & 0.1328 & 0.1329 & 0.1341 & 0.1391 & 0.1493 & 0.1837 & 0.2278 & 0.3794 \\
\bottomrule
\end{tabular}
}
\endgroup
\end{table}
\begin{table}[!htbp]
\centering
\caption{Mean relative spectral errors $\mathcal E_{sp}$ for the Gaussian factor signal with $m=500$, $n=250$, $p=0.2$ and $r=5$. Each entry averages 100 repetitions.}
\label{tab:supp:factor_specnorm}
\begingroup
\small
\setstretch{1.05}
\setlength{\tabcolsep}{3.2pt}
\resizebox{\linewidth}{!}{%
\begin{tabular}{@{}llrrrrrrrr@{}}
\toprule
$r$ & Estimator & $\sigma=0$ & $\sigma=0.001$ & $\sigma=0.01$ & $\sigma=0.05$ & $\sigma=0.1$ & $\sigma=0.5$ & $\sigma=1$ & $\sigma=2$ \\
\midrule
\multirow{4}{*}{$5$} & $\widetilde{\M}_{\mathrm{fac}}$ & 0.0000 & 0.0002 & 0.0015 & 0.0074 & 0.0148 & 0.0725 & 0.1340 & 0.2118 \\
 & $\widehat{\M}_{\mathrm{fac}}$ & 0.0000 & 0.0002 & 0.0019 & 0.0084 & 0.0161 & 0.0728 & 0.1383 & 0.2708 \\
 & $\widetilde{\M}_{\mathrm{nuc}}$ & 0.0000 & 0.0002 & 0.0015 & 0.0075 & 0.0149 & 0.0725 & 0.1332 & 0.2095 \\
 & $\widehat{\M}_{\mathrm{nuc}}$ & 0.2459 & 0.2460 & 0.2469 & 0.2508 & 0.2633 & 0.2668 & 0.2552 & 0.3246 \\
\bottomrule
\end{tabular}
}
\endgroup
\end{table}
\begin{table}[!htbp]
\centering
\caption{Mean relative maximum errors $\mathcal E_\infty$ for the Gaussian factor signal with $m=500$, $n=250$, $p=0.2$ and $r=5$. Each entry averages 100 repetitions.}
\label{tab:supp:factor_maxnorm}
\begingroup
\small
\setstretch{1.05}
\setlength{\tabcolsep}{3.2pt}
\resizebox{\linewidth}{!}{%
\begin{tabular}{@{}llrrrrrrrr@{}}
\toprule
$r$ & Estimator & $\sigma=0$ & $\sigma=0.001$ & $\sigma=0.01$ & $\sigma=0.05$ & $\sigma=0.1$ & $\sigma=0.5$ & $\sigma=1$ & $\sigma=2$ \\
\midrule
\multirow{4}{*}{$5$} & $\widetilde{\M}_{\mathrm{fac}}$ & 0.0001 & 0.0002 & 0.0018 & 0.0091 & 0.0182 & 0.0880 & 0.1598 & 0.2458 \\
 & $\widehat{\M}_{\mathrm{fac}}$ & 0.0000 & 0.0009 & 0.0051 & 0.0156 & 0.0257 & 0.0912 & 0.1621 & 0.3048 \\
 & $\widetilde{\M}_{\mathrm{nuc}}$ & 0.0001 & 0.0002 & 0.0019 & 0.0095 & 0.0186 & 0.0881 & 0.1591 & 0.2436 \\
 & $\widehat{\M}_{\mathrm{nuc}}$ & 0.1957 & 0.1960 & 0.1977 & 0.2027 & 0.2150 & 0.2355 & 0.2511 & 0.3625 \\
\bottomrule
\end{tabular}
}
\endgroup
\end{table}

% BEGIN MAJOR REVISION S17
\begin{revision}
\subsection{Misspecified working rank}
We fix $m=500$, $n=250$, $p=0.2$ and $\sigma=1$, and use $r\in\{5,10,20\}$. For each true rank, the working rank is fixed in turn at $s=r,r+1,\ldots,r+6$. Only the number of recovery updates is selected by the stopping rule. Thus this experiment assesses sensitivity to overestimating rank rather than automatic rank selection.
Tables~\ref{tab:supp:misrank_fnorm}, \ref{tab:supp:misrank_specnorm} and \ref{tab:supp:misrank_maxnorm} give the mean errors over 100 repetitions for each $(r,s)$ pair. All three errors increase with the working rank, but the proposed factor estimator increases more slowly than least squares. The ratios of the Frobenius error at $s=r+6$ to that at $s=r$ are $1.510$, $1.318$ and $1.197$ for $r=5,10,20$, compared with $2.060$, $1.740$ and $1.680$ for the baseline.

\end{revision}
% END MAJOR REVISION S17
\begin{table}[!htbp]
\centering
\caption{Mean relative Frobenius errors $\mathcal E_F$ of the proposed factor estimator and its least-squares baseline under working-rank misspecification. The settings are $m=500$, $n=250$, $p=0.2$ and $\sigma=1$, with 100 repetitions per configuration.}
\label{tab:supp:misrank_fnorm}
\begingroup
\small
\setstretch{1.05}
\setlength{\tabcolsep}{3.2pt}
\begin{tabular}{@{}llrrrrrrr@{}}
\toprule
$r$ & Estimator & $s=r$ & $s=r+1$ & $s=r+2$ & $s=r+3$ & $s=r+4$ & $s=r+5$ & $s=r+6$ \\
\midrule
\multirow{2}{*}{$5$} & $\widetilde{\M}_{\mathrm{fac}}$ & 0.0692 & 0.0762 & 0.0824 & 0.0883 & 0.0939 & 0.0993 & 0.1045 \\
 & $\widehat{\M}_{\mathrm{fac}}$ & 0.0766 & 0.0909 & 0.1051 & 0.1186 & 0.1318 & 0.1451 & 0.1578 \\
\midrule
\multirow{2}{*}{$10$} & $\widetilde{\M}_{\mathrm{fac}}$ & 0.0703 & 0.0743 & 0.0782 & 0.0819 & 0.0855 & 0.0891 & 0.0926 \\
 & $\widehat{\M}_{\mathrm{fac}}$ & 0.0801 & 0.0896 & 0.0995 & 0.1095 & 0.1194 & 0.1293 & 0.1394 \\
\midrule
\multirow{2}{*}{$20$} & $\widetilde{\M}_{\mathrm{fac}}$ & 0.0766 & 0.0792 & 0.0817 & 0.0842 & 0.0867 & 0.0892 & 0.0917 \\
 & $\widehat{\M}_{\mathrm{fac}}$ & 0.0993 & 0.1109 & 0.1217 & 0.1329 & 0.1444 & 0.1559 & 0.1668 \\
\bottomrule
\end{tabular}
\endgroup
\end{table}
\begin{table}[!htbp]
\centering
\caption{Mean relative spectral errors $\mathcal E_{sp}$ of the proposed factor estimator and its least-squares baseline under working-rank misspecification. The settings are $m=500$, $n=250$, $p=0.2$ and $\sigma=1$, with 100 repetitions per configuration.}
\label{tab:supp:misrank_specnorm}
\begingroup
\small
\setstretch{1.05}
\setlength{\tabcolsep}{3.2pt}
\begin{tabular}{@{}llrrrrrrr@{}}
\toprule
$r$ & Estimator & $s=r$ & $s=r+1$ & $s=r+2$ & $s=r+3$ & $s=r+4$ & $s=r+5$ & $s=r+6$ \\
\midrule
\multirow{2}{*}{$5$} & $\widetilde{\M}_{\mathrm{fac}}$ & 0.0487 & 0.0556 & 0.0560 & 0.0567 & 0.0574 & 0.0581 & 0.0590 \\
 & $\widehat{\M}_{\mathrm{fac}}$ & 0.0539 & 0.0799 & 0.0852 & 0.0889 & 0.0926 & 0.0961 & 0.0989 \\
\midrule
\multirow{2}{*}{$10$} & $\widetilde{\M}_{\mathrm{fac}}$ & 0.0508 & 0.0572 & 0.0581 & 0.0589 & 0.0597 & 0.0606 & 0.0614 \\
 & $\widehat{\M}_{\mathrm{fac}}$ & 0.0580 & 0.0819 & 0.0887 & 0.0932 & 0.0971 & 0.1010 & 0.1047 \\
\midrule
\multirow{2}{*}{$20$} & $\widetilde{\M}_{\mathrm{fac}}$ & 0.0590 & 0.0630 & 0.0647 & 0.0659 & 0.0669 & 0.0680 & 0.0689 \\
 & $\widehat{\M}_{\mathrm{fac}}$ & 0.0773 & 0.1011 & 0.1116 & 0.1185 & 0.1254 & 0.1320 & 0.1369 \\
\bottomrule
\end{tabular}
\endgroup
\end{table}
\begin{table}[!htbp]
\centering
\caption{Mean relative maximum errors $\mathcal E_\infty$ of the proposed factor estimator and its least-squares baseline under working-rank misspecification. The settings are $m=500$, $n=250$, $p=0.2$ and $\sigma=1$, with 100 repetitions per configuration.}
\label{tab:supp:misrank_maxnorm}
\begingroup
\small
\setstretch{1.05}
\setlength{\tabcolsep}{3.2pt}
\begin{tabular}{@{}llrrrrrrr@{}}
\toprule
$r$ & Estimator & $s=r$ & $s=r+1$ & $s=r+2$ & $s=r+3$ & $s=r+4$ & $s=r+5$ & $s=r+6$ \\
\midrule
\multirow{2}{*}{$5$} & $\widetilde{\M}_{\mathrm{fac}}$ & 0.0640 & 0.0733 & 0.0804 & 0.0866 & 0.0923 & 0.0964 & 0.1010 \\
 & $\widehat{\M}_{\mathrm{fac}}$ & 0.0721 & 0.0952 & 0.1173 & 0.1363 & 0.1549 & 0.1715 & 0.1911 \\
\midrule
\multirow{2}{*}{$10$} & $\widetilde{\M}_{\mathrm{fac}}$ & 0.0685 & 0.0741 & 0.0786 & 0.0827 & 0.0876 & 0.0905 & 0.0945 \\
 & $\widehat{\M}_{\mathrm{fac}}$ & 0.0809 & 0.1003 & 0.1171 & 0.1351 & 0.1496 & 0.1654 & 0.1823 \\
\midrule
\multirow{2}{*}{$20$} & $\widetilde{\M}_{\mathrm{fac}}$ & 0.0807 & 0.0846 & 0.0881 & 0.0912 & 0.0939 & 0.0965 & 0.0998 \\
 & $\widehat{\M}_{\mathrm{fac}}$ & 0.1166 & 0.1426 & 0.1630 & 0.1783 & 0.1985 & 0.2158 & 0.2316 \\
\bottomrule
\end{tabular}
\endgroup
\end{table}

\bibliographystyle{abbrvnat}
\bibliography{reference}
\end{document}
